\documentclass[11pt,letterpaper]{article}
\usepackage{xr-hyper}


\usepackage[margin=1in]{geometry}
\usepackage[T1]{fontenc}
\usepackage{placeins}
\usepackage{amsmath,amssymb,amsthm,mathtools}
\usepackage{aliascnt}
\usepackage{newtxtext,newtxmath}
\usepackage{microtype}
\usepackage{booktabs}
\usepackage{array}
\usepackage{enumitem}
\usepackage[table]{xcolor}
\usepackage{url}
\usepackage{algorithm}
\usepackage{algpseudocode}
\usepackage{hyperref}
\makeatletter
\def\theHALG@line{\thealgorithm.\arabic{ALG@line}}
\makeatother
\usepackage[nameinlink,capitalize,noabbrev]{cleveref}

\hypersetup{
  pdftitle={Reed--Solomon Codes Beyond Johnson: Efficient Decoding and Smaller Cryptographic Proofs},
  pdfauthor={Quang Dao, Scott Duke Kominers, and Justin Thaler},
  pdfsubject={Reed--Solomon list decoding and mutual correlated agreement via hidden Hasse derivatives},
  pdfkeywords={Reed--Solomon codes, list decoding, correlated agreement, Hasse derivatives, interpolation},
  colorlinks=true,
  linkcolor=blue!55!black,
  citecolor=blue!55!black,
  urlcolor=blue!55!black
}

\newtheorem{theorem}{Theorem}[section]
\newaliascnt{lemma}{theorem}
\newtheorem{lemma}[lemma]{Lemma}
\aliascntresetthe{lemma}
\newaliascnt{proposition}{theorem}
\newtheorem{proposition}[proposition]{Proposition}
\aliascntresetthe{proposition}
\newaliascnt{corollary}{theorem}
\newtheorem{corollary}[corollary]{Corollary}
\aliascntresetthe{corollary}
\newaliascnt{claim}{theorem}
\newtheorem{claim}[claim]{Claim}
\aliascntresetthe{claim}
\theoremstyle{definition}
\newaliascnt{definition}{theorem}
\newtheorem{definition}[definition]{Definition}
\aliascntresetthe{definition}
\newaliascnt{remark}{theorem}
\newtheorem{remark}[remark]{Remark}
\aliascntresetthe{remark}

\newcommand{\F}{\mathbb{F}}
\newcommand{\R}{\mathbb{R}}
\newcommand{\N}{\mathbb{N}}
\newcommand{\Z}{\mathbb{Z}}
\newcommand{\E}[1]{\mathbb{E}\!\left[#1\right]}
\newcommand{\Prob}{\mathbb{P}}
\newcommand{\RS}{\operatorname{RS}}
\newcommand{\wtdeg}{\operatorname{wdeg}}
\newcommand{\agree}{\operatorname{agr}}
\newcommand{\rank}{\operatorname{rank}}
\newcommand{\poly}{\operatorname{poly}}
\newcommand{\List}{\mathcal{L}}
\newcommand{\Qspace}{\mathcal{Q}}
\newcommand{\positive}[1]{\left(#1\right)_{+}}
\newcommand{\Bpartial}{B_{\partial}}
\newcommand{\Bjet}{B_{\mathrm{jet}}}
\newcommand{\BZ}{B_Z}
\newcommand{\lambdapartial}{\lambda_{\partial}}
\newcommand{\lambdaagr}{\lambda_{\mathrm{agr}}}
\newcommand{\afirst}{a_1}
\newcommand{\rhoswitch}{\rho_c}
\newcommand{\gcurve}{\eta_1}
\newcommand{\gcurvefinite}{\bar{\eta}_1}
\newcommand{\gJohnson}{\eta_0}
\newcommand{\Brecon}{B_{\mathrm{rec}}}
\newcommand{\IntMargin}{\mathsf{Margin}_{\mathrm{int}}}

\title{\bfseries Reed--Solomon Codes Beyond Johnson:\\
Efficient Decoding and Smaller Cryptographic Proofs}
\author{
  Quang Dao\thanks{Carnegie Mellon University. \texttt{qvd@andrew.cmu.edu}.}
  \and
  Scott Duke Kominers\thanks{Harvard University and a16z crypto research.
    \texttt{kominers@fas.harvard.edu}.}
  \and
  Justin Thaler\thanks{a16z crypto research and Georgetown University.
    \texttt{jthaler@a16z.com}.}
}
\date{September 2026}
 \externaldocument[M-][nocite]{rs-beyond-johnson}[rs-beyond-johnson.pdf]
\title{\bfseries Reed--Solomon List Decoding:\\
Certificates, Refinements, and Obstructions}
\hypersetup{
 pdftitle={Reed--Solomon List Decoding: Certificates, Refinements, and Obstructions},
 pdfsubject={Companion results on finite interpolation, differential roots, and list-size obstructions},
 pdfkeywords={Reed--Solomon codes, finite certificates, differential roots, resonance, Frobenius, list-size lower bounds}
}
\begin{document}
\maketitle
\begin{abstract}
Hidden-derivative interpolation for Reed--Solomon codes raises quantitative
questions beyond the existence of an all-rate decoder: how to compute
finite certificates, how characteristic affects coefficient lifting, and
how large agreement lists can be near capacity.
This companion collects exact rank formulas, differential-root refinements,
and combinatorial constructions addressing these questions.
We give finite block decompositions of the local interpolation map,
fixed-parameter certificates of orders zero and one, and a fixed-gap
boundary asymptotic for a one-sided support.
We also analyze descent to the data field, resonant Taylor coefficients,
and randomized regular witnesses.
Coefficient fibers and multiplicative cosets give large agreement lists,
while a Frobenius-plane construction obstructs enumeration from higher-derivative-supported equations
in fixed characteristic.
The all-rate decoding theorem and the geometric list-size and correlated
agreement results are developed in the separate main paper.
 \end{abstract}
\section{Introduction}
\label{sec:companion-introduction}

Our main paper, \emph{Reed--Solomon Codes Beyond Johnson: Efficient Decoding
and Smaller Cryptographic Proofs}, develops explicit
list-decoding and agreement bounds from ordinary and hidden-derivative
interpolation, together with a deterministic decoder.
This companion collects the auxiliary results that support further study
of the method. Its purpose is to retain the finite formulas, alternative
root arguments, and obstructions without reproducing the main decoding theorem or
assembling its proof. The two manuscripts have separate introductions and result
statements, and share basic definitions and the differential-root core.

\paragraph{Finite interpolation.}
\Cref{sec:exact-certificate} gives the exact dimension of a derivative-
weighted monomial space, the kernel of the full local map, an exact rank
for an enlarged map, and a block decomposition of the actual map.
\Cref{sec:one-sided-estimates} preserves the one-sided dimension and
coordinate estimates behind the earlier support analysis.
These formulas distinguish an exact rank from a coordinate-count upper
bound and yield a finite interpolation certificate.
\Cref{sec:optimized-parameters} preserves fixed-parameter low-order
examples and the logarithmically shrinking-gap consequence of the main
paper's certificate. \Cref{sec:quantitative-refinements} preserves finite
first-order and Johnson constants, exact Johnson tuning, and the proved
dominance over an earlier full-differentiation envelope.
\Cref{app:boundary-asymptotics} analyzes a different, one-sided support
with the rate and gap held fixed as the derivative order grows.

\paragraph{Root finding and characteristic.}
\Cref{sec:root-finding} records the regular-lifting and separant-chain lemmas
used by the main paper, then develops alternative all-root and stratified bounds.
\Cref{sec:primitive-normalization,sec:refinements} treat primitive
normalization, descent to the field generated by the data, resonant
coefficient lifting, and randomized regular witnesses.
The extension to characteristic growing linearly with the block length
is a corollary of the main paper's interpolation certificate and the
resonant root bound. The Frobenius obstruction in
\Cref{sec:characteristic} concerns the present support and its
unfiltered differential roots, not every possible list decoder.

\paragraph{Agreement lists near capacity.}
\Cref{sec:sharpness} gives elementary upper bounds, coefficient-fiber
lower bounds for arbitrary evaluation sets, and multiplicative-coset
constructions. The distinctions between fixed gap, varying gap, and
sufficiently large block length are part of the statements.
These examples constrain list-size guarantees without giving matching
lower bounds on the derivative order of hidden-derivative interpolation.

\paragraph{Reading the companion.}
We recall the polynomial and local-interpolation conventions, so most
auxiliary arguments can be read here directly.
References explicitly marked ``of the main paper'' link to
\href{rs-beyond-johnson.pdf}{the accompanying main PDF}; their numbering belongs to
that manuscript. In particular, the support certificate is an external
input where cited. The exact counting and source-audit appendices are
retained here. The main paper's verification paragraph describes the
checked Lean scope; it does not assert formal verification of all the
auxiliary results collected in this companion.

\paragraph{Authorship and assistance.}
This companion was separated from the same AI-assisted research draft
as the main paper. Its account of assistance and author responsibility
applies to both manuscripts. The authors are responsible for all
statements and proofs retained here.
 \section{Preliminaries}
\label{sec:preliminaries}

We collect the notation and background used by the interpolation,
root-counting, and correlated-agreement arguments. The discussion fixes
our agreement and Hasse-derivative conventions, records the backward
Taylor identity, and introduces the Taylor-jet and algebraic-geometric
language needed later.

Throughout the paper, \(\F_q\) denotes a finite field,
\(\N=\{0,1,2,\ldots\}\), and \([n]=\{1,\ldots,n\}\).
The local interpolation argument in
\Cref{sec:hidden-derivative-interpolation} works in arbitrary
characteristic; root-finding statements specify their characteristic
hypotheses.  We write \(\F_q[X]_{\leq D}\) for the vector space of univariate
polynomials of degree at most \(D\), including the zero polynomial.
We take $\deg 0=-\infty$, so the degree inequalities below also include it.

\paragraph{Counting and probability notation.}
For a real number $x$, its positive part is $x_+:=\max\{x,0\}$.
Thus $(u-z)_+^2$ means $\bigl(\max\{u-z,0\}\bigr)^2$; in the
integrals below the integrand is zero whenever $z\ge u$.
The indicator $\mathbf1_{\{\mathcal P\}}$ equals one when the condition
$\mathcal P$ holds and zero otherwise. We write $\#S$ or $|S|$ for
the cardinality of a finite set $S$. For a nonnegative vector $x$,
$|x|$ denotes the sum of its coordinates; for a scalar it denotes
absolute value.

We write $\Pr[\mathcal A]$ for the probability of an event and
$\E{U}$ for the expectation of a random variable $U$, under the
specified distribution. In particular, $\E{U^2}$ is the second moment,
whereas $\bigl(\E{U}\bigr)^2$ is the square of the mean.
The variance is $\operatorname{Var}(U)=\E{(U-\E{U})^2}$.
For a Euclidean region $S$, $\operatorname{vol}(S)$ denotes its Lebesgue
volume; a uniform point of $S$ is sampled from normalized volume.

Unless a base is displayed, $\log$ denotes the natural logarithm.
We use the harmonic sums
$H_t=\sum_{j=1}^t j^{-1}$ and $H_t^{(2)}=\sum_{j=1}^t j^{-2}$,
with both sums zero at $t=0$. A subscript on asymptotic notation,
such as $O_\delta(\cdot)$, indicates that its implicit constant may
depend on that parameter. Section-specific notation is introduced
where its role in the argument becomes relevant.

\subsection{Reed--Solomon codes and agreement}

Fix pairwise distinct evaluation points
\(\boldsymbol{\alpha}=(\alpha_1,\ldots,\alpha_n)\in\F_q^n\).  The Reed--Solomon
code of message dimension \(k\) on this evaluation set is
\begin{equation}
  \RS_{\boldsymbol{\alpha}}[n,k]
  :=
  \bigl\{
    (P(\alpha_1),\ldots,P(\alpha_n)):
    P\in\F_q[X],\ \deg P<k
  \bigr\}.
  \label{eq:rs-code}
\end{equation}
Thus the code has rate \(\rho=k/n\); no multiplicative or additive structure is
assumed of the evaluation set.

For \(u,v\in\F_q^n\), their agreement is the number of coordinates on which
they coincide:
\begin{equation}
  \agree(u,v):=\bigl|\{i\in[n]:u_i=v_i\}\bigr|.
  \label{eq:agreement}
\end{equation}
Given a received word \(y\in\F_q^n\) and an integer \(A\), the agreement list
at ambient degree \(D\) is
\begin{equation}
  \List_D(y,A)
  :=
  \bigl\{
    P\in\F_q[X]_{\leq D}:
    \agree((P(\alpha_i))_{i=1}^n,y)\geq A
  \bigr\}.
  \label{eq:agreement-list}
\end{equation}
The decoder will first work in such an ambient space and then retain only the
polynomials of degree less than the actual message dimension.

\subsection{Hasse derivatives and multiplicity}

We use the standard Hasse-derivative convention~\cite{DKSS13}.

\begin{definition}[Hasse derivative]
\label{def:hasse-derivative}
For \(P\in\F_q[X]\), its \(j\)-th Hasse derivative \(P^{[j]}\) is defined by
the polynomial identity
\begin{equation}
  P(X+T)=\sum_{j\geq 0}P^{[j]}(X)T^j.
  \label{eq:hasse-taylor}
\end{equation}
Equivalently, if \(P(X)=\sum_i p_iX^i\), then
\begin{equation}
  P^{[j]}(X)=\sum_{i\geq j}\binom{i}{j}p_iX^{i-j}.
  \label{eq:hasse-coefficients}
\end{equation}
\end{definition}
The definition is valid in every characteristic and gives
\(\deg P^{[j]}\leq D-j\) whenever \(\deg P\leq D\), with a zero derivative
when \(j>D\).

Ordinary Taylor expansion uses derivatives at $X$ to obtain the value
at $X+T$. The backward identity below uses derivatives at $X+T$ to
recover the value at $X$. The hidden-derivative constraints use this
inversion of
\Cref{eq:hasse-taylor}.  We include the short proof because every sign and
power of \(T\) enters the local multiplicity calculation.

\begin{lemma}[Backward Taylor identity]
\label{lem:backward-taylor}
Every \(P\in\F_q[X]\) satisfies
\begin{equation}
  P(X)=\sum_{j\geq 0}(-T)^jP^{[j]}(X+T)
  \qquad\text{in }\F_q[X,T].
  \label{eq:backward-taylor}
\end{equation}
Consequently, for every \(d\geq 0\),
\begin{equation}
  P(X+T)
  \equiv
  P(X)+\sum_{j=1}^{d}(-1)^{j+1}T^jP^{[j]}(X+T)
  \pmod{T^{d+1}}.
  \label{eq:backward-taylor-truncated}
\end{equation}
\end{lemma}

\begin{proof}
Apply \Cref{eq:hasse-taylor} with \(X\) replaced by \(X+T\) and the increment
replaced by \(-T\).  This gives \Cref{eq:backward-taylor}.  Isolating its
\(j=0\) term and discarding the terms with \(j>d\) gives
\Cref{eq:backward-taylor-truncated}.
\end{proof}

\begin{definition}[Multiplicity]
\label{def:root-multiplicity}
For \(F\in\F_q[X]\) and \(\alpha\in\F_q\), the multiplicity of \(F\) at
\(\alpha\) is
\begin{equation}
  \operatorname{mult}_{\alpha}(F)
  :=\max\{s\in\N:T^s\text{ divides }F(\alpha+T)\}.
  \label{eq:root-multiplicity}
\end{equation}
We set \(\operatorname{mult}_{\alpha}(0)=\infty\).
\end{definition}

We use the standard fact that a nonzero univariate polynomial cannot have more
roots, counted with multiplicity, than its degree.  In the form needed below,
if \(\alpha_1,\ldots,\alpha_t\) are distinct and
\(\operatorname{mult}_{\alpha_i}(F)\geq m\) for every \(i\), then either
\(F=0\) or \(\deg F\geq mt\).

For a bivariate polynomial $Q(X,Y)$, multiplicity at least $s$ at
$(\alpha,y)$ means that every coefficient of $U^aV^b$ with $a+b<s$
vanishes in $Q(\alpha+U,y+V)$. This is the notion used by the
order-zero interpolation argument. If $P(\alpha)=y$, both $T$ and
$P(\alpha+T)-y$ are divisible by $T$. Substituting them for $U,V$
therefore shows that $Q(X,P(X))$ has a root of multiplicity at least
$s$ at $X=\alpha$.

\subsection{Taylor data and polynomial differential equations}
\label{sec:taylor-terminology}

\paragraph{Taylor coefficients and jets.}
At a chosen point $a$, the identity
$P(a+T)=\sum_j P^{[j]}(a)T^j$ writes the polynomial in powers of
$X-a$. We call $a$ the \emph{Taylor center} and the values
$P^{[j]}(a)$ its \emph{Taylor coefficients} there. An order-$r$
\emph{jet} is simply the finite tuple
\[
 (P(a),P^{[1]}(a),\ldots,P^{[r]}(a)).
\]
Equivalently, it records $P(a+T)$ modulo $T^{r+1}$.
Thus ``initial Taylor data'' and ``initial jet'' refer to these first
$r+1$ coefficients; neither term requires analytic power series or
convergence.

\paragraph{Differential equations and jet degree.}
We encode a polynomial differential equation by a nonzero polynomial
$Q(X,Y_0,\ldots,Y_r)$. A polynomial $P$ solves it when substituting
$Y_j=P^{[j]}(X)$ makes $Q$ identically zero. The $Y_j$ are called
\emph{jet variables} because, at a fixed $X=a$, they stand for the
entries of the jet. The \emph{order} is the largest index $r$ on
which $Q$ actually depends. Its \emph{total jet degree} is its total
degree in $Y_0,\ldots,Y_r$, ignoring $X$ and any challenge variable.
These are different quantities: $Q=Y_1-Y_0^2$ has order one and
jet degree two.

\paragraph{The separant and regular Taylor data.}
For an equation of actual order $r$, its \emph{separant} is the
ordinary partial derivative
\[
 S=\frac{\partial Q}{\partial Y_r}.
\]
A solution's jet at $a$ is \emph{regular for this equation} if
$S(a,P(a),\ldots,P^{[r]}(a))\ne0$; otherwise we call it singular.
The separant matters because it multiplies each new unknown coefficient
in the reconstruction equation. Under the characteristic hypotheses
used below, a regular jet therefore determines successive coefficients
uniquely. Singular solutions are handled by passing to the separant
equation and repeating; this is the \emph{separant chain}.

\paragraph{Taylor reconstruction and charts.}
A \emph{Taylor recurrence} is the equation used to compute the next
coefficient from those already known. A \emph{lift} extends an initial
jet to a longer coefficient sequence satisfying the equation to the
required order. The resulting formulas express the later coefficients
as rational functions of the initial ones. We call this
parameterization a \emph{Taylor chart}, restricted to the initial data
where its denominators are nonzero. In particular, a chart is a map
from a short tuple of coefficients to the reconstructed polynomial
coefficients; its degree bounds control the subsequent counting argument.

\subsection{Weighted degree}

\begin{definition}[Weighted degree]
\label{def:weighted-degree}
For nonnegative integer weights \(w_0,\ldots,w_s\), the weighted degree of a
monomial \(Z_0^{e_0}\cdots Z_s^{e_s}\) is
\(\sum_{i=0}^s w_ie_i\).  The weighted degree
\(\wtdeg_{w_0,\ldots,w_s}Q\) of a nonzero polynomial \(Q\) is the maximum
weighted degree of one of its monomials with nonzero coefficient.
We set $\wtdeg_{w_0,\ldots,w_s}(0)=-\infty$.
\end{definition}

The weights relevant to polynomial differential equations record the actual
degree lost under Hasse differentiation.  Namely, if \(d\leq D\),
\(Q\in\F_q[X,Y_0,\ldots,Y_d]\), and \(P\in\F_q[X]_{\leq D}\), then
\begin{equation}
  \deg_X Q\bigl(X,P(X),P^{[1]}(X),\ldots,P^{[d]}(X)\bigr)
  \leq
  \wtdeg_{1,D,D-1,\ldots,D-d}Q.
  \label{eq:weighted-specialization-degree}
\end{equation}
Indeed, the variable \(X\) contributes degree one and the substitution for
\(Y_j\) contributes degree at most \(D-j\).  This elementary accounting is
the reason to distinguish the actual derivative weights from the common
upper bound $D$ used in the weighted-support certificate.


\subsection{Terminology for counting polynomial families}
\label{sec:geometric-terminology}

The counting proofs use sets of coefficient tuples satisfying polynomial
equations. We recall the geometric terminology here; the specific
counting inequalities are stated where they are used.

An \emph{affine algebraic set} is a common zero set of polynomials in
finitely many coordinates. We work over an algebraically closed field,
meaning that every nonconstant univariate polynomial splits into linear
factors. A \emph{locus} simply means the set satisfying the stated
conditions; these may include inequalities such as $S\ne0$.
A \emph{hypersurface} is the zero set of one nonconstant polynomial;
a \emph{hyperplane} is the case of a degree-one equation.
The \emph{Zariski closure} of a set is the smallest algebraic set
containing it. An \emph{irreducible component} is a maximal piece
that cannot be written as a union of two proper algebraic subsets.

Intuitively, dimension counts the independent parameters in such a
family. The \emph{degree} of an irreducible $s$-dimensional algebraic
set is the number of points, counted with multiplicity, obtained by
intersecting its projective closure (which also includes points at
infinity) with a linear subspace of
codimension $s$ in general position. Here general position means that
the choices avoid the algebraic conditions causing a degenerate
intersection. A point and a linear space each have degree one.
For a set with several components, our \emph{total degree} sums their
degrees, including components of different dimensions. The degree
bounds limit how many pieces can remain as we impose agreement equations.

A \emph{rational map} has coordinates that are ratios of polynomials;
it is used only where its denominators are nonzero. Its
\emph{image closure} is the Zariski closure of the values it takes.
An \emph{incidence} is a point together with an equation it satisfies.
Thus counting agreement incidences means counting pairs consisting of
a candidate and a coordinate where it agrees with the received word.
 \paragraph{A staircase count.}
For the finite certificates and boundary asymptotics, we use
\begin{equation}
 \mathsf G(z)=\sum_{j\ge0}(z-j)_+,
 \qquad (t)_+=\max\{t,0\}.
 \label{eq:normalized-staircase-function}
\end{equation}
Only finitely many summands are nonzero.
The count $\mathsf N_D(L)=\#\{(u,a)\in\N^2:u+Da<L\}$ satisfies
$\mathsf N_D(L)=D\mathsf G(L/D)$ for integer $L$ and $D\ge1$.
 \section{Hidden-derivative interpolation}
\label{sec:hidden-derivative-interpolation}
We use hidden-derivative interpolation to construct a nonzero equation
satisfied by every sufficiently close message. We call the polynomial
representing such an equation a hidden-derivative explainer, or simply
an explainer.
\begin{definition}[Explainer]
\label{def:explainer}
Fix a field $F$, distinct evaluation points $\alpha_1,\ldots,\alpha_n\in F$,
a received word $w\in F^n$, integers $K\ge1$ and $1\le A\le n$,
and an order bound $d\ge0$.
A nonzero polynomial $Q\in F[X,Y_0,\ldots,Y_d]$ is an
\emph{explainer} for $w$ at agreement threshold $A$ and degree bound $K$
if every $P\in F[X]$ with $\deg P<K$ satisfies
\[
 \bigl|\{i:P(\alpha_i)=w_i\}\bigr|\ge A
 \quad\Longrightarrow\quad
 Q\bigl(X,P,P^{[1]},\ldots,P^{[d]}\bigr)\equiv0.
\]
\end{definition}

The implication is one-way: an explainer can have polynomial solutions
that are far from the received word. Agreement must still be imposed
when recovering or counting the requested messages.
The ambient degree bound $K$ is part of the definition and may exceed
the actual message dimension $k$; explaining all degree-$<K$ candidates also explains
the degree-$<k$ candidates.
 \subsection{Local interpolation constraints}
\label{sec:local-interpolation}
BCPZZ's backward-Taylor substitution turns agreement into multiplicity
without requiring derivative evaluations~\cite[Section~3.1]{BCPZZ26}.
Fix an order $d\ge0$ and a target multiplicity $m\ge1$. For formal
variables $Y=(Y_1,\ldots,Y_d)$, write
$B_d(T,Y)=\sum_{j=1}^d(-1)^{j+1}T^jY_j$; an empty sum is zero.
The identity in \Cref{eq:backward-taylor-truncated} says that, at an
agreement, the remainder after these $d$ terms is divisible by $T^{d+1}$.
We encode that divisibility directly in the local substitution.

\begin{definition}[Local hidden-derivative constraints]
\label{def:local-constraints}
For $(\alpha,y)\in\F^2$ and $Q\in\F[X,Y_0,\ldots,Y_d]$, put
\begin{equation}
 \Phi_{\alpha,y}^{d,m}(Q):=
 Q(\alpha+T,y+B_d(T,Y)+T^{d+1}V,Y_1,\ldots,Y_d)\pmod{T^m}
 \label{eq:local-substitution}
\end{equation}
in $\F[T,V,Y_1,\ldots,Y_d]/(T^m)$. The multiplicity-$m$, order-$d$
constraints at $(\alpha,y)$ are
\begin{equation}
 \Phi_{\alpha,y}^{d,m}(Q)=0.
 \label{eq:local-constraints}
\end{equation}
\end{definition}
These are homogeneous linear conditions on any fixed interpolation space.
They require every coefficient in $V,Y_1,\ldots,Y_d$ to be divisible by
$T^m$ before reduction.

\begin{lemma}[Local multiplicity {\cite[Lemma~3.1]{BCPZZ26}}]
\label{lem:local-multiplicity}
If $P(\alpha)=y$ and $Q$ satisfies \eqref{eq:local-constraints}, then
\begin{equation}
 Q(\alpha+T,P(\alpha+T),P^{[1]}(\alpha+T),\ldots,P^{[d]}(\alpha+T))
 \equiv0\pmod{T^m}.
 \label{eq:local-multiplicity-conclusion}
\end{equation}
Thus each agreement gives a root of multiplicity at least $m$ after
substituting the candidate and its derivatives into $Q$.
\end{lemma}
\begin{proof}
The backward Taylor identity makes
\begin{equation}
 V_P(T)=\frac{P(\alpha+T)-y-
 \sum_{j=1}^d(-1)^{j+1}T^jP^{[j]}(\alpha+T)}{T^{d+1}}
 \label{eq:local-remainder}
\end{equation}
a polynomial. Substitute $Y_j=P^{[j]}(\alpha+T)$ and $V=V_P(T)$
into \eqref{eq:local-constraints}. The second input to $Q$ becomes
$P(\alpha+T)$, proving the claim in every characteristic.
\end{proof}

For the rank calculations we retain the equivalent remainder variable
$E=T^dV$. Substitute $Y_0=y+B_d(T,Y)+TE$ and retain the rows
\begin{equation}
 T^iE^uY^{\mathbf j}\quad\text{with }i+du<m.
 \label{eq:local-e-row-cutoff}
\end{equation}
Rows with $i<u$ are structurally zero, since each $E$ comes with a factor
of $T$. On the remaining rows, the substitution relabels
\begin{equation}
 (i,u,\mathbf j)\longmapsto(i+du,u,\mathbf j).
 \label{eq:local-row-relabeling}
\end{equation}
Its image is exactly the possible nonzero rows of
\eqref{eq:local-substitution}, whose $T$-degrees are at least $(d+1)u$
and below $m$. Thus it preserves every matrix entry, kernel, rank, and
jet-degree grouping. If $y$ depends on a challenge, its challenge degrees
are preserved as well. We can therefore compute ranks in the $E$
coordinates while using \Cref{lem:local-multiplicity} for every order.
 
\subsection{Obtaining an explainer from the local constraints}
We now combine the conditions from all received positions. Enough
available coefficients will guarantee a nonzero interpolant; enough
agreements will force its specialization to vanish identically. The
result is an explainer for the received word. The following principle
separates these requirements so that we can next choose a support
satisfying them.


Let \(\Qspace\subseteq\F_q[X,Y_0,\ldots,Y_d]\) be a finite-dimensional
monomial space.  The local constraints define a linear map
\(\Phi_{\alpha,y}\) from \(\Qspace\) to the vector of coefficients required to
vanish in \Cref{eq:local-constraints}.  Write
\(\rank_{\mathrm{loc}}\) for a common upper bound on
\(\rank(\Phi_{\alpha,y})\), uniform over \((\alpha,y)\).

The next proposition abstracts the final dimension-versus-rank and multiplicity step of
\cite[Proposition~3.13]{BCPZZ26}. The ambient-degree formulation distinguishes the message dimension from the
larger space used for interpolation; the underlying argument is theirs.

\begin{proposition}[BCPZZ global interpolation principle, in ambient-degree form]
\label{prop:global-interpolation-principle}
Fix distinct \(\alpha_1,\ldots,\alpha_n\in\F_q\), a received word
\(y=(y_1,\ldots,y_n)\), and an agreement threshold \(A\).  Suppose that:
\begin{enumerate}[label=\textup{(\roman*)}]
  \item every \(Q\in\Qspace\) and \(P\in\F_q[X]_{\leq D}\) satisfy
  \begin{equation}
    \deg_X Q(X,P,P^{[1]},\ldots,P^{[d]})<mA;
    \label{eq:global-specialization-degree-hypothesis}
  \end{equation}
  \item the local constraint map at each received point has rank at most
  \(\rank_{\mathrm{loc}}\); and
  \item \(\dim\Qspace>n\rank_{\mathrm{loc}}\).
\end{enumerate}
Then there is a nonzero \(Q\in\Qspace\) satisfying the local constraints at
every \((\alpha_i,y_i)\).  Every polynomial
\(P\in\F_q[X]_{\leq D}\) that agrees with \(y\) in at least \(A\) positions
then satisfies the polynomial identity
\begin{equation}
  Q(X,P(X),P^{[1]}(X),\ldots,P^{[d]}(X))\equiv 0.
  \label{eq:global-differential-identity}
\end{equation}
\end{proposition}

\begin{proof}
Stack the \(n\) local maps into one homogeneous linear map on \(\Qspace\).
Its rank is at most \(n\rank_{\mathrm{loc}}\), so the strict dimension
inequality gives a nonzero element \(Q\) in its kernel.

Fix a polynomial \(P\) with at least \(A\) agreement positions and set
\begin{equation}
  F_P(X):=Q(X,P(X),P^{[1]}(X),\ldots,P^{[d]}(X)).
  \label{eq:specialized-interpolant}
\end{equation}
At each agreement position, \Cref{lem:local-multiplicity} gives a root of
\(F_P\) of multiplicity at least \(m\).  The roots are distinct, so a nonzero
\(F_P\) would have degree at least \(mA\).  This contradicts
\Cref{eq:global-specialization-degree-hypothesis}; hence \(F_P=0\).
\end{proof}
  \section{Exact finite interpolation refinements}
\label{sec:exact-certificate}

We now choose an interpolation space and evaluate both sides of the dimension
versus rank comparison without asymptotic estimates. BCPZZ introduce the
weighted derivative support, factor the local map through an enlarged map
\(\Gamma\), and exhibit independent subspaces of \(\ker\Gamma\). We retain
these objects but strengthen their analysis in four ways. Each derivative
variable is charged its actual specialization degree, no separate upper bound is
placed on the ordinary degree in \(Y_2,\ldots,Y_d\), the exhibited kernel
subspaces are proved to exhaust \(\ker\Gamma\), and the actual local map is
decomposed into exact finite blocks. The resulting test is a finite inequality
in integers.

Fix integers
\begin{equation}
  1\leq A\leq n,
  \qquad
  1\leq d<D,
  \qquad
  m\geq 1,
  \qquad
  M\geq 0,
  \qquad
  W\geq 0.
  \label{eq:certificate-parameters}
\end{equation}
For \(c=(c_2,\ldots,c_d)\in\N^{d-1}\), define
\begin{equation}
  |c|:=\sum_{j=2}^{d}c_j,
  \qquad
  \omega(c):=\sum_{j=2}^{d}(j-1)c_j,
  \qquad
  Y^c:=\prod_{j=2}^{d}Y_j^{c_j}.
  \label{eq:omega-definition}
\end{equation}
We use the weighted-simplex count
\begin{equation}
  \Lambda_d(z)
  :=
  \bigl|
    \{c\in\N^{d-1}:\omega(c)\leq z\}
  \bigr|,
  \label{eq:lambda-definition}
\end{equation}
with the convention \(\Lambda_d(z)=0\) for \(z<0\).

\subsection{The derivative-weighted interpolation space}

\begin{definition}[Derivative-weighted interpolation space]
\label{def:cap-free-interpolation-space}
Let \(\Qspace(D,A,d,m,M,W)\) be the span of all monomials
\(X^aY_0^{b_0}Y_1^{b_1}Y^c\) indexed by nonnegative integers
\(a,b_0,b_1\) and \(c\in\N^{d-1}\) such that
\begin{align}
  b_1&\leq M,
  &\omega(c)&\leq W,
  \label{eq:cap-free-support}\\
  a+Db_0+(D-1)b_1+
  \sum_{j=2}^{d}(D-j)c_j&<mA.
  &&
  \label{eq:cap-free-weighted-degree}
\end{align}
\end{definition}
Because \(D-j>0\) for every \(j\leq d\),
\Cref{eq:cap-free-weighted-degree} makes this a finite-dimensional space even
without a separate bound on \(|c|\).

\begin{lemma}[Degree bounds]
\label{lem:cap-free-degree-bounds}
Every \(Q\in\Qspace(D,A,d,m,M,W)\) has the following properties:
\begin{align}
  \wtdeg_{1,D,D-1,\ldots,D-d}Q&<mA,
  \label{eq:cap-free-weighted-bound}\\
  \deg_X Q(X,P,P^{[1]},\ldots,P^{[d]})&<mA
  \quad\text{for every }P\in\F_q[X]_{\leq D},
  \label{eq:cap-free-specialization-bound}\\
  \deg_{Y_0,\ldots,Y_d}Q
  &\leq
  \left\lfloor\frac{mA-1}{D-d}\right\rfloor.
  \label{eq:cap-free-total-y-degree}
\end{align}
\end{lemma}

\begin{proof}
The first claim is the defining inequality.  The second follows from
\Cref{eq:weighted-specialization-degree}.  For the last claim, every
\(Y_j\) has weight at least \(D-d\).  A monomial of total \(Y\)-degree
\(s\) therefore has weighted degree at least \((D-d)s\); the strict integer
inequality in \Cref{eq:cap-free-weighted-degree} gives
\((D-d)s\leq mA-1\).
\end{proof}

To state the dimension exactly, define the staircase count
\begin{equation}
  \mathsf N_D(t)
  :=
  \sum_{b\geq 0}\positive{t-Db}
  =
  \begin{cases}
    \displaystyle
    \sum_{b=0}^{\lfloor(t-1)/D\rfloor}(t-Db),&t>0,\\[1.2ex]
    0,&t\leq 0.
  \end{cases}
  \label{eq:staircase-count}
\end{equation}
For integer \(t\), this is the number of pairs
\((a,b)\in\N^2\) satisfying \(a+Db<t\).

\begin{proposition}[Exact dimension of the interpolation space]
\label{prop:exact-cap-free-dimension}
The interpolation space in \Cref{def:cap-free-interpolation-space} has
dimension
\begin{equation}
  \dim\Qspace(D,A,d,m,M,W)
  =
  \sum_{\substack{c\in\N^{d-1}\\\omega(c)\leq W}}
  \ \sum_{b_1=0}^{M}
  \mathsf N_D\left(
    mA-(D-1)b_1-
    \sum_{j=2}^{d}(D-j)c_j
  \right).
  \label{eq:exact-cap-free-dimension}
\end{equation}
\end{proposition}

\begin{proof}
Fix \(c\) and \(b_1\) satisfying \Cref{eq:cap-free-support}.  The remaining
indices \((a,b_0)\) satisfy
\begin{equation}
  a+Db_0
  <
  mA-(D-1)b_1-
  \sum_{j=2}^{d}(D-j)c_j.
  \label{eq:dimension-residual-inequality}
\end{equation}
By \Cref{eq:staircase-count}, their number is the corresponding summand in
\Cref{eq:exact-cap-free-dimension}.  Distinct index tuples give distinct
monomials, so summing over \(c\) and \(b_1\) gives equality.
\end{proof}

\subsection{The full local kernel}

Before passing to the enlarged map, it is useful to record the exact kernel
of the normalized local substitution on the full polynomial ring.  This
description is independent of any interpolation support and is new to this
work.  Put
\begin{equation}
  B:=Y_0-\sum_{j=1}^{d}(-1)^{j+1}X^jY_j.
  \label{eq:backward-taylor-coordinate}
\end{equation}
After normalizing the received point to \((0,0)\) and renaming the direct
remainder variable $V$ from \Cref{eq:local-substitution} as $Z$, the local
map is
\begin{equation}
  X\longmapsto T,\qquad
  B\longmapsto T^{d+1}Z,\qquad
  Y_j\longmapsto Y_j
  \quad(1\leq j\leq d),
  \label{eq:normalized-full-local-map}
\end{equation}
with the target reduced modulo \(T^m\).
Denote this full-ring map by \(\Phi^{\rm full}_{0,0}\).

\begin{proposition}[Exact kernel of the full local map]
\label{prop:full-local-kernel}
The kernel of \eqref{eq:normalized-full-local-map} is the monomial ideal in
the coordinates \((X,B,Y_1,\ldots,Y_d)\)
\begin{equation}
  \mathcal I_m=
  \left(
    X^{\max\{0,m-(d+1)b\}}B^b:
    0\leq b\leq\left\lceil\frac{m}{d+1}\right\rceil
  \right).
  \label{eq:full-local-kernel-ideal}
\end{equation}
Consequently, for every finite-dimensional polynomial space
\(\mathcal U\),
\begin{equation}
  \rank(\Phi^{\rm full}_{0,0}|_{\mathcal U})
  =\dim\mathcal U-\dim(\mathcal U\cap\mathcal I_m).
  \label{eq:full-local-rank-intersection}
\end{equation}
At a general received point \((\alpha,y)\), the same statement holds after
replacing \(X\) by \(X-\alpha\) and \(B\) by
\[
  B_{\alpha,y}:=
  Y_0-y-\sum_{j=1}^{d}(-1)^{j+1}(X-\alpha)^jY_j.
\]
\end{proposition}

\begin{proof}
The change of coordinates from \(Y_0\) to \(B\) is triangular.  Hence the
monomials \(X^aB^bY_1^{c_1}\cdots Y_d^{c_d}\) form a basis of the source.
Under \eqref{eq:normalized-full-local-map}, such a monomial becomes
\[
  T^{a+(d+1)b}Z^bY_1^{c_1}\cdots Y_d^{c_d}.
\]
Distinct source monomials that survive modulo \(T^m\) have distinct target
monomials.  The kernel is therefore spanned by those with
\(a+(d+1)b\geq m\), and the displayed generators are exactly the minimal
monomials with that property.  Restricting a linear map to \(\mathcal U\)
gives \eqref{eq:full-local-rank-intersection}.  Translation gives the final
statement.
\end{proof}

\subsection{The exact enlarged-map rank}

Fix one received point \((\alpha,y)\).  Following~\cite{BCPZZ26}, first make
only the substitutions
\begin{equation}
  X=\alpha+T,
  \qquad
  Y_0=y+TU,
  \label{eq:rank-first-substitution}
\end{equation}
and reduce modulo \(T^m\).  The resulting map \(\Psi_{\alpha,y}\) takes
\(\Qspace(D,A,d,m,M,W)\) into the space
\begin{equation}
  \mathcal V
  :=
  \operatorname{span}\left\{
    T^rU^aY_1^bY^c:
    \begin{array}{l}
      0\leq r<m,\ 0\leq a\leq r,\ 0\leq b\leq M,\\
      c\in\N^{d-1},\ \omega(c)\leq W+r
    \end{array}
  \right\}.
  \label{eq:rank-space-v}
\end{equation}
The relaxed bound \(\omega(c)\leq W+r\) contains the image and is chosen so
that the kernel below remains inside \(\mathcal V\).

To compute the rank, we return to the equivalent $E$ coordinates from
\Cref{eq:local-e-row-cutoff}. Solving $Y_0=y+TU$ for $E$ gives
\begin{equation}
  E=U-\sum_{j=1}^{d}(-1)^{j+1}T^{j-1}Y_j.
  \label{eq:e-in-u}
\end{equation}
Let \(\Gamma\) rewrite a polynomial in \(\mathcal V\) in the variables
\(T,E,Y_1,\ldots,Y_d\), equivalently by substituting
\(U=E+\sum_{j=1}^{d}(-1)^{j+1}T^{j-1}Y_j\), and then retain exactly the
coefficients of monomials
\(T^iE^bY_1^{e_1}\cdots Y_d^{e_d}\) for which \(i+db<m\).
We identify its nonzero target rows with those of the direct remainder map
through \Cref{eq:local-row-relabeling}. With this identification, the local
constraint map factors as
\begin{equation}
  \Phi_{\alpha,y}=\Gamma\circ\Psi_{\alpha,y},
  \label{eq:local-map-factorization}
\end{equation}
and hence \(\rank(\Phi_{\alpha,y})\leq\rank(\Gamma)\).

For \(0\leq r<m\), define the least remainder exponent that crosses the local
multiplicity threshold by
\begin{equation}
  h_r:=\left\lceil\frac{m-r}{d}\right\rceil.
  \label{eq:hr-definition}
\end{equation}
For each such \(r\), expand \(E\) according to \Cref{eq:e-in-u} and define
the subspace
\begin{equation}
  \mathcal K_r
  :=
  T^rE^{h_r}\cdot
  \operatorname{span}\left\{
    U^aY_1^bY^c:
    \begin{array}{l}
      0\leq a\leq r-h_r,\ 0\leq b\leq M-h_r,\\
      c\in\N^{d-1},\ \omega(c)\leq W+r
    \end{array}
  \right\},
  \label{eq:kernel-space-kr}
\end{equation}
where a negative upper bound makes the corresponding span zero.
The subspaces \(\mathcal K_r\) are the kernel subspaces constructed in
\cite[Lemma~3.11]{BCPZZ26}. BCPZZ prove that their direct sum lies in
\(\ker\Gamma\) and use its dimension as a lower bound. The following result
retains the positive parts that disappear in their coarser asymptotic analysis
and strengthens the inclusion to an equality.

\begin{lemma}[Exact rank of the enlarged map]
\label{lem:exact-local-rank}
The kernel of \(\Gamma\) is
\begin{equation}
  \ker\Gamma
  =
  \bigoplus_{r=0}^{m-1}\mathcal K_r,
  \label{eq:exact-gamma-kernel}
\end{equation}
where the spaces \(\mathcal K_r\) are defined in
\Cref{eq:kernel-space-kr}.  This identity holds in every characteristic.
Consequently,
\begin{equation}
  \rank(\Gamma)
  =
  \mathsf R_d(m,M,W),
  \label{eq:exact-gamma-rank}
\end{equation}
where
\begin{equation}
  \mathsf R_d(m,M,W)
  :=
  \sum_{r=0}^{m-1}\Lambda_d(W+r)
  \left(
    (r+1)(M+1)
    -\positive{r-h_r+1}\positive{M-h_r+1}
  \right).
  \label{eq:exact-local-rank}
\end{equation}
In particular, for every received point
\((\alpha,y)\),
\begin{equation}
  \rank(\Phi_{\alpha,y})
  \leq
  \mathsf R_d(m,M,W).
  \label{eq:local-rank-bound}
\end{equation}
\end{lemma}

\begin{proof}
All calculations take place in
\(\F_q[T,U,Y_1,\ldots,Y_d]/(T^m)\).  The monomial basis in
\Cref{eq:rank-space-v} gives
\begin{equation}
  \dim\mathcal V
  =
  \sum_{r=0}^{m-1}(r+1)(M+1)\Lambda_d(W+r).
  \label{eq:dimension-v}
\end{equation}

We first check that \(\mathcal K_r\subseteq\mathcal V\).  In a term of the expansion of
\(E^{h_r}\), suppose the variables \(Y_j\), for \(j\geq 2\), contribute the
extra exponent vector \(e\).  They contribute
\(\ell=\omega(e)\) powers of \(T\).  The resulting monomial has \(T\)-degree
\(r+\ell\), \(U\)-degree at most \(r\), \(Y_1\)-degree at most \(M\), and
\(\omega\)-weight at most \(W+r+\ell\).  It therefore lies in \(\mathcal V\),
or vanishes modulo \(T^m\) when \(r+\ell\geq m\).

Next rewrite an element of \(\mathcal K_r\) in the variables used by
\(\Gamma\).  Every resulting monomial contains at least \(T^rE^{h_r}\).
Since \(r+dh_r\geq m\), none of these monomials is retained by \(\Gamma\).
Thus \(\mathcal K_r\subseteq\ker\Gamma\).

The dimension of this kernel subspace is
\begin{equation}
  \dim\mathcal K_r
  =
  \positive{r-h_r+1}\positive{M-h_r+1}\Lambda_d(W+r).
  \label{eq:dimension-kr}
\end{equation}
To verify both this equality and independence across different \(r\), take a
nonzero relation among the \(\mathcal K_r\) and choose its smallest
\(T\)-degree \(r_\star\).  Modulo \(T\),
\(E\equiv U-Y_1\), so the coefficient of \(T^{r_\star}\) contributed by
\(\mathcal K_{r_\star}\) is \((U-Y_1)^{h_{r_\star}}g\) for a nonzero
polynomial \(g\).  This coefficient is nonzero and cannot be cancelled by a
space with larger index.  The same observation shows that multiplication by
\((U-Y_1)^{h_r}\) introduces no dependence within \(\mathcal K_r\).

It follows that the direct sum of the \(\mathcal K_r\)'s lies in
\(\ker\Gamma\).  We prove the reverse inclusion by successively removing
the lowest nonzero \(T\)-coefficient.  Write an arbitrary
\(F\in\mathcal V\) uniquely as
\begin{equation}
  F=\sum_{r=0}^{m-1}T^rf_r(U,Y_1,\ldots,Y_d),
  \label{eq:gamma-t-adic-expansion}
\end{equation}
where every monomial of \(f_r\) has \(U\)-degree at most \(r\),
\(Y_1\)-degree at most \(M\), and \(\omega\)-weight at most \(W+r\) in
\(Y_2,\ldots,Y_d\).

Suppose that \(F\in\ker\Gamma\) is nonzero, and let \(r\) be the least
index for which \(f_r\ne0\).  After rewriting
\[
  U=E+Y_1+\sum_{j=2}^{d}(-1)^{j+1}T^{j-1}Y_j,
\]
the coefficient of \(T^r\) is
\[
  f_r(E+Y_1,Y_1,\ldots,Y_d).
\]
The map \(\Gamma\) retains its coefficients of \(E^b\) exactly when
\(r+db<m\), or equivalently when \(0\leq b<h_r\).  Since
\(F\in\ker\Gamma\), all of these coefficients vanish.  Therefore
\begin{equation}
  E^{h_r}\mid f_r(E+Y_1,Y_1,\ldots,Y_d),
  \qquad
  (U-Y_1)^{h_r}\mid f_r(U,Y_1,\ldots,Y_d).
  \label{eq:gamma-leading-divisibility}
\end{equation}
Write
\(f_r=(U-Y_1)^{h_r}g_r\).  The divisor is monic of degree \(h_r\)
when viewed separately as a polynomial in either \(U\) or \(Y_1\).
Hence
\[
  \deg_U g_r\leq r-h_r,
  \qquad
  \deg_{Y_1}g_r\leq M-h_r.
\]
Multiplication by \((U-Y_1)^{h_r}\) does not mix monomials in
\(Y_2,\ldots,Y_d\). Decomposing coefficientwise in those monomials
therefore shows that every monomial of \(g_r\) has \(\omega\)-weight at
most \(W+r\). Thus
\[
  G_r:=T^rE^{h_r}g_r
\]
belongs to \(\mathcal K_r\), where on the right \(E\) denotes the
polynomial in \Cref{eq:e-in-u}.  Its coefficient of \(T^r\) is exactly
\(f_r\).  Both \(F\) and \(G_r\) lie in \(\ker\Gamma\), so
\(F-G_r\) has strictly larger \(T\)-adic order and remains in the kernel.
Iteration terminates after at most \(m\) steps and expresses \(F\) as an
element of \(\sum_r\mathcal K_r\).  The independence already proved makes
this sum direct, establishing \Cref{eq:exact-gamma-kernel}.

Finally, subtracting the dimensions in \Cref{eq:dimension-kr} from
\Cref{eq:dimension-v} gives
\(\rank(\Gamma)=\mathsf R_d(m,M,W)\).  The factorization
\Cref{eq:local-map-factorization} gives the last assertion.
\end{proof}

The equality in \Cref{eq:exact-gamma-rank} concerns the enlarged,
point-independent map \(\Gamma\).  It does not assert that
\(\Phi_{\alpha,y}\) has the same rank: the first substitution may have
proper image, and
\[
  \rank(\Phi_{\alpha,y})
  =\rank\!\left(\Gamma\big|_{\operatorname{im}\Psi_{\alpha,y}}\right)
  \leq\rank(\Gamma).
\]
The next proposition, new to this work, gives an exact block decomposition of
this restricted map.

\subsection{A block decomposition of the actual local map}

For nonnegative integers \(u,t_0,\ldots,t_d\), define the integer
multinomial coefficient
\begin{equation}
  \operatorname{Mult}(u;t_0,\ldots,t_d)
  :=
  \begin{cases}
    \displaystyle\frac{u!}{t_0!\cdots t_d!},
      &t_0+\cdots+t_d=u,\\[1.2ex]
    0,&\text{otherwise},
  \end{cases}
  \label{eq:local-block-multinomial}
\end{equation}
with value zero as well if any \(t_j\) is negative.

\begin{proposition}[Exact block decomposition of the local rank]
\label{prop:exact-local-block-decomposition}
For a source monomial
\[
  X^xY_0^uY_1^bY^c
  \in\Qspace(D,A,d,m,M,W),
\]
define its two grades by
\begin{equation}
  \ell=u+b+|c|,
  \qquad
  \nu=x-b-\sum_{j=2}^{d}jc_j.
  \label{eq:source-local-grades}
\end{equation}
For an output monomial
\(T^rZ^sY_1^eY^z\), where
\(z=(z_2,\ldots,z_d)\in\N^{d-1}\), define
\begin{equation}
  \ell=s+e+|z|,
  \qquad
  \nu=r-(d+1)s-e-\sum_{j=2}^{d}jz_j.
  \label{eq:target-local-grades}
\end{equation}
For each \((\ell,\nu)\in\N\times\mathbb Z\), let
\(\mathsf C_{\ell,\nu}\) be the matrix whose columns are the source
monomials of grade \((\ell,\nu)\), whose rows are the output monomials of
that grade satisfying
\begin{equation}
  0\leq r<m,
  \qquad
  r\geq (d+1)s,
  \label{eq:local-block-row-range}
\end{equation}
and whose entry in the row
\(T^rZ^sY_1^eY^z\) and column
\(X^xY_0^uY_1^bY^c\) is the reduction in \(\F_q\) of
\begin{equation}
  (-1)^{\sum_{j=2}^{d}(j+1)(z_j-c_j)}
  \operatorname{Mult}
  \left(u;s,e-b,z_2-c_2,\ldots,z_d-c_d\right).
  \label{eq:local-block-entry}
\end{equation}
When \(d=1\), the final list of differences in
\Cref{eq:local-block-entry} is empty.
Then the exact rank of the actual local constraint map is independent of
\((\alpha,y)\) and satisfies
\begin{equation}
  \rank_{\F_q}(\Phi_{\alpha,y})
  =
  \sum_{\ell,\nu}\rank_{\F_q}(\mathsf C_{\ell,\nu}).
  \label{eq:exact-actual-local-rank}
\end{equation}
The sum is finite.  Moreover, a source monomial belongs to
\(\mathsf C_{\ell,\nu}\) only if
\begin{equation}
  D\ell+\nu<mA.
  \label{eq:block-grade-weight}
\end{equation}
\end{proposition}

\begin{proof}
First observe that the rank is independent of the received pair.  The
translation
\[
  Q(X,Y_0,Y_1,\ldots,Y_d)
  \longmapsto
  Q(X+\alpha,Y_0+y,Y_1,\ldots,Y_d)
\]
is an automorphism of \(\Qspace(D,A,d,m,M,W)\): expanding the two translated
variables can only lower their exponents and hence their weighted degree.
Composing the normalized local map at \((0,0)\) with this automorphism gives
the map at \((\alpha,y)\).

At \((0,0)\), the local substitution is
\[
  X=T,
  \qquad
  Y_0=
  \sum_{j=1}^{d}(-1)^{j+1}T^jY_j+TE.
\]
Rows with $i<s$ are identically zero because every occurrence of $E$ in
the local substitution comes with a factor of $T$. Delete these rows and
relabel each remaining target coordinate by replacing \(E\) with \(T^dZ\).
This is injective on the remaining target basis because
\[
  T^iE^sY_1^eY^z
  \longmapsto
  T^{i+ds}Z^sY_1^eY^z,
  \qquad i+ds<m.
\]
Its image consists exactly of the rows with $r\ge(d+1)s$ in
\Cref{eq:local-block-row-range}. Thus it preserves rank and changes the
normalized substitution to
\begin{equation}
  Y_0=
  \sum_{j=1}^{d}(-1)^{j+1}T^jY_j+T^{d+1}Z
  \pmod{T^m}.
  \label{eq:local-block-substitution}
\end{equation}

Expand a source monomial.  To obtain
\(T^rZ^sY_1^eY^z\), the \(u\) factors of \(Y_0\) must contribute
\(s\) copies of \(T^{d+1}Z\), \(e-b\) copies of \(TY_1\), and
\(z_j-c_j\) copies of
\((-1)^{j+1}T^jY_j\) for each \(2\leq j\leq d\).  The coefficient is
exactly \Cref{eq:local-block-entry}, and the resulting exponent of \(T\) is
\begin{equation}
  r=x+(d+1)s+(e-b)+\sum_{j=2}^{d}j(z_j-c_j).
  \label{eq:local-block-t-degree}
\end{equation}
The choice counts also give
\[
  u=s+(e-b)+\sum_{j=2}^{d}(z_j-c_j).
\]
These two identities say precisely that both grades in
\Cref{eq:source-local-grades,eq:target-local-grades} are preserved.
Conversely, when the lower arguments of the multinomial coefficient are
nonnegative, these identities describe every term in the expansion.

Distinct grades have disjoint source bases and disjoint target coordinates.
The matrix of the local map is therefore block diagonal with blocks
\(\mathsf C_{\ell,\nu}\), proving
\Cref{eq:exact-actual-local-rank}.  Finally, the source weighted degree is
\[
  x+Du+(D-1)b+\sum_{j=2}^{d}(D-j)c_j
  =D\ell+\nu,
\]
which proves \Cref{eq:block-grade-weight} and also makes the collection of
nonempty blocks finite.
\end{proof}

The block decomposition is not confined to this interpolation space.  For any
monomial interpolation support, including the weighted support of the main paper
(\Cref{M-thm:weighted-support}), one obtains the exact local matrix by retaining
precisely the columns belonging to that support and deleting empty rows.

Because the block entries are integers, ranks computed over
\(\mathbb Q\) are characteristic-independent upper bounds for the ranks in
\Cref{eq:exact-actual-local-rank}.  A rank computed modulo one prime is a
lower bound for the rational rank, not by itself an upper bound in other
characteristics.  The exact block sum may be strictly smaller than
\(\mathsf R_d(m,M,W)\); the latter remains the convenient closed-form
upper bound used below.  In particular, the exact block sum in
\Cref{eq:exact-actual-local-rank} may replace
\(\mathsf R_d(m,M,W)\) in any finite interpolation certificate over a
specified field.

\subsection{The finite certificate}

We now state a closed-form sufficient interpolation test.  Its left side is
the exact dimension, and its right side is the exact rank of the enlarged
map \(\Gamma\), hence an explicit upper bound on the actual local rank.  The
exact block sum from \Cref{prop:exact-local-block-decomposition} gives a
sharper field-specific test whenever its block rank is smaller. All of these quantities are
finite, so no limit or hidden constant enters the assertion.

\begin{theorem}[Finite interpolation certificate]
\label{thm:finite-interpolation-certificate}
Assume the parameters in \Cref{eq:certificate-parameters} satisfy
\begin{multline}
  \sum_{\substack{c\in\N^{d-1}\\\omega(c)\leq W}}
  \ \sum_{b_1=0}^{M}
  \mathsf N_D\left(
    mA-(D-1)b_1-
    \sum_{j=2}^{d}(D-j)c_j
  \right)
  \\
  >
  n\sum_{r=0}^{m-1}\Lambda_d(W+r)
  \left(
    (r+1)(M+1)
    -\positive{r-h_r+1}\positive{M-h_r+1}
  \right).
  \label{eq:finite-interpolation-certificate}
\end{multline}
Then, for every choice of distinct evaluation points
\(\alpha_1,\ldots,\alpha_n\in\F_q\) and every received word
\(y\in\F_q^n\), there is a nonzero
\(Q\in\Qspace(D,A,d,m,M,W)\) satisfying the local constraints at all received
points.  Every \(P\in\F_q[X]_{\leq D}\) with agreement at least \(A\) then
satisfies
\begin{equation}
  Q(X,P(X),P^{[1]}(X),\ldots,P^{[d]}(X))\equiv 0.
  \label{eq:certificate-differential-equation}
\end{equation}
The interpolant also satisfies the total \(Y\)-degree bound in
\Cref{eq:cap-free-total-y-degree}.
\end{theorem}

\begin{proof}
By \Cref{prop:exact-cap-free-dimension}, the left side of
\Cref{eq:finite-interpolation-certificate} is exactly \(\dim\Qspace\).
By \Cref{lem:exact-local-rank}, each received point imposes rank at most the
sum on the right after removing the factor \(n\).  Finally,
\Cref{lem:cap-free-degree-bounds} supplies the specialization-degree
hypothesis of \Cref{prop:global-interpolation-principle}.  That proposition
gives the nonzero interpolant and \Cref{eq:certificate-differential-equation}.
\end{proof}
 \section{One-sided supports and ambient-dimension estimates}
\label{sec:one-sided-estimates}

This section retains the dimension and reachable-coordinate estimates for
an alternative one-sided interpolation support. They isolate the effect of
ambient padding and weighted-simplex enlargement, and supply context for
\Cref{app:boundary-asymptotics}. The sharper weighted-support certificate
is proved in the main paper; its proof is not repeated here.

Fix a gap \(\delta\in(0,1)\), and define the baseline ambient rate
\[
  \rho_0:=\frac{\delta(1-\delta)}2.
\]
For a code of dimension \(k\) and block length \(n\), set
\begin{equation}
  A:=k+\lceil\delta n\rceil,
  \qquad
  K:=\max\{k,\lceil\rho_0 n\rceil\},
  \qquad
  D:=K-1.
  \label{eq:adaptive-ambient}
\end{equation}
Throughout this section we assume \(A\le n\).

\subsection{Uniform ambient geometry}

Write \(r=k/n\) and
\[
  \rho(r):=\max\{r,\rho_0\},
  \qquad
  \eta(r):=r+\delta.
\]
The factor \(1/2\) in this baseline leaves extra room at the low-rate
endpoint. It is not needed for the qualitative all-rate result. The
weighted-support proof in \Cref{M-sec:optimized-parameters} of the main paper uses the larger
baseline \(\delta/2\), implemented by \(\lfloor\delta n/2\rfloor\), and a
separate endpoint analysis.

\begin{lemma}[Adaptive ambient bounds]
\label{lem:adaptive-ambient-bounds}
Suppose \(A\le n\). Then \(0\le r\le1-\delta\), \(k\le K<A\), and
\begin{equation}
  \rho(r)\ge\frac{\delta(1-\delta)}2,
  \qquad
  \frac{\eta(r)}{\rho(r)}\ge\frac1{1-\delta},
  \qquad
  1-\frac{\rho(r)}{\eta(r)}\ge\delta.
  \label{eq:uniform-ambient-geometry}
\end{equation}
Moreover,
\[
  \frac{A}{D}>\frac{\eta(r)}{\rho(r)}
\]
whenever \(D>0\), and \(D/n\ge\rho_0-1/n\).
\end{lemma}

\begin{proof}
The inequality \(A\le n\) and the definition of \(A\) imply
\(k+\delta n\le n\), hence \(r\le1-\delta\). If \(K=k\), then \(K<A\) because
\(\lceil\delta n\rceil\ge1\). If \(K=\lceil\rho_0n\rceil\), then
\(K\le\lceil\delta n\rceil<A\), where the strict inequality uses \(k\ge1\).

The first inequality in \eqref{eq:uniform-ambient-geometry} is immediate. If \(r\ge\rho_0\), then
\[
  \frac{\eta(r)}{\rho(r)}=\frac{r+\delta}{r}
  \ge \frac{1}{1-\delta},
\]
because \(r\le1-\delta\). If \(r\le\rho_0\), then
\[
  \frac{\eta(r)}{\rho(r)}
  =\frac{2(r+\delta)}{\delta(1-\delta)}
  \ge \frac{2}{1-\delta}
  >\frac{1}{1-\delta}.
\]
The third inequality is equivalent to \(\rho(r)\le(1-\delta)(r+\delta)\). For
\(r\ge\rho_0\), this reduces to \(r\le1-\delta\); for \(r\le\rho_0\), it follows from
\(\rho_0=\delta(1-\delta)/2\le(1-\delta)(r+\delta)\).

Finally, \(A/n\ge\eta(r)\), while \(D/n<\rho(r)\): in the branch \(K=k\), we have
\(D=k-1<rn\), and in the other branch \(D=\lceil\rho_0n\rceil-1<\rho_0n\).
The lower bound on \(D/n\) follows from \(K\ge\lceil\rho_0n\rceil\).
\end{proof}

The additive padding \(K=k+\lfloor\lambda\delta n\rfloor\) also gives uniform separation, but it
unnecessarily increases the ambient degree when \(k\) is already linear in \(n\). The adaptive choice
\eqref{eq:adaptive-ambient} uses padding only near rate zero and makes the elementary uniform bounds
particularly simple.  More generally, every fixed baseline in
\((0,\delta(1-\delta)]\) works qualitatively. We retain the present
half-baseline for this existence proof and for the one-sided optimization in
\Cref{app:boundary-asymptotics}; it is not the ambient choice in the
weighted-support certificate.

\subsection{Dimension and coordinate estimates}

We now compare the source dimension with the number of reachable local coordinates. Unlike BCPZZ's
large-kernel analysis, the proof counts the local output coordinates
directly. This shorter comparison isolates the two asymptotic effects that
matter: the formal remainder exponent has about \(m/d\) choices, while
enlarging the weighted simplex costs only \(d^{1/a}\).

Let
\[
  \gamma:=\frac1{1-\delta},
\]
and choose fixed real numbers \(1<a<b<\gamma\). For an integer \(d\ge3\), put
\[
  m=d^3,
  \qquad
  H=H_{d-1}:=\sum_{j=1}^{d-1}\frac1j,
  \qquad
  W=\left\lfloor\frac{adm}{H}\right\rfloor,
  \qquad
  C=\lceil bm\rceil.
\]
For \(c=(c_2,\ldots,c_d)\in\Z_{\ge0}^{d-1}\), recall
\[
  |c|=\sum_{j=2}^dc_j,
  \qquad
  \omega(c)=\sum_{j=2}^d(j-1)c_j,
\]
and define
\[
  \mathcal B=\{c\in\Z_{\ge0}^{d-1}:\omega(c)\le W,\ |c|\le C\}.
\]
Let \(\Qspace_{\mathrm{coarse}}\) be the span of the monomials
\begin{equation}
  X^xY_0^{b_0}Y_1^{b_1}Y^c,
  \qquad
  c\in\mathcal B,\quad 0\le b_1\le m,\quad x,b_0\in\N,
  \qquad
  x+D(b_0+b_1+|c|)<mA.
  \label{eq:coarse-interpolation-space}
\end{equation}
Every specialization at a polynomial of degree at most \(D\) has degree
less than \(mA\).

The next dimension calculation reformulates the counting argument in
\cite[Lemma~3.2]{BCPZZ26}, with the ambient degree substituted for the actual
message degree and the strict inequalities and rounding retained. It is
governed by the function
\begin{equation}
  f(u):=\frac12\int_0^1\positive{u-z}^{2}\,dz
  =
  \begin{cases}
    0,&u\le0,\\
    u^3/6,&0\le u\le1,\\
    (3u^2-3u+1)/6,&u\ge1.
  \end{cases}
  \label{eq:dimension-function}
\end{equation}

\begin{lemma}[Coarse dimension lower bound]
\label{lem:coarse-dimension}
Suppose \(D>d\). For all sufficiently large \(d\), uniformly over all \(n,k\) with \(A\le n\),
\[
  \dim\Qspace_{\mathrm{coarse}}
  \ge
  |\mathcal B|D m^3
  f\left(\frac{A}{D}-\frac{C}{m}\right).
\]
Consequently, for all sufficiently large \(n\),
\begin{equation}
  \frac{\dim\Qspace_{\mathrm{coarse}}}{n}
  \ge
  \frac{\rho_0}{2}|\mathcal B|m^3
  f\left(\frac{\gamma-b}{2}\right).
  \label{eq:uniform-dimension-lower}
\end{equation}
\end{lemma}

\begin{proof}
For fixed \(c,b_1\), put
\(L=mA-D(|c|+b_1)\). The number of pairs \((x,b_0)\in\Z_{\ge0}^2\) with
\(x+Db_0<L\) is at least \(L_+^2/(2D)\). Since \(|c|\le C\), summing over
\(0\le b_1\le m\) and comparing the decreasing sum with its integral gives
\[
  \dim\Qspace_{\mathrm{coarse}}
  \ge |\mathcal B|\frac D2
  \sum_{b_1=0}^m\positive{mA/D-C-b_1}^{2}
  \ge |\mathcal B|Dm^3f(A/D-C/m).
\]
By \Cref{lem:adaptive-ambient-bounds}, \(A/D>\gamma\). Since \(C/m\to b\) as
\(d\to\infty\), the argument of \(f\) is at least \((\gamma-b)/2\) for large \(d\).
The same lemma gives \(D/n\ge\rho_0/2\) for large \(n\), proving
\eqref{eq:uniform-dimension-lower}.
\end{proof}

Define
\[
  \Lambda(z):=\#\{c\in\Z_{\ge0}^{d-1}:\omega(c)\le z\}.
\]

\begin{lemma}[Weighted-simplex enlargement]
\label{lem:uniform-simplex-enlargement}
For fixed \(1<a<b\) and \(m=d^3\),
\begin{equation}
  \frac{\Lambda(W+m)}{|\mathcal B|}=O_{a,b}(d^{1/a}).
  \label{eq:uniform-simplex-enlargement}
\end{equation}
\end{lemma}

\begin{proof}
Consider the weighted simplex
\[
  S_W=\left\{x\in\mathbb R_{\ge0}^{d-1}:\sum_{j=1}^{d-1}jx_j\le W\right\}.
\]
Its volume is \(W^{d-1}/((d-1)!)^2\), and a uniform point \(X\in S_W\) satisfies
\(\E{|X|}=WH/d\le am\). Markov's inequality shows that the portion with
\(|X|\le C\) has volume at least \((1-a/b)\operatorname{vol}(S_W)\). Taking coordinatewise floors
maps this portion into \(\mathcal B\), so
\[
  |\mathcal B|\ge(1-a/b)\frac{W^{d-1}}{((d-1)!)^2}.
\]
Conversely, the disjoint unit cubes based at the lattice points counted by \(\Lambda(W+m)\) lie in
the simplex of radius \(W+m+\binom d2\). Hence
\[
  \Lambda(W+m)
  \le\frac{(W+m+\binom d2)^{d-1}}{((d-1)!)^2}.
\]
Since \(W=(1+o(1))adm/H\), their ratio is at most
\[
  \frac b{b-a}
  \left(1+\frac{m+\binom d2}{W}\right)^{d-1}
  =O_{a,b}(e^{H/a})=O_{a,b}(d^{1/a}).
\]
\end{proof}

\begin{lemma}[Direct local-coordinate bound]
\label{lem:direct-local-coordinate-bound}
For every received pair \((\alpha,y)\),
\begin{equation}
  \rank\!\left(
    \Phi_{\alpha,y}|_{\Qspace_{\mathrm{coarse}}}
  \right)
  \le
  \left\lceil\frac{mA}{D}\right\rceil
  \sum_{r=0}^{m-1}
  \left\lceil\frac{m-r}{d+1}\right\rceil
  \Lambda(W+r).
  \label{eq:direct-local-coordinate-bound}
\end{equation}
If \(m=d^3\), \(d\ge3\), and \(n\ge2/\rho_0\), then
\begin{equation}
  \rank\!\left(
    \Phi_{\alpha,y}|_{\Qspace_{\mathrm{coarse}}}
  \right)
  \le
  \left(\frac2{\rho_0}+1\right)
  \frac{m^3}{d}\Lambda(W+m).
  \label{eq:direct-local-coordinate-coarse}
\end{equation}
Consequently,
\begin{equation}
  \frac{
    \rank(\Phi_{\alpha,y}|_{\Qspace_{\mathrm{coarse}}})
  }{|\mathcal B|m^3}
  =
  O_{\delta,a,b}\left(d^{-(1-1/a)}\right).
  \label{eq:direct-local-power-saving}
\end{equation}
\end{lemma}

\begin{proof}
Expand a source monomial after the local substitution and consider a
retained output coordinate
\[
  T^iE^hY_1^eY^z.
\]
Constants arising from \(\alpha\) and \(y\) only lower source exponents.
Every coordinate therefore satisfies
\begin{equation}
  i\ge h,\qquad
  \omega(z)\le W+i-h,\qquad
  i+dh<m,\qquad
  h+e+|z|<\frac{mA}{D}.
  \label{eq:coarse-reachable-coordinates}
\end{equation}
Put \(r=i-h\). The second inequality allows at most \(\Lambda(W+r)\)
choices for \(z\). The local truncation becomes
\(r+(d+1)h<m\), so \(h\) has at most
\(\lceil(m-r)/(d+1)\rceil\) possible values. Finally, after fixing
\((r,h,z)\), the last inequality allows at most
\(\lceil mA/D\rceil\) possible values for \(e\). The rank is at most the
number of reachable output coordinates, which proves
\eqref{eq:direct-local-coordinate-bound}.

For \(m=d^3\) and \(d\ge3\),
\[
  \sum_{r=0}^{m-1}
  \left\lceil\frac{m-r}{d+1}\right\rceil
  \le \frac{m(m+1)}{2d}+m
  \le \frac{m^2}{d}.
\]
Moreover, \(D\ge\rho_0n/2\) when \(n\ge2/\rho_0\), while \(A\le n\).
Thus
\[
  \left\lceil\frac{mA}{D}\right\rceil
  \le\left(\frac2{\rho_0}+1\right)m.
\]
Monotonicity of \(\Lambda\) gives
\eqref{eq:direct-local-coordinate-coarse}. Dividing by
\(|\mathcal B|m^3\) and applying
\Cref{lem:uniform-simplex-enlargement} proves
\eqref{eq:direct-local-power-saving}.
\end{proof}

The one-sided weighted support and the comparison between a weighted simplex
and its enlargement are adapted from BCPZZ. The ambient padding and the
direct coordinate count are new.


\begin{corollary}[Existence of a one-sided interpolation certificate]
\label{thm:uniform-interpolation}
For every fixed \(\delta\in(0,1)\), there exist integers
\[
  d=d(\delta),\qquad m=m(\delta),\qquad W=W(\delta),\qquad C=C(\delta),
\]
and \(N=N(\delta)\) such that, for every \(n\ge N\), every
\(1\le k\le n\) with \(A=k+\lceil\delta n\rceil\le n\), and the adaptive
ambient dimension \eqref{eq:adaptive-ambient},
\[
  \dim\Qspace_{\mathrm{coarse}}
  >
  n\,
  \rank\!\left(
    \Phi_{\alpha,y}|_{\Qspace_{\mathrm{coarse}}}
  \right)
\]
for every received pair \((\alpha,y)\). Hence a nonzero common interpolant
exists, and every degree-at-most-\(D\) polynomial with at least \(A\)
agreements satisfies its differential equation. The same derivative order
\(d(\delta)\) works for every \(k\).
\end{corollary}

\begin{proof}
Choose \(1<a<b<1/(1-\delta)\). The constant on the right-hand side of
\eqref{eq:uniform-dimension-lower} is positive. Since \(1-1/a>0\), the right-hand side of
\eqref{eq:direct-local-power-saving} tends to zero as \(d\to\infty\).
Choose one sufficiently large \(d\), set \(m=d^3\), and choose \(W,C\) as
above. For all sufficiently large \(n\), the lower bound for
\(\dim\Qspace_{\mathrm{coarse}}/n\) in
\eqref{eq:uniform-dimension-lower} is strictly larger than the local rank
bound in \eqref{eq:direct-local-power-saving}. Stacking the \(n\) local maps
therefore leaves a nonzero common kernel. The specialization-degree bound
following \eqref{eq:coarse-interpolation-space}, together with
\Cref{prop:global-interpolation-principle}, proves the differential identity.
All remaining requirements, including \(D>d\) and the individual degree bounds needed by root finding,
hold after increasing \(N\), because \(D/n\ge\rho_0-1/n\) and every interpolation parameter is fixed.
\end{proof}
 \section{Small-order certificates and shrinking gaps}
\label{sec:optimized-parameters}

The main manuscript now gives first-order certificates at gaps down to
$0.2195$, together with linear list bounds and quadratic MCA bounds
(\Cref{M-sec:low-order-interpolation}). The certificates below retain
simpler boundary cases and their original exact-rank calculations.

The following fixed-parameter certificates are not needed for the strongest
derivative-order statement. They remain useful when one needs every
interpolation parameter, including the multiplicity and jet degrees, to be
independent of the block length.

\begin{theorem}[Fixed-parameter small-order certificates]
\label{thm:small-derivative-orders}
The following interpolation certificates hold uniformly over the rate.
\begin{enumerate}
  \item Let $n\geq2$, $A_0=k+\lceil n/2\rceil\leq n$, and $D_0=k-1$.  The order-zero space
  \begin{equation}
    \mathcal Q_0=
    \operatorname{span}\{X^uY_0^b:b\in\{0,1\},\ u+D_0b<A_0\}
    \label{eq:order-zero-space}
  \end{equation}
  has dimension $2A_0-D_0>n$, while each local map has rank one.  The agreement list at threshold
  $A_0$ has at most one member.

  \item Let $n\geq512$, $A_1=k+\lceil n/4\rceil\leq n$, and
  \[
    K_1=\max\{k,\lfloor n/16\rfloor\},
    \qquad
    D_1=K_1-1,
    \qquad
    \rho=D_1/n,
    \qquad
    \gamma=A_1/D_1.
  \]
  With derivative order $d=1$, multiplicity $m=64$, and $Y_1$-degree bound $M=16$, let
  \begin{equation}
    \mathcal Q_1=
    \operatorname{span}\{X^uY_0^aY_1^b:
      0\leq b\leq16,\ u+D_1a+(D_1-1)b<64A_1\}.
    \label{eq:order-one-space}
  \end{equation}
  Its local rank is exactly $28152$ in every characteristic, and
  \begin{equation}
    \frac{\dim\mathcal Q_1}{n}
    \geq\rho\sum_{b=0}^{16}\mathsf G(64\gamma-b)
    \geq\frac{585072}{19}>28152.
    \label{eq:order-one-dimension}
  \end{equation}
  Thus a nonzero order-one interpolant exists at every rate.
\end{enumerate}
\end{theorem}

\begin{proof}
For the order-zero assertion,
\[
  \dim\mathcal Q_0=A_0+(A_0-D_0)=2A_0-D_0>n.
\]
The local map is evaluation at one received pair, so its rank is one because the space contains the
constant polynomial.  A nonzero interpolant has the form $Q_0(X)+Q_1(X)Y_0$.  Here $Q_1\neq0$:
otherwise the nonzero polynomial $Q_0$, of degree less than $A_0\leq n$, would vanish at all $n$
distinct evaluation points.  Any listed polynomial must therefore equal $-Q_0/Q_1$, proving uniqueness.

We turn to order one.  First, $\gamma\geq4/3$.  If $K_1=k$, this follows from
$k\leq3n/4$ and $A_1\geq k+n/4$.  If $K_1>k$, then $D_1<n/16$ and
$A_1\geq n/4$, so in fact $\gamma>4$.

Translations in $X$ and $Y_0$ act triangularly on \eqref{eq:order-one-space}, so they preserve the
local rank.  Normalize the received pair to $(0,0)$.  At $T$-degree
$s\in\{0,\ldots,63\}$, the intermediate
space is the rectangle
\[
  \operatorname{span}\{T^sU^aY_1^b:0\leq a\leq s,\ 0\leq b\leq16\}.
\]
Every displayed monomial occurs in the image of the source: use
$X^{s-a}Y_0^aY_1^b$.  Its source weight is at most
\[
  s+(D_1-1)(s+16)\leq79D_1-16<64A_1,
\]
where the last inequality follows from $\gamma\geq4/3$.  The remaining substitution is
$E=U-Y_1$.  At $T$-degree $s$, its kernel on the rectangle consists exactly of the polynomials
divisible by $(U-Y_1)^{64-s}$.  Writing $h=64-s$, the total intermediate dimension is
$17\cdot64\cdot65/2=35360$, while the kernel has dimension
\begin{equation}
  \sum_{h=1}^{16}(65-2h)(17-h)=7208.
  \label{eq:order-one-kernel-sum}
\end{equation}
Thus the local rank is $35360-7208=28152$.  This argument uses divisibility by a monic polynomial,
so it is valid in every characteristic.

It remains to verify the uniform dimension lower bound.  For an integer $L$, let
$\mathsf N_D(L)$ count $(u,a)\in\N^2$ with $u+Da<L$.  Then
\[
  \mathsf N_D(L)=D\mathsf G(L/D).
\]
Consequently,
\[
  \dim\mathcal Q_1
  =\sum_{b=0}^{16}\mathsf N_{D_1}(64A_1-(D_1-1)b)
  \geq D_1\sum_{b=0}^{16}\mathsf G(64\gamma-b),
\]
which is the first inequality in \eqref{eq:order-one-dimension}.

Put $S(z)=\sum_{b=0}^{16}\mathsf G(z-b)$.  For $z\geq16$ and $j=\lfloor z\rfloor$, direct summation
gives
\begin{equation}
  S(z)=17(j-7)z-\frac{17}{2}j(j+1)+680.
  \label{eq:order-one-staircase}
\end{equation}
If $K_1=k$, then $\gamma\geq1+1/(4\rho)$, so it suffices to minimize
$F(\rho)=\rho S(64+16/\rho)$.  On an interval where
$j=\lfloor64+16/\rho\rfloor$ is constant, its derivative has the sign of
\[
  -j^2+127j-816.
\]
The sign is positive for $85\leq j\leq120$ and negative for $j\geq121$.  Hence the minimum is at
$64+16/\rho=121$, or $\rho=16/57$, where
\[
  \rho S(64+16/\rho)=\frac{16}{57}S(121)=\frac{585072}{19}.
\]
If $K_1>k$, then $\rho<1/16$ and $\gamma\geq1/(4\rho)$.  Formula
\eqref{eq:order-one-staircase} shows that $\rho S(16/\rho)$ is decreasing on
$(0,1/16]$, and therefore
\[
  \rho S(16/\rho)\geq\frac1{16}S(256)=\frac{65637}{2}>
  \frac{585072}{19}.
\]
This proves \eqref{eq:order-one-dimension} and the theorem.
\end{proof}

\begin{corollary}[A logarithmically vanishing gap]
\label{cor:logarithmically-vanishing-gap}
Fix $C>3$, let $\delta_n=C/\log n$, and suppose $n\leq q=\Theta(n)$ is prime. For all sufficiently
large $n$, combining the main paper's weighted-support certificate with the
all-root enumeration procedure in \Cref{sec:root-finding} gives decoding
uniformly over the rate at gap $\delta_n$, with
\[
  d=n^{1.5/C+o(1)},
  \qquad
  m=n^{3/C+o(1)}=o(n).
\]
Its list size and running time are at most
\[
  \exp\!\left(n^{1.5/C+o(1)}\right).
\]
\end{corollary}

\begin{proof}
Substitution into \eqref{M-eq:main-derivative-order} of the main paper gives the displayed estimates for $d$ and $m$.
Write $b=\Bjet(\delta_n)$. The parameter formulas give
$m\delta_n^{-2}=n^{3/C+o(1)}$ and
$b\delta_n^{-2}=n^{3/C+o(1)}$. Thus $C>3$ makes both quantities
$o(n)$, so the current finite hypothesis
\[
 n\geq\max\{\lceil m/\delta_n^2\rceil,
              \lceil4b/\delta_n^2\rceil\}
\]
holds eventually.
The interpolation space has $q^{O(d)}$ monomials and constraints.  The differential root-finding
bound of \Cref{sec:root-finding}, with extension degree two, gives list size $q^{2d+O(1)}$.  Since
$q=\Theta(n)$, both quantities are at most
$\exp(n^{1.5/C+o(1)})$; polynomial-time linear algebra in their dense dimensions has the same
asymptotic form.
\end{proof}
 \section{Finite quantitative refinements}
\label{sec:quantitative-refinements}

This section preserves finite constants and comparisons that supplement
the main paper's asymptotic bounds. References prefixed by ``main paper''
point to the corresponding statements in that manuscript.

\subsection{Concrete first-order constants}

The exact support test of main-paper
\Cref{M-prop:tuned-first-order} gives much smaller parameters than the
uniform margin estimate. \Cref{tab:first-order-gap-constants} uses
$n=65{,}536$, $k=\rho n$, $\gcurve=0.01$, and
$A=\lceil n(\afirst(\rho)+\gcurve)\rceil$.
The derivative-weighted support and its row-height test certify every
displayed tuple; main-paper \Cref{M-lem:first-order-factorwise} gives the
list and exceptional-count bounds. Here $M<D=k-1$, so the positive
characteristic requirement is just $p>D$. The bounds hold for every
prescribed evaluation set and received line over such a field.

\begin{table}[H]
\centering\small
\begin{tabular}{c r c r c}
\toprule
$\rho$ & $A$ & $(m,M,B,H)$ & list upper bound & $|\mathcal E|$ upper bound\\
\midrule
$1/16$ & $14{,}950$ & $(20,12,72,736)$ & $71{,}457{,}911$ & $7.576\cdot10^{15}$\\
$1/4$ & $31{,}379$ & $(28,12,53,671)$ & $117{,}645{,}549$ & $2.583\cdot10^{16}$\\
$1/2$ & $45{,}869$ & $(24,6,33,399)$ & $55{,}807{,}550$ & $1.097\cdot10^{16}$\\
\bottomrule
\end{tabular}
\caption{First-order constants at gap $\gcurve=0.01$. All displayed
bounds are rounded upward; the MCA error is at most $|\mathcal E|/q$.
These agreements lie below Johnson, so this table addresses a different
threshold from \Cref{tab:johnson-finite-constants}.}
\label{tab:first-order-gap-constants}
\end{table}

The search used $m\le64$, integers $M$ within four of
$\lfloor\lambdapartial m\rfloor$, and integers $B$ within ten of
$\lceil mA/D\rceil$. It selected these supports using the coarse
factorwise envelope within that finite range. The table reevaluates them
with main-paper \eqref{M-eq:first-order-sharp-mca}, minimizing its threshold
$L$ exactly; it does not assert global optimality. The ordinary-tail
alternative can have smaller exact constants even though the factorwise
estimate has better asymptotic gap dependence.

\subsection{Dominance over full differentiation}

The main paper's ordinary-tail transfer retains the received-curve degree
directly. The next result records its comparison with the coarser envelope
obtained by differentiating through the ordinary tail at an ambient
reconstruction degree.

\begin{corollary}[Dominance over full differentiation]
\label{lem:capped-first-order-transfer}
Let $2\le k\le K\le n$ and retain the remaining hypotheses and notation of
main-paper \Cref{M-lem:first-order-hybrid}. For every actual derivative
degree $0\le d\le\Bpartial$ and every $k\le L\le A$, the curve bound
$E_d^{(\ell)}(L)$, reconstructed through the actual message degree
$D=k-1$, is no larger than the following full-differentiation envelope.
Let $J_1'$ and $F_1'$ be the degree budgets
$D_{\mathrm{joint}}(\Bpartial)$ and
$D_{\mathrm{fiber}}(\Bpartial)$ in that lemma, evaluated at $D'=K-1$.
With $\tau'=2D'-1$, $b_0=\Bjet-\Bpartial$, and
\begin{align*}
 J_0'&=\ell\frac{b_0(b_0+1)}2
       +\BZ(b_0+\tau'b_0^2),&
 F_0'&=\frac{b_0(b_0+1)}2,
\end{align*}
define
\[
 E_{\rm full}^{(K,\ell)}(L)=\BZ
 +\frac{n-L+1}{A-L+1}
       \left(\frac{n-k+1}{A-k+1}J_1'+J_0'\right)
 +\ell(n-L)\left(\frac{n-k+1}{L-k+1}F_1'+F_0'\right).
\]
Then $E_d^{(\ell)}(L)\le E_{\rm full}^{(K,\ell)}(L)$ for every
$d,L$, and consequently
\begin{equation}
 \max_{0\le d\le\Bpartial}\min_{k\le L\le A}E_d^{(\ell)}(L)
 \le \min_{k\le L\le A}E_{\rm full}^{(K,\ell)}(L),
 \label{eq:capped-first-order-transfer}
\end{equation}
so every certificate obtained from that envelope remains valid.
\end{corollary}
\begin{proof}
First compare the terminal tails at $D=k-1$ and put $\tau=2D-1$.
Full differentiation through an ordinary tail of degree $b$ has fiber
and joint degrees
\[
 F_{\rm old}=\frac{b(b+1)}2,\qquad
 J_{\rm old}=\ell\frac{b(b+1)}2+\BZ(b+\tau b^2).
\]
The ordinary transfer instead has $F_{\rm new}=b$ and
$J_{\rm new}=\ell b+\BZ\Psi_D(b)$. For $b\ge1$,
\[
 b+\tau b^2-\Psi_D(b)-2(b-1)
 =(b-1)(\tau(b-1)-1)-2(b-2D-1)_+\ge0.
\]
Thus the joint-degree saving, whose incidence multiplier is at least one,
pays the extra $2(b-1)\BZ$ preliminary exclusions. The fiber term also
decreases. Minimizing the new tail over its independent threshold can only
improve it. For $b=0$, both tails cost $\BZ$.

Add the actual order-one stages. Promoting a missing order-one stage only
increases the bound because $D_{\mathrm{fiber}}(j,r;b,c)\ge j$ and
$2bc-c^2\ge b$, and all order-one incidence factors are at least their
order-zero counterparts. Finally increase the reconstruction degree from
$D=k-1$ to $K-1$. Both arguments in the definition of $c$, as well as $b$
and $\ell+(2D-1)\BZ$, are nondecreasing in $D$; so are $rb+(j-r)c$ and
$2bc-c^2$. This proves the pointwise comparison for every $d,L$. Taking
first the minimum over $L$ and then the maximum over the possible actual
degree gives the displayed inequality. If $k=K=n$, the full-code endpoint
was handled directly in main-paper \Cref{M-lem:first-order-hybrid}.
\end{proof}

\subsection{Johnson constants and integer tuning}

The two list estimates in main-paper \eqref{M-eq:johnson-list-bound} are
classical. Writing $x=\sqrt{\rho_-}$, the pairwise-agreement Johnson
argument gives
\[
 |\List_D(y,A)|\le\frac{1-x}{2\gJohnson}+\frac{1+x}{4x}.
\]
Its leading constant is $1-x$ times that of the main paper's closed
interpolation recipe.

Retain $a$, $\rho_-$, $t$, and $B_{\rm cf}$ from the closed bound in
main-paper \eqref{M-eq:johnson-characteristic-free-bound}. Let
$m_{\rm B}=\max\{\lceil\sqrt{\rho_-}/\gJohnson\rceil,3\}$ and
$t_{\rm B}=m_{\rm B}+1/2$. BCHKS give the following bound
\cite[full version, Theorem~4.6]{BCHKS25}:
\[
 B_{\rm BCHKS}
 =\frac{2t_{\rm B}^5+3t_{\rm B}(1-a)\rho_-}
        {3\rho_-^{3/2}}n
  +\frac{t_{\rm B}}{\sqrt{\rho_-}}.
\]
Since $t\le t_{\rm B}$, $t_{\rm B}\ge7/2$, and $0<\rho_-<1$, comparison
with the positive leading term gives
\[
 \frac{B_{\rm cf}}{B_{\rm BCHKS}}
 <\frac{(8/3)nt^3/\rho_-}{(2/3)nt_{\rm B}^5/\rho_-^{3/2}}
 \le\frac{4\sqrt{\rho_-}}{t_{\rm B}^2}<\frac{16}{49}.
\]
This compares the printed real bounds before truncation by the field size;
either probability bound can be vacuous when $q$ is too small.

The closed recipe in the main paper need not minimize the exceptional
count. Retaining the Guruswami--Sudan constraints and BCHKS weighted
challenge count~\cite{GS99,BCHKS25}, we instead choose integer cutoffs
directly.

\begin{proposition}[Exact ordinary interpolation certificate]
\label{prop:johnson-finite-certificate}
Under the code and agreement hypotheses of main-paper
\Cref{M-thm:quantitative-johnson-mca}, choose integers $m\ge1$ and
$1\le B\le\lfloor(mA-1)/D\rfloor$. Put $u=\min\{B,m-1\}$ and
\[
 N=\sum_{j=0}^B(mA-Dj),\quad W=\sum_{j=0}^B j(mA-Dj),\quad
 R=\sum_{b=0}^u(m-b),\quad T=\sum_{b=0}^u b(m-b).
\]
If $N>nR$, set
\begin{equation}
 H=\max\left\{B,\left\lfloor\frac{W-nT}{N-nR}\right\rfloor\right\}.
 \label{eq:johnson-finite-height}
\end{equation}
There is a symbolic explainer with jet degree at most $B$, challenge
degree at most $H$, and weighted specialization degree less than $mA$.
It is nonzero at every challenge and sound after every field extension.
Consequently the list bound is at most $B$, and main-paper
\eqref{M-eq:ordinary-transfer-bound} applies with
$(\Bjet,\BZ)=(B,H)$.
\end{proposition}
\begin{proof}
Give the coefficient of each base monomial $X^iY^j$, with $i+Dj<mA$
and $j\le B$, challenge degree at most $H-j$. A multiplicity output row
of $Y$-degree $b$ and $X$-degree less than $m-b$ has entries of challenge
degree at most $j-b$, with zero entries when $j<b$.
Apply main-paper \Cref{M-lem:polynomial-kernel-height} with column cost
$j$ and row cost $b$. Since $H\ge B$, its unknown and equation counts are
$(H+1)N-W$ and $n((H+1)R-T)$. The chosen height gives the strict surplus
\[
 (H+1)(N-nR)>W-nT.
\]
The lemma supplies a primitive kernel vector, hence an interpolant
nonzero at every challenge. The ordinary multiplicity argument of
main-paper \Cref{M-prop:johnson-symbolic-interpolant} proves the explaining
identity.
\end{proof}

To identify the leading constant, write $x=\sqrt{D/n}$,
$\widehat a=A/n$, and $\widehat\eta=\widehat a-x$. For the full cutoff
$B=\lfloor(mA-1)/D\rfloor$, the exact sums give
\begin{align*}
 (N-nR)/n
 &=\tfrac12m^2(\widehat a^2/x^2-1)
       -\tfrac12m(1-\widehat a)+O_x(1),\\
 (W-nT)/n
 &=\tfrac16m^3(\widehat a^3/x^4-1)+O_x(m).
\end{align*}
At fixed reduced rate and small gap, choose
$m\sim cx(1-x)/(2\widehat\eta)$ for fixed $c>1$.
Then the surplus is positive and \eqref{eq:johnson-finite-height} gives
\[
 B\sim\frac{c(1-x)}{2\widehat\eta},\qquad
 H\sim\frac{c^2x(1-x)^2}{12(c-1)\widehat\eta^2}.
\]
The leading transfer term is proportional to $BH$; minimizing
$c^3/(c-1)$ at $c=3/2$ gives
\[
 BH\sim\frac{9x(1-x)^3}{32\widehat\eta^3},\qquad
 \frac{BH}{1/(24\widehat\eta^3)}
 =\frac{27}{4}x(1-x)^3+o(1)\le\frac{729}{1024}+o(1).
\]
The denominator is the closed recipe's leading product. Thus tuning
improves the constant at every fixed rate, while the inverse-gap exponent
remains three; finite claims use the integer sums above.

\Cref{tab:johnson-finite-constants} gives exact examples for
$n=65{,}536$, $\gJohnson=0.01$, and
$A=\lceil n(\sqrt{D/n}+0.01)\rceil$. The bounded search checks
$1\le m\le4m_0$ and every permitted $B$, where $m_0$ is the closed
multiplicity; it makes no global optimality claim. The MCA columns use the
unified transfer and are rounded down to valid integer bounds; the list
column uses main-paper \eqref{M-eq:johnson-list-bound}.
\begin{table}[ht]
\centering\small
\begin{tabular}{c r c r r r}
\toprule
$k/n$ & $A$ & selected $(m,B,H)$ & list & closed $|\mathcal E|$ & tuned $|\mathcal E|$\\
\midrule
$1/16$ & $17{,}038$ & $(14,57,568)$ & 38 & $4{,}046{,}431{,}259$ & $2{,}498{,}629{,}121$\\
$1/4$ & $33{,}423$ & $(18,36,504)$ & 25 & $8{,}278{,}798{,}980$ & $3{,}383{,}852{,}708$\\
$1/2$ & $46{,}996$ & $(15,21,234)$ & 15 & $13{,}537{,}540{,}252$ & $1{,}448{,}631{,}664$\\
\bottomrule
\end{tabular}
\caption{Finite Johnson constants. Divide exceptional counts by $q$
for the MCA error probability.}
\label{tab:johnson-finite-constants}
\end{table}
 \subsection{Regular lifting and the separant chain}
\label{sec:root-finding}

We count the solutions of a differential equation that agree with the
received word, using Kopparty's reconstruction from a regular initial Taylor
jet~\cite[Theorem~4.3, Theorem~4.4, and Corollary~4.5]{Kopparty15}.
Below the characteristic, the required binomial coefficients are
nonzero, so such a jet has at most one lift. An explicit derivative
chain controls singular solutions; coefficient-space incidence then
gives a list bound independent of the field size.

Recall the Hasse-derivative convention of
\Cref{def:hasse-derivative} and the weighted degree of
\Cref{def:weighted-degree}. For nonzero
\(Q\in\F_q[X,Y_0,\ldots,Y_r]\), write
\[
  \deg_Y Q:=\deg_{Y_0,\ldots,Y_r}Q,
  \qquad
  \ell_D(Q):=\wtdeg_{1,D,D-1,\ldots,D-r}Q.
\]
By \eqref{eq:weighted-specialization-degree}, for every
\(P\in\F_q[X]_{\le D}\),
\begin{equation}
  \deg_X Q(X,P,P^{[1]},\ldots,P^{[r]})\le \ell_D(Q).
  \label{eq:root-specialization-degree}
\end{equation}

\subsection{Unique lifting below the characteristic}

We first show why regular initial Taylor data determine at most one
solution. Below the characteristic, each new coefficient enters the
recurrence with a nonzero multiplier, so we can recover it uniquely.

\begin{lemma}[Unique lift from a regular jet]
\label{lem:unique-regular-lift}
Let \(E\) be a field of characteristic zero or \(p>D\), let \(0\le s\le D\), and let
\[
  R\in E[X,Y_0,\ldots,Y_s].
\]
Fix \(\alpha,\beta_0,\ldots,\beta_s\in E\) such that
\begin{equation}
  R(\alpha,\beta_0,\ldots,\beta_s)=0
  \quad\text{and}\quad
  \frac{\partial R}{\partial Y_s}
  (\alpha,\beta_0,\ldots,\beta_s)\ne0.
  \label{eq:regular-initial-jet}
\end{equation}
There is at most one \(P\in E[X]_{\le D}\) satisfying
\begin{equation}
  R(X,P,P^{[1]},\ldots,P^{[s]})\equiv0
  \label{eq:regular-differential-equation}
\end{equation}
and \(P^{[j]}(\alpha)=\beta_j\) for every \(0\le j\le s\).
Moreover, this polynomial can be reconstructed, or its nonexistence
detected, using a number of field operations polynomial in \(D\),
\(\ell_D(R)\), and the sparse support size and binary exponent lengths of \(R\).
Here coefficients in the arbitrary field $E$ are accessed through field
operations and equality tests; no bit encoding of $E$ is assumed.
\end{lemma}

\begin{proof}
Write $P(\alpha+T)=\sum_{i=0}^Dp_iT^i$, with $p_i=\beta_i$ for
$i\le s$. For $s<i\le D$, changing $p_i$ by $\gamma$ changes the
$j$-th Hasse derivative by
\begin{equation}
 \gamma\binom ijT^{i-j}.
 \label{eq:hasse-lift-perturbation}
\end{equation}
The first affected coefficient in the differential equation is at
$T^{i-s}$, and its change is exactly
\begin{equation}
 \gamma\binom is\frac{\partial R}{\partial Y_s}
 (\alpha,\beta_0,\ldots,\beta_s).
 \label{eq:unique-lift-coefficient}
\end{equation}
Terms with two changed factors have order at least $2(i-s)>i-s$;
a lower derivative or a coefficient $p_j$ with $j>i$ affects only
higher orders. The displayed multiplier is nonzero by regularity and
the characteristic hypothesis. Solving this linear equation successively
therefore fixes every remaining coefficient. Truncated polynomial
arithmetic and a final check of \eqref{eq:regular-differential-equation}
give the asserted reconstruction or rejection and its polynomial cost.
\end{proof}

\subsection{Counting and enumerating solutions}

A solution can be singular for the original equation yet regular for a
later derivative. We use one deterministic chain to find that stage.
Fix a nonzero $Q\in F[X,Y_0,\ldots,Y_r]$ whose individual jet degrees
are smaller than the positive characteristic, if any.
Starting from $Q_0=Q$, form a deterministic derivative chain. If $Q_i$
depends on a jet variable, let
\begin{equation}
  s_i:=\max\{j:Q_i\text{ depends on }Y_j\},
  \qquad
  Q_{i+1}:=\frac{\partial Q_i}{\partial Y_{s_i}}.
  \label{eq:root-derivative-chain}
\end{equation}
Every member is nonzero by the degree hypothesis, and each derivative
lowers total jet degree. The chain therefore ends at a nonzero
$Q_t\in F[X]$.

Every solution has a first nonzero specialization along this chain;
the preceding equation vanishes and its separant does not.
Under the hypotheses of \Cref{lem:unique-regular-lift}, a regular initial
jet determines at most one solution. Enumerating all jets counts all
differential roots; restricting them by received-word agreement gives
the field-independent lists needed for decoding.
 For the algorithmic statements below, \(\operatorname{size}(Q)\) denotes the bit length of the sparse
coefficient-list representation of \(Q\), with exponents written in binary.

Fix a finite base field \(\F_q\), a differential polynomial \(Q\), and
an ambient degree \(D\), and write
\begin{equation}
  \mathcal R_D(Q)
  :=
  \left\{
    P\in\F_q[X]_{\le D}:
    Q(X,P,P^{[1]},\ldots,P^{[r]})\equiv0
  \right\}.
  \label{eq:differential-root-set}
\end{equation}

Define its instance-sensitive singularity degree by
\begin{equation}
  H_{\rm ch}(Q):=
  \max_{1\le i\le t}\ell_D(Q_i),
  \label{eq:chain-singularity-degree}
\end{equation}
with value zero when $t=0$. The ambient degree $D$ is suppressed in the
notation $H_{\rm ch}$. Differentiating once lowers weighted degree by
at least $D-r$, so
\begin{equation}
  H_{\rm ch}(Q)
  \le \max\{0,\ell_D(Q)-(D-r)\}.
  \label{eq:chain-singularity-generic-bound}
\end{equation}
More precisely, if \(s=\max\{j:Q\text{ depends on }Y_j\}\), then all
later members of the chain have no larger weighted degree, and therefore
\begin{equation}
  H_{\rm ch}(Q)=
  \ell_D\!\left(\frac{\partial Q}{\partial Y_s}\right).
  \label{eq:chain-singularity-first-separant}
\end{equation}


\begin{theorem}[Instance-sensitive refinement of Kopparty's root bound]
\label{thm:root-bound}
Let $E_0=\F_q$ have characteristic $p$, let $0\le r\le D<p$, and let
\[
  Q\in E_0[X,Y_0,\ldots,Y_r]\setminus\{0\}
\]
have individual $Y_j$-degrees smaller than $p$, and put $\Delta=\deg_YQ$.
Let $E/E_0$ be an extension of size $S=q^e$, and
assume $S>H_{\rm ch}(Q)$. Then
\begin{equation}
  |\mathcal R_D(Q)|
  \le
  \frac{S}{S-H_{\rm ch}(Q)}\,\Delta S^r.
  \label{eq:chain-root-bound}
\end{equation}
All solutions can be found deterministically in
\[
  (\Delta+1)S^{r+2}
  \poly\bigl(\operatorname{size}(Q),D,\ell_D(Q),e,\log q\bigr)
\]
bit operations.
\end{theorem}

\begin{proof}
If $Q$ is independent of the jet variables, it is a nonzero polynomial in
$X$ and has no solution. Otherwise the chain terminates at a nonzero
univariate polynomial. For a solution $P$, let $j(P)\ge1$ be the first
index such that
\[
  Q_{j(P)}(X,P,P^{[1]},\ldots,P^{[r]})\not\equiv0.
\]
The preceding specialization is identically zero. The nonzero one has degree
at most $H_{\rm ch}(Q)$, and is therefore nonzero at at least
$S-H_{\rm ch}(Q)$ values $\alpha\in E$. Call such an $\alpha$ a regular
witness for $P$.

Map a regular pair $(P,\alpha)$ to
\[
  \bigl(\alpha,P(\alpha),P^{[1]}(\alpha),\ldots,P^{[r]}(\alpha)\bigr)
  \in E^{r+2}.
\]
This map is injective. Indeed, two solutions with different values of
$j(P)$ contradict the first nonzero chain evaluation at their common jet.
If the indices agree, put $F=Q_{j(P)-1}$ and $s=s_{j(P)-1}$. The two
polynomials solve the same equation $F=0$, have the same order-$r$ jet,
and satisfy
\[
  \frac{\partial F}{\partial Y_s}
  (\alpha,P(\alpha),\ldots,P^{[s]}(\alpha))\ne0.
\]
They coincide by \Cref{lem:unique-regular-lift}.

Fix $\alpha\in E$.
If $Q(\alpha,Y_0,\ldots,Y_r)=0$ as a polynomial, specialization of $X$
commutes with differentiation in every $Y_j$. Inductively, every
$Q_i(\alpha,Y_0,\ldots,Y_r)$ is therefore zero.
No solution has a regular witness at this $\alpha$.
Otherwise, $Q(\alpha,Y_0,\ldots,Y_r)$ is a nonzero polynomial of total
degree at most $\Delta$. Every witness jet is a zero of this polynomial.
By the Schwartz--Zippel lemma~\cite{Schwartz80}, applied with
$h=r+1$ variables and total-degree bound $b=\Delta$, there are at most
$\Delta S^r$ such jets. Thus there are at most $\Delta S^{r+1}$ regular
witnesses, while each solution has at least $S-H_{\rm ch}(Q)$. Double
counting proves \eqref{eq:chain-root-bound}.

For the algorithm, enumerate $\alpha\in E$ and every jet
$(\beta_0,\ldots,\beta_r)\in E^{r+1}$. Discard a pair unless
$Q_0(\alpha,\boldsymbol\beta)=0$ and some later chain value is nonzero.
Let $j\ge1$ be its first nonzero chain index, set $F=Q_{j-1}$ and
$s=s_{j-1}$, and use \Cref{lem:unique-regular-lift} with
$(\alpha,\beta_0,\ldots,\beta_s)$ to reconstruct a degree-at-most-$D$
solution of $F=0$, if one exists. Check all prescribed order-$r$ jet values,
the original identity $Q=0$, and membership of every coefficient in $E_0$.
Every true solution has a regular witness, so it is reconstructed at least
once. Sort coefficient vectors and remove duplicates. There are $S^{r+2}$ point-jet pairs, and every
pair requires evaluating at most $\Delta+1$ members of the derivative
chain. Every remaining operation is polynomial in the displayed
parameters.
\end{proof}
 \begin{remark}[The exponent for a single equation]
\label{rem:root-exponent-optimal}
The equation $Q=Y_r$ has exactly $q^r$ solutions of degree less than $r$,
namely all polynomials of degree at most $r-1$. Thus no bound for a generic
single order-$r$ differential equation can improve the base-field exponent
below $r$.
\end{remark}

\begin{theorem}[Generic stratified differential root bound]
\label{thm:root-bound-stratified}
Let \(q\) be prime, let \(0\le r\le D<q\), and let
\[
  Q\in\F_q[X,Y_0,\ldots,Y_r]\setminus\{0\}.
\]
Put
\begin{equation}
  \Delta:=\deg_Y Q,
  \qquad
  \ell:=\ell_D(Q),
  \qquad
  h:=\max\{0,\ell-(D-r)\},
  \label{eq:root-bound-parameters}
\end{equation}
and assume \(\Delta<q\).  Let \(e\ge1\), set \(S=q^e\), and suppose \(S>h\).
Then
\begin{equation}
  |\mathcal R_D(Q)|
  \le
  \frac{S}{S-h}\,\Delta S^r.
  \label{eq:sharpened-root-bound}
\end{equation}
\end{theorem}

\begin{proof}
If \(\Delta=0\), then \(Q\) is a nonzero polynomial in \(X\), so the root
set is empty.  Assume henceforth that \(\Delta>0\).

Set \(E:=\F_S\).  We first bound all solutions in \(E[X]_{\le D}\); the
desired \(\F_q\)-solutions form a subset.  Consider any nonzero iterated
partial derivative \(R\) of \(Q\), allowing \(R=Q\).  If \(R\) is
independent of all \(Y\)-variables, then \(R\in E[X]\setminus\{0\}\) and has
no solutions.  Otherwise, let \(s\) be the largest index such that \(R\)
depends on \(Y_s\).  No assumption is made that \(R\) depends on the
intermediate variables.  Put
\[
  t:=\deg_{Y_s}R,
  \qquad
  \delta_R:=\deg_YR,
\]
and define
\begin{equation}
  R_a:=\frac{\partial^aR}{\partial Y_s^a},
  \qquad 0\le a\le t.
  \label{eq:partial-derivative-chain}
\end{equation}
Since \(t\le\delta_R\le\Delta<q\), the polynomial
\begin{equation}
  R_t=t!\operatorname{LC}_{Y_s}(R)
  \label{eq:terminal-partial-derivative}
\end{equation}
is nonzero and does not depend on \(Y_s\).

Write \(\mathcal Z_D(R)\) for the set of \(P\in E[X]_{\le D}\) satisfying
\[
  R(X,P,\ldots,P^{[s]})\equiv0,
\]
where variables absent from \(R\) are simply ignored.  These solutions split
into regular strata and a residual singular stratum.  For \(0\le a<t\), the
\(a\)-th regular stratum consists of those \(P\in\mathcal Z_D(R)\) for which
\begin{equation}
  R_b(X,P,\ldots,P^{[s]})\equiv0
  \quad(0\le b\le a),
  \qquad
  R_{a+1}(X,P,\ldots,P^{[s]})\not\equiv0.
  \label{eq:regular-stratum}
\end{equation}
The remaining solutions make every \(R_1,\ldots,R_t\) vanish identically
and therefore belong to \(\mathcal Z_D(R_t)\).  Since \(R_t\) has active
order strictly below \(s\), or is independent of all \(Y\)-variables, this
is the residual equation in the recursion.

Fix a solution \(P\) in a regular stratum and write
\[
  H_P(X):=R_{a+1}(X,P,\ldots,P^{[s]}).
\]
Every such separant is a nonzero positive-order partial derivative of the
original \(Q\).  Differentiating with respect to \(Y_j\) lowers weighted
degree by at least \(D-j\ge D-r\).  Together with
\eqref{eq:root-specialization-degree}, this gives
\begin{equation}
  \deg H_P\le \ell-(D-r)\le h.
  \label{eq:separant-degree-bound}
\end{equation}
Thus \(H_P(\alpha)\ne0\) for at least \(S-h\) values \(\alpha\in E\).

We next count the regular solution-point pairs \((P,\alpha)\), jointly
across all regular strata.  Fix \(\alpha\in E\) and
\(\beta_0,\ldots,\beta_{s-1}\in E\), and consider the univariate polynomial
\begin{equation}
  f(Z):=R(\alpha,\beta_0,\ldots,\beta_{s-1},Z).
  \label{eq:fiber-polynomial}
\end{equation}
For a pair in the \(a\)-th stratum, with
\(z=P^{[s]}(\alpha)\), condition \eqref{eq:regular-stratum} says that
\[
  f(z)=f'(z)=\cdots=f^{(a)}(z)=0,
  \qquad
  f^{(a+1)}(z)\ne0.
\]
Since \(\deg f\le t<q\), these are precisely the roots of \(f\) classified
by their exact multiplicity.  Across all \(a<t\), there are at most \(t\)
possible values of \(z\).  For each resulting regular initial jet,
\Cref{lem:unique-regular-lift}, applied to the equation \(R_a=0\), gives at
most one degree-at-most-\(D\) solution.  Hence the number of regular
solution-point pairs is at most
\begin{equation}
  S\cdot tS^s.
  \label{eq:regular-pair-upper-bound}
\end{equation}
On the other hand, every regular solution contributes at least \(S-h\)
pairs.  The number of regular solutions is therefore at most
\begin{equation}
  \frac{S}{S-h}\,tS^s.
  \label{eq:regular-solution-bound}
\end{equation}

It remains to account for the residual equation \(R_t=0\).  Its total
\(Y\)-degree satisfies
\begin{equation}
  \deg_YR_t\le\delta_R-t.
  \label{eq:residual-total-degree}
\end{equation}
The residual \(R_t\) is again an iterated partial derivative of \(Q\), so
the same separant-degree bound applies at every lower active order.
If \(s=0\), then \(R_t\) has no \(Y\)-variables and contributes no residual
solutions, so \eqref{eq:regular-solution-bound} already gives the desired
bound.  Suppose \(s\ge1\).  Induction on the active order, with the base
case of a nonzero polynomial in \(E[X]\), now gives
\begin{align}
  |\mathcal Z_D(R)|
  &\le
  \frac{S}{S-h}
  \left(tS^s+(\delta_R-t)S^{s-1}\right) \\
  &\le
  \frac{S}{S-h}\,\delta_RS^s.
  \label{eq:root-bound-induction}
\end{align}
If the active order of \(R_t\) is \(s'<s\), the induction
hypothesis gives the residual contribution
\(\frac{S}{S-h}(\delta_R-t)S^{s'}\), which is no larger than the term
displayed above; if \(R_t\) has no \(Y\)-variables, its residual contribution
is zero.  Applying \eqref{eq:root-bound-induction} to \(R=Q\), whose active
order is at most \(r\), proves \eqref{eq:sharpened-root-bound}.
\end{proof}

The preceding proof is constructive.  We state the resulting algorithm
separately because later sections use its running time as well as its list
bound.

\begin{corollary}[Constructive root enumeration]
\label{cor:root-enumeration}
Under the hypotheses of \Cref{thm:root-bound-stratified}, one can enumerate
\(\mathcal R_D(Q)\) deterministically in
\begin{equation}
  (\Delta+1)S^{r+2}
  \poly\bigl(\operatorname{size}(Q),D,\ell,e,\log q\bigr)
  \label{eq:root-enumeration-runtime}
\end{equation}
bit operations. In particular, if \(e,r,\Delta\) are fixed and
\(\operatorname{size}(Q),D,\ell\le q^{O(1)}\), the running time is
\(q^{e(r+2)+O(1)}\).
\end{corollary}

\begin{proof}
If \(\Delta=0\), then \(Q\) is a nonzero polynomial in \(X\), so the root set is empty and can be returned
after reading \(Q\).  Assume \(\Delta>0\).
Follow the residual chain \eqref{eq:terminal-partial-derivative}.  At each
polynomial \(R\) and each regular stratum, enumerate \(\alpha\in E\) and all
initial jets in \(E^{s+1}\), retain the regular jets, and apply the
reconstruction in \Cref{lem:unique-regular-lift} to the equation \(R_a=0\).
Verify the original polynomial identity, retain only polynomials with
coefficients in \(\F_q\), and remove duplicates.  The degrees \(t\)
encountered along the residual chain sum to at most \(\Delta\), so there are
at most \(\Delta\) regular strata.  Each stratum enumerates at most
\(S^{r+2}\) point-jet pairs.  Arithmetic and all identity checks take time
polynomial in the displayed parameters.  A representation of \(E\) can be
constructed deterministically within the same bound, for example by
enumerating monic degree-\(e\) polynomials and testing irreducibility.
\end{proof}

\subsection{Specialization to the interpolant}

We finally record the form used by the decoder.  It exposes the two
field-size regimes without repeating the interpolation proof.

\begin{corollary}[Root bound for the hidden-derivative interpolant]
\label{cor:interpolant-root-bound}
Let \(D>d\), and let
\[
  Q\in\F_q[X,Y_0,\ldots,Y_d]\setminus\{0\}
\]
be an interpolant satisfying
\begin{equation}
  L_Q:=\ell_D(Q)\le mA-1.
  \label{eq:interpolant-root-weight}
\end{equation}
Put
\begin{equation}
  \Delta_Q:=\deg_YQ,
  \qquad
  H_Q:=\max\{0,L_Q-(D-d)\}.
  \label{eq:interpolant-root-parameters}
\end{equation}
Assume \(q\) is prime, \(d\le D<q\), and \(\Delta_Q<q\).  For every
\(e\ge1\) with \(q^e>H_Q\),
\begin{equation}
  |\mathcal R_D(Q)|
  \le
  \frac{q^e}{q^e-H_Q}\,\Delta_Qq^{ed}.
  \label{eq:interpolant-root-list-bound}
\end{equation}
Moreover,
\begin{equation}
  \Delta_Q
  \le
  \left\lfloor\frac{mA-1}{D-d}\right\rfloor.
  \label{eq:interpolant-total-degree-bound}
\end{equation}
If \(D=K-1\), then
\begin{equation}
  H_Q\le\max\{0,mA-K+d\}.
  \label{eq:interpolant-separant-bound}
\end{equation}
In particular, \(q^e\ge2H_Q\) implies
\begin{equation}
  |\mathcal R_D(Q)|\le2\Delta_Qq^{ed}.
  \label{eq:interpolant-root-list-bound-two}
\end{equation}
The solutions can be enumerated within the running time of
\Cref{cor:root-enumeration}.
\end{corollary}

\begin{proof}
Every \(Y_j\) has weight at least \(D-d\).  Therefore the weighted-degree
bound \eqref{eq:interpolant-root-weight} implies
\eqref{eq:interpolant-total-degree-bound}.  The remaining claims follow from
\Cref{thm:root-bound}, using
\[
  L_Q-(D-d)\le mA-1-(D-d)=mA-K+d
\]
when \(D=K-1\).
\end{proof}

\begin{remark}[Characteristic and source conventions]
\label{rem:kopparty-source-conventions}
Kopparty's Theorem~4.3 is stated with
\(\deg_{Y_j}Q<q\)~\cite{Kopparty15}.  Its surrounding application is over a
prime-order field.  Read over a field of characteristic \(p\), the recursive
differentiation step requires \(\deg_{Y_j}Q<p\), since it uses
\(t!\operatorname{LC}_{Y_j}(Q)\ne0\); the unique-lifting step likewise
requires \(D<p\).  Our assumptions that \(q\) is prime and
\(D,\Delta<q\) make both requirements explicit.  The restriction on \(D\)
is substantive: over \(\F_q\), the equation \(P^{[1]}=0\) has the
\(q^2\) solutions \(P=aX^q+b\) of degree at most \(q\).  More generally,
in characteristic \(p\), every polynomial \(P(X)=H(X^p)\) satisfies
\(P^{[1]}=0\).  Once \(D\ge p\), the initial-jet uniqueness used above
therefore fails.  There is also a
harmless typographical slip in the proof on p.~167: the nonvanishing
displayed for \(Q\) should instead be assigned to the separant \(H\), while
\(Q(\alpha,P(\alpha),\ldots)=0\) follows from the polynomial identity.
The algorithm and its induction use these intended relations.
\end{remark}
 \subsection{Primitive normalization}
\label{sec:primitive-normalization}

An optional preprocessing step can improve instance-specific parameters.
View \(Q\) as a polynomial in the jet variables with
coefficients in \(\F_q[X]\), let \(g(X)\) be the greatest common divisor of
its nonzero coefficients, and put \(Q^{\rm prim}=Q/g\). Since
\(\F_q[X]\) is an integral domain,
\begin{equation}
  \mathcal R_D(Q)=\mathcal R_D(Q^{\rm prim}).
  \label{eq:primitive-part-same-roots}
\end{equation}
Passing to the primitive part cannot increase the jet degrees or weighted
degree. The theorem below applies directly to either polynomial.
If this optional normalization is used, its cost must be charged to the
original input: grouping $M$ sparse monomials by their jet exponents,
computing coefficient gcds, and performing divisions takes
$\poly(M,\ell_D(Q),\log q)$ bit operations. The resulting coefficient
lists contain at most $M(\ell_D(Q)+1)$ monomials.
Our decoder does not require this preprocessing.
 \section{Field descent and further refinements}
\label{sec:refinements}

The interpolation constraints depend only on the evaluation points and the
received values. This lets us replace the ambient field by the subfield
generated by the instance. We then return to
Kopparty's coefficient-lifting argument. Kopparty already branches at
resonant Taylor indices in his general solver~\cite[Corollary~4.5]{Kopparty15}.
We use that mechanism to give a refined root count for the present
interpolants when the ambient degree is not below the characteristic.

\subsection{Descent to the data field}

Let \(p\) be the characteristic of \(\F_q\).  For an instance
\((\boldsymbol\alpha,y)\), define its \emph{data field} by
\begin{equation}
  F_0:=\F_p(\alpha_1,\ldots,\alpha_n,y_1,\ldots,y_n)
  \subseteq\F_q,
  \qquad q_0:=|F_0|.
  \label{eq:data-field}
\end{equation}
\begin{theorem}[Data-field descent]
\label{thm:data-field-descent}
Let \(\alpha_1,\ldots,\alpha_n\in\F_q\) be distinct, let
\(y\in\F_q^n\), and let \(1\leq k\leq A\leq n\).  Set
\[
  \List_{<k}(y,A)
  :=
  \{P\in\F_q[X]:\deg P<k,\
      \agree((P(\alpha_i))_{i=1}^n,y)\geq A\}.
\]
Then the following statements hold.
\begin{enumerate}[label=\textup{(\roman*)}]
  \item Every polynomial in \(\List_{<k}(y,A)\) belongs to \(F_0[X]\),
        and \(q_0\geq n\).
  \item Every local interpolation system in
        \Cref{sec:hidden-derivative-interpolation,sec:exact-certificate}
        has coefficients in \(F_0\).  Whenever the dimension comparison
        gives a nonzero interpolant over \(\F_q\), it also gives one over
        \(F_0\).  Root finding may therefore be performed over extensions
        of \(F_0\), and its field-size dependence may be expressed in
        terms of \(q_0\) rather than \(q\).
\end{enumerate}
Given a standard explicit representation of \(\F_q\), a basis and
arithmetic for \(F_0\) can be constructed using a number of operations
polynomial in \(n\) and \(\log q\).
\end{theorem}

\begin{proof}
Fix \(P\in\List_{<k}(y,A)\).  Choose any \(k\) agreement positions.
Lagrange interpolation at these positions expresses \(P\) using only the
corresponding \(\alpha_i\)'s, the values \(y_i\), and inverses of their
nonzero pairwise differences.  Every such quantity lies in \(F_0\), which
proves the descent in (i).  The field \(F_0\) contains the \(n\) distinct
evaluation points, so \(q_0\geq n\).

For (ii), inspect the substitution in \Cref{def:local-constraints}.  All
matrix entries are polynomial expressions in the \(\alpha_i\)'s and
\(y_i\)'s with coefficients in the prime field, and hence lie in \(F_0\).
The rank of a matrix does not change after extending its coefficient field:
its rank is the largest size of a nonzero minor.  Thus a homogeneous system
over \(F_0\) that has a nonzero kernel vector over \(\F_q\) already has a
nonzero kernel vector over \(F_0\).  The multiplicity argument is an
identity over the same field.  By (i), every requested message is also
defined over \(F_0\), so the subsequent differential root search may take
\(F_0\) as its base field.

It remains to justify the algorithmic assertion.  Write
\(s=[\F_q:\F_p]\), and start with the \(\F_p\)-span of \(1\) and all input
elements.  Repeatedly enlarge the span whenever multiplying a current basis
element by an input generator leaves the span.  Each enlargement raises the
dimension, so at most \(s\) enlargements occur.  At termination the span is
stable under all input generators, hence equals the \(\F_p\)-algebra they
generate.  A finite subalgebra of a field is itself a field, so this algebra
is \(F_0\): multiplication by any nonzero element is injective on the
finite algebra and therefore surjective.  Linear algebra over \(\F_p\)
supplies a basis, multiplication
matrices, and inverses within the stated complexity.
\end{proof}

\subsection{Resonant coefficient lifting}

The unique-lifting argument in \Cref{lem:unique-regular-lift} solves for the
Taylor coefficient \(p_i\) by dividing by \(\binom is\), where \(s\) is the
active derivative order.  When the characteristic does not exceed the
ambient degree, the vanishing of this binomial coefficient does not destroy
the recursion. It makes \(p_i\) a free branch, as in Kopparty's general
coefficient-lifting procedure~\cite[Corollary~4.5]{Kopparty15}. Our new
contribution is to convert these branches into the global root bounds below
and specialize them to hidden-derivative interpolants. The following quantity
counts those branches, including the \(s\) coefficients below the initial
jet order:
\begin{equation}
  \kappa_s(D,p)
  :=
  \bigl|
    \{0\leq i\leq D:\tbinom is=0\text{ in characteristic }p\}
  \bigr|,
  \qquad
  \kappa(D,p,d):=\max_{0\leq s\leq d}\kappa_s(D,p).
  \label{eq:resonance-count}
\end{equation}

\begin{lemma}[Lifting with resonances]
\label{lem:resonant-lifting}
Let \(E\) be a finite field of characteristic \(p\) and cardinality \(S\).
Let \(0\leq s\leq D\), and let
\(R\in E[X,Y_0,\ldots,Y_s]\).  Suppose that
\[
  R(\alpha,\beta_0,\ldots,\beta_s)=0,
  \qquad
  \frac{\partial R}{\partial Y_s}
    (\alpha,\beta_0,\ldots,\beta_s)\neq0.
\]
Among the polynomials \(P\in E[X]_{\leq D}\) with
\(P^{[j]}(\alpha)=\beta_j\) for \(0\leq j\leq s\), at most
\begin{equation}
  S^{\kappa_s(D,p)-s}
  \label{eq:resonant-lift-count}
\end{equation}
satisfy
\(R(X,P,P^{[1]},\ldots,P^{[s]})\equiv0\).  All of them can be found by
successive coefficient lifting and branching only at the indices
\(i>s\) for which \(\binom is=0\) in \(E\).
\end{lemma}

\begin{proof}
Write \(P(\alpha+T)=\sum_{i=0}^Dp_iT^i\). The initial jet fixes
\(p_0,\ldots,p_s\). Suppose the coefficients through \(p_{i-1}\) have
been chosen. Changing \(p_i\) by \(c\) changes no coefficient below degree
\(T^{i-s}\), and the change in degree \(T^{i-s}\) is
\[
  c\binom is
  \frac{\partial R}{\partial Y_s}
    (\alpha,\beta_0,\ldots,\beta_s).
\]
If \(\binom is\neq0\), at most one value of \(c\) can make the coefficient
vanish.  If \(\binom is=0\), retaining every value of \(c\in E\) loses no
solution.  The binomial coefficient also vanishes for every \(i<s\), while
it equals one for \(i=s\).  The number of branching indices above \(s\) is
therefore \(\kappa_s(D,p)-s\).  Continuing through \(i=D\) and checking the
full identity proves the claim.
\end{proof}

We can now rerun the stratified root count without the hypothesis \(D<p\).
The remaining characteristic condition concerns only the individual
degrees of the differential polynomial.

\begin{theorem}[Differential root bound with resonances]
\label{thm:resonant-root-bound}
Let \(F_0\) be a finite field of cardinality \(q_0\) and characteristic
\(p\).  Let \(d\leq D\), and let
\[
  Q\in F_0[X,Y_0,\ldots,Y_d]\setminus\{0\}
\]
have individual \(Y_j\)-degrees smaller than \(p\).  Put
\begin{equation}
  \Delta:=\deg_YQ,
  \qquad
  \ell:=\ell_D(Q),
  \qquad
  h:=\max\{0,\ell-(D-d)\},
  \qquad
  \kappa:=\kappa(D,p,d).
  \label{eq:resonant-root-parameters}
\end{equation}
Let \(E/F_0\) be an extension of cardinality \(S=q_0^e\), and assume
\(S>h\).  Then
\begin{equation}
  \left|
    \{P\in F_0[X]_{\leq D}:
      Q(X,P,P^{[1]},\ldots,P^{[d]})\equiv0\}
  \right|
  \leq
  \frac{S}{S-h}\,\Delta S^\kappa.
  \label{eq:resonant-root-bound}
\end{equation}
All such polynomials can be enumerated deterministically in
\begin{equation}
  (\Delta+1)S^{\kappa+2}
  \poly\bigl(\operatorname{size}(Q),D,\ell,e,\log q_0\bigr)
  \label{eq:resonant-enumeration-time}
\end{equation}
bit operations.
\end{theorem}

\begin{proof}
We count all solutions over \(E\), which can only enlarge the desired set.
Consider any nonzero iterated partial derivative \(R\) of \(Q\).  If
\(R\) contains no jet variable, then it is a nonzero polynomial in \(X\)
and has no differential root.  Otherwise, let \(s\) be its largest active
jet order and set \(t=\deg_{Y_s}R\).  The hypothesis \(t<p\) implies that
\(R_t:=\partial^tR/\partial Y_s^t\) is nonzero and has active order below
\(s\).

As in the proof of \Cref{thm:root-bound-stratified}, the roots of \(R\)
split into \(t\) regular multiplicity strata and the residual roots of
\(R_t\).  For a root in a regular stratum, the corresponding separant is a
nonzero positive-order partial derivative of \(Q\).  Its specialization has
degree at most \(h\), so the root has at least \(S-h\) regular witnesses
\(\alpha\in E\).

Fix one such \(\alpha\) and the lower jet
\((\beta_0,\ldots,\beta_{s-1})\).  The possible values of \(\beta_s\),
counted jointly over the regular strata, are the roots of the nonzero
univariate fiber
\[
  R(\alpha,\beta_0,\ldots,\beta_{s-1},Z),
\]
classified by exact multiplicity.  There are at most \(t\) of them.
For each regular initial jet, \Cref{lem:resonant-lifting} leaves at most
\(S^{\kappa_s(D,p)-s}\) lifts.  Hence the number of regular
solution-witness pairs for \(R\) is at most
\[
  S\cdot S^s\cdot t\cdot
  S^{\kappa_s(D,p)-s}
  \leq tS^{\kappa+1}.
\]
Double counting the witnesses bounds the regular roots by
\(\frac{S}{S-h}tS^\kappa\).

The total \(Y\)-degree of the residual polynomial \(R_t\) is at most
\(\deg_YR-t\).  Induction on the active order, starting with a nonzero
polynomial in \(X\), therefore gives
\[
  |\mathcal Z_D(R)|
  \leq
  \frac{S}{S-h}
  \bigl(t+(\deg_YR-t)\bigr)S^\kappa.
\]
Taking \(R=Q\) proves \eqref{eq:resonant-root-bound}.

For enumeration, follow the same residual recursion.  At each regular
stratum, enumerate \(\alpha\), the lower jet, and \(\beta_s\), and branch
at precisely the resonant indices in \Cref{lem:resonant-lifting}.  The
number of resulting trials per stratum is at most \(S^{\kappa+2}\).  The
degrees encountered along the residual recursion sum to at most \(\Delta\).
Checking the original identity, retaining the polynomials over \(F_0\),
and removing duplicates gives \eqref{eq:resonant-enumeration-time}.
\end{proof}

The resonance count has a closed form in the range used by the decoder.

\begin{proposition}[Closed form for the resonance count]
\label{prop:resonance-closed-form}
Assume \(d<p\), and write \(D=up+v\) with \(0\leq v<p\).  Then, for
\(0\leq s\leq d\),
\begin{equation}
  \kappa_s(D,p)=us+\min\{s,v+1\},
  \qquad
  \kappa(D,p,d)=ud+\min\{d,v+1\}.
  \label{eq:resonance-closed-form}
\end{equation}
In particular, \(D<p\) gives \(\kappa=d\), recovering the exponent in
\Cref{thm:root-bound-stratified}.
\end{proposition}

\begin{proof}
For \(i=ap+b\) with \(0\leq b<p\), the characteristic-\(p\) identity
\[
  (1+Z)^i=(1+Z^p)^a(1+Z)^b
\]
shows that \(\binom is=\binom bs\) for \(s<p\).  The latter coefficient is
nonzero exactly when \(b\geq s\).  Each of the \(u\) complete residue
blocks therefore contributes \(s\) zeros, and the last block contributes
\(\min\{s,v+1\}\).  This proves the first formula.  Its right side is
nondecreasing in \(s\), so the maximum over \(s\leq d\) occurs at \(d\).
If \(D<p\), then \(u=0\) and \(v=D\geq d\), giving \(\kappa=d\).
\end{proof}

The maximum resonance count and the single worst separant degree in
\Cref{thm:resonant-root-bound} are convenient but can be unnecessarily
pessimistic.  The residual chain yields the following term-by-term form.

\begin{proposition}[Residual-chain root bound]
\label{prop:residual-chain-root-bound}
Under the hypotheses of \Cref{thm:resonant-root-bound}, define
\(R_0=Q\).  While \(R_j\) depends on a jet variable, let
\[
  s_j=\max\{s:R_j\text{ depends on }Y_s\},
  \qquad
  t_j=\deg_{Y_{s_j}}R_j,
  \qquad
  R_{j+1}=
  \frac{\partial^{t_j}R_j}{\partial Y_{s_j}^{t_j}},
\]
and put
\[
  h_j=
  \ell_D\!\left(\frac{\partial R_j}{\partial Y_{s_j}}\right).
\]
If \(S>\max_jh_j\), then
\begin{equation}
  |\mathcal R_D(Q)|
  \leq
  \sum_j
  \frac{S}{S-h_j}\,
  t_jS^{\kappa_{s_j}(D,p)}.
  \label{eq:residual-chain-root-bound}
\end{equation}
\end{proposition}

\begin{proof}
At stage \(j\), the roots outside the residual equation \(R_{j+1}=0\)
split among the \(t_j\) regular multiplicity strata.  Every corresponding
separant is a positive \(Y_{s_j}\)-derivative of \(R_j\), so its
specialization degree is at most \(h_j\).  The witness double count in the
proof of \Cref{thm:resonant-root-bound}, now retaining the active order,
therefore bounds these roots by
\[
  \frac{S}{S-h_j}\,t_jS^{\kappa_{s_j}(D,p)}.
\]
The remaining roots satisfy \(R_{j+1}=0\).  Iterating until the residual is
a nonzero polynomial in \(X\), which has no differential roots, and adding
the regular contributions proves \eqref{eq:residual-chain-root-bound}.
\end{proof}

\begin{proposition}[Exactness of the resonance exponent]
\label{prop:resonance-exponent-exact}
For \(0\leq s\leq D\), the equation \(Y_s=0\) over a field of cardinality
\(q_0\) and characteristic \(p\) has exactly
\begin{equation}
  q_0^{\kappa_s(D,p)}
  \label{eq:resonance-exponent-exact}
\end{equation}
polynomial solutions of degree at most \(D\).
\end{proposition}

\begin{proof}
Writing \(P=\sum_{i=0}^Da_iX^i\) gives
\[
  P^{[s]}=\sum_{i=s}^D a_i\binom isX^{i-s}.
\]
Distinct \(i\)'s give distinct monomials. Thus \(a_i\) is forced to vanish
when \(\binom is\neq0\) and is free when \(\binom is=0\), including every
\(i<s\).  The number of free coefficients is precisely
\(\kappa_s(D,p)\).
\end{proof}

The closed form lets us decode when the characteristic grows only linearly
with the block length, at a cost polynomial in the field size. We use the
fixed-parameter certificates in this statement so that the individual jet
degrees remain bounded independently of \(n\).

\begin{corollary}[Extension to linearly growing characteristic]
\label{cor:linear-characteristic-decoding}
Fix \(\delta\in(0,1)\) and \(c>0\).  There is \(N_{\delta,c}\) and a deterministic
algorithm such that, for every \(n\geq N_{\delta,c}\), every finite field
\(\F_q\) of characteristic \(p\geq cn\) and cardinality \(q\geq n\), every
integer \(1\leq k\leq n\), every
set of \(n\) distinct evaluation points, and every received word, the
algorithm outputs exactly all degree-\(<k\) polynomials having at least
\(k+\lceil\delta n\rceil\) agreements.  Its list size and running time are
\(q^{O_{\delta,c}(1)}\).

More precisely, for the interpolant selected in the proof, of derivative
order \(d\) and total jet degree \(\Delta_Q\), the differential-root stage has list
bound
\begin{equation}
  2\Delta_Q q^{\,2d\lceil1/c\rceil}
  \label{eq:linear-characteristic-root-bound}
\end{equation}
for all sufficiently large \(n\), with the convention that the exponent is
zero when \(d=0\).
\end{corollary}

\begin{proof}
Put \(A=k+\lceil\delta n\rceil\). If \(A>n\), return the empty list.
Otherwise the certified threshold used below is at most \(A\) and at most
\(n\).
If \(\delta\geq1/2\), use the order-zero fixed-parameter certificate in
\Cref{thm:small-derivative-orders}.  If \(1/4\leq\delta<1/2\), use its
order-one certificate at threshold \(k+\lceil n/4\rceil\).  If
\(0<\delta<1/4\), use the weighted-support certificate of the main paper
(\Cref{M-thm:weighted-support}).  In all three
cases, \(d\) and every individual jet degree are bounded in terms of
\(\delta\) alone. Interpolation and multiplicity are characteristic-free.
Once \(n\) is sufficiently large, \(p\) exceeds \(d\) and every individual
jet degree, so the resonant root theorem applies.

Since \(D<n\) and \(p\geq cn\), writing \(D=up+v\) gives
\[
  u<1/c,
  \qquad
  \kappa(D,p,d)\leq d(u+1)\leq d\lceil1/c\rceil
\]
by \Cref{prop:resonance-closed-form}.  The weighted specialization degree
is \(O_\delta(n)\), whereas \(q^2\geq n^2\); hence the extension
\(S=q^2\) satisfies \(S\geq2h\) for all sufficiently large \(n\).
Applying \Cref{thm:resonant-root-bound} proves
\eqref{eq:linear-characteristic-root-bound}; its constructive part and the
polynomial-size interpolation systems give the stated running time.  Final
degree and agreement filtering proves exactness.
\end{proof}

\subsection{Random regular witnesses}

Deterministic enumeration spends one factor of \(S\) on the top entry of a
regular jet.  Randomly sampling the witness point and finding the roots of
the resulting univariate fiber removes that factor.  We state the result as
a one-sided randomized algorithm: every output is checked, while
completeness holds with prescribed probability.

\begin{theorem}[Randomized regular-witness enumeration]
\label{thm:random-regular-witnesses}
Adopt the hypotheses and notation of \Cref{thm:resonant-root-bound}, and
assume \(S\geq2h\).  Given arithmetic in \(E\), for every
\(\varepsilon\in(0,1)\), there is a
randomized algorithm that outputs only members of the differential root set
in \eqref{eq:resonant-root-bound} and outputs every member with probability
at least \(1-\varepsilon\).  Its expected running time is
\begin{equation}
  S^\kappa
  \poly\bigl(
    \operatorname{size}(Q),D,\ell,\Delta,e,\log q_0,
    \log(1/\varepsilon)
  \bigr).
  \label{eq:random-witness-time}
\end{equation}
The same algorithm may check a degree restriction and an agreement
threshold before output, without changing the displayed field-size
exponent.
\end{theorem}

\begin{proof}
The root bound gives at most
\(M_0:=\max\{1,2\Delta S^\kappa\}\) solutions.  Every solution has at least
\(S-h\geq S/2\) regular witness points in the stratum where it leaves the
residual chain.  Sample
\[
  T:=\left\lceil\log_2(M_0/\varepsilon)\right\rceil
\]
independent uniform elements of \(E\).  A fixed solution has no sampled
regular witness with probability at most \(2^{-T}\).  A union bound over all
solutions makes the probability of missing any solution at most
\(\varepsilon\).

It remains to explain how one sampled point is processed without enumerating
the top jet entry.  At each residual polynomial \(R\) of active order
\(s\), enumerate the \(S^s\) lower jets and form the fiber
\[
  f(Z)=R(\alpha,\beta_0,\ldots,\beta_{s-1},Z).
\]
If \(f=0\), this lower jet gives no regular witness for \(R\) and may be
skipped.  Otherwise, find all roots of \(f\) in \(E\), classify them by
exact multiplicity, and apply the resonant lifting procedure. The roots can
be found by a Las Vegas computation polynomial in \(\deg f\), the
coefficient-list size of \(f\), and \(\log S\):
compute \(\gcd(f,Z^S-Z)\) by modular exponentiation, then recursively split
the resulting square-free product of linear factors.  For a factor of
degree \(r\), evaluation of a uniformly random polynomial of degree below
\(r\) at its \(r\) distinct roots gives independent uniform field values,
by Vandermonde interpolation.  In odd
characteristic, random quadratic-character tests separate any two factors
with constant probability.  In characteristic two, the absolute trace to
\(\F_2\) gives the analogous split.  Repeating only failed splits preserves
correctness and gives polynomial expected time.

For a fiber of degree \(t\), at most \(t\) top jet values survive.  The
lower-jet choices and the resonant branches number at most
\[
  S^sS^{\kappa_s(D,p)-s}\leq S^\kappa.
\]
The degrees of the fibers along the residual chain sum to at most
\(\Delta\). The number of samples is
\(T=O(\log\Delta+\kappa\log S+\log(1/\varepsilon))\), which is absorbed by
the polynomial factor in \eqref{eq:random-witness-time}. This proves the
expected running-time bound. Every candidate is verified against the original
identity, its base field, and any requested agreement condition, so the
algorithm has no false outputs.

In the decoder, the extension degree is at most two.  When \(e=2\), sample
monic quadratics over \(F_0\) and test irreducibility by a greatest common
divisor with \(Z^{q_0}-Z\).  Exactly \((q_0^2-q_0)/2\) of the \(q_0^2\)
monic quadratics are irreducible, as follows by pairing the elements of
\(\F_{q_0^2}\setminus F_0\) under Frobenius conjugation.  Thus this
constructs \(E\) in expected time
polynomial in \(\log q_0\) and does not change
\eqref{eq:random-witness-time}.
\end{proof}

For an interpolant produced from the received instance, the multiplicity
argument places the requested agreement list inside its differential root
set.  Combining
\Cref{thm:data-field-descent,thm:resonant-root-bound,thm:agreement-upper-bounds}
therefore gives three simultaneous list bounds.  One may take the minimum of the differential
bound \eqref{eq:resonant-root-bound}, the subset bound
\eqref{eq:subset-agreement-bound}, and, when its denominator is positive,
the Johnson-type bound \eqref{eq:johnson-agreement-bound}.  Over a prime
field with \(D<q\), the data field is the full field and
\(\kappa=d\).  Over a proper extension field, descent can reduce the base
of the exponential root-finding term from \(q\) to \(q_0\).
 \section{Small characteristic and the Frobenius obstruction}
\label{sec:characteristic}

The local interpolation and multiplicity arguments are identities over an
arbitrary field.  Small characteristic enters only when the differential
roots are extracted.  This section separates two failures of coefficient
lifting from an obstruction for differential polynomials whose monomials all contain
a higher derivative. The Frobenius decomposition used in the analysis is
standard and is proved here for completeness. The obstruction identifies
a support restriction that can prevent efficient enumeration of roots.

\subsection{Two failures of coefficient lifting}

There are two distinct reasons that the below-characteristic proof in
\Cref{sec:root-finding} cannot be applied verbatim.

\begin{proposition}[Characteristic failures]
\label{prop:characteristic-failures}
Over a field of characteristic \(p\), the following statements hold.
\begin{enumerate}[label=\textup{(\roman*)}]
  \item Formal differentiation of a nonconstant jet polynomial can vanish.
        For example, \(\partial(Y_1^p)/\partial Y_1=0\). Thus a nonzero
        recursive derivative chain is guaranteed by the hypothesis that
        every individual \(Y_j\)-degree is smaller than \(p\).
  \item Even with a nonzero separant, an initial jet need not determine a
        unique polynomial once \(D\geq p\).  The equation \(Y_1=0\) is
        satisfied by every \(P(X)=H(X^p)\), since \(P^{[1]}=0\).
\end{enumerate}
\end{proposition}

\begin{proof}
The first assertion follows from \(pY_1^{p-1}=0\).  For the second, the
identity
\((X+T)^p=X^p+T^p\) shows that the coefficient of \(T\) in
\(H((X+T)^p)\) is zero.  Hence the first Hasse derivative vanishes, even
though the separant of \(Y_1\) is one.  Equivalently, the coefficient
lifting step at Taylor degree \(p\) encounters \(\binom p1=0\).
\end{proof}

The first failure is excluded by the individual-degree hypothesis in
\Cref{thm:resonant-root-bound}.  The second is measured exactly by the
resonance count \(\kappa(D,p,d)\).  In a fixed small characteristic, that
count can grow linearly with \(D\), so the resulting enumeration may still
be exponential in the block length.

\subsection{Frobenius coordinates}

The behavior of Hasse derivatives in small characteristic is transparent
after grouping monomials by their exponent modulo a power of the
characteristic. We record the standard decomposition and its short proof to
fix the degree conventions used in the obstruction.

\begin{lemma}[Frobenius decomposition]
\label{lem:frobenius-decomposition}
Let \(F\) have characteristic \(p\), and let \(M=p^u\).  Every polynomial
\(P\in F[X]_{\leq D}\) has a unique representation
\begin{equation}
  P(X)=\sum_{a=0}^{M-1}X^aH_a(X^M),
  \qquad
  \deg H_a\leq\left\lfloor\frac{D-a}{M}\right\rfloor
  \label{eq:frobenius-decomposition}
\end{equation}
with the convention that \(H_a=0\) when \(a>D\).  For \(0\leq j<M\), its
Hasse derivatives are
\begin{equation}
  P^{[j]}(X)
  =
  \sum_{a=j}^{M-1}\binom ajX^{a-j}H_a(X^M).
  \label{eq:frobenius-hasse-derivatives}
\end{equation}
\end{lemma}

\begin{proof}
Collect the monomials of \(P\) according to their exponent modulo \(M\).
This gives \eqref{eq:frobenius-decomposition}, its degree bounds, and its
uniqueness.  Since \(M\) is a power of \(p\),
\[
  (X+T)^M=X^M+T^M.
\]
For \(j<M\), the factor \(H_a(X^M+T^M)\) contributes no positive power
\(T^j\).  Taking the coefficient of \(T^j\) in
\((X+T)^aH_a((X+T)^M)\) therefore gives
\(\binom ajX^{a-j}H_a(X^M)\).  Summing over \(a\) proves
\eqref{eq:frobenius-hasse-derivatives}.
\end{proof}

The simultaneous kernel of the derivatives of orders \(2,\ldots,d\) has
the same simple form in every characteristic.

\begin{lemma}[The higher-derivative Frobenius kernel]
\label{lem:higher-derivative-kernel}
Let \(F\) have characteristic \(p\), let \(d\geq2\), and let \(M=p^u\) be
the least power of \(p\) exceeding \(d\).  Then
\begin{equation}
  P^{[j]}=0\text{ for every }2\leq j\leq d
  \quad\Longleftrightarrow\quad
  P(X)=H_0(X^M)+XH_1(X^M)
  \label{eq:higher-derivative-kernel}
\end{equation}
for uniquely determined \(H_0,H_1\in F[X]\).
\end{lemma}

\begin{proof}
The reverse implication follows at once from
\eqref{eq:frobenius-hasse-derivatives}, since only the residue classes zero
and one occur and \(d<M\).  For the forward implication, the minimality of
\(M\) gives \(p^{u-1}\leq d\).  Hence
\(P^{[p^t]}=0\) for every \(1\leq t<u\).  If
\(a=\sum_{t=0}^{u-1}a_tp^t\), the characteristic-\(p\) factorization
\[
  (1+Z)^a=\prod_{t=0}^{u-1}(1+Z^{p^t})^{a_t}
\]
gives \(\binom a{p^t}=a_t\) in \(F\): factors with lower index have
combined degree at most \(p^t-1\), and factors with higher index have
degree larger than \(p^t\). In
\eqref{eq:frobenius-hasse-derivatives}, distinct residue classes \(a\)
contribute to distinct residue classes \(a-p^t\), so no cancellation is
possible.  Thus every residue class in the support of \(P\) has
\(a_t=0\) for \(1\leq t<u\).  If \(p=2\), its remaining digit is already
zero or one.  If \(p>2\), the same factorization and the equation
\(P^{[2]}=0\) give \(\binom{a_0}{2}=0\), so again
\(a_0\in\{0,1\}\). The Frobenius
decomposition now gives the stated form and its uniqueness.
\end{proof}

\subsection{An obstruction for higher-derivative-supported equations}

The preceding Frobenius identities force many roots of any equation
whose monomials all contain a higher derivative. We examine the consequence
for a decoder that enumerates all differential roots before filtering by agreement.

Call a differential polynomial \emph{higher-derivative supported} if
\begin{equation}
  Q\in(Y_2,\ldots,Y_d)
  \subseteq F[X,Y_0,\ldots,Y_d].
  \label{eq:higher-derivative-supported}
\end{equation}
Equivalently, every monomial of \(Q\) contains some \(Y_j\) with
\(j\geq2\). The support in \Cref{M-thm:weighted-support} of the main
paper does not impose this restriction, so the obstruction below does
not apply to every interpolant it supplies.

\begin{theorem}[Frobenius-plane obstruction]
\label{thm:frobenius-plane-obstruction}
Let \(F=\F_q\) have characteristic \(p\), let \(d\geq2\), and let \(M\) be
the least power of the characteristic exceeding \(d\).  Let
\[
  Q\in F[X,Y_0,\ldots,Y_d]
\]
be higher-derivative supported.  For every pair \(H_0,H_1\in F[X]\), the
polynomial
\[
  P(X)=H_0(X^M)+XH_1(X^M)
\]
satisfies
\begin{equation}
  Q(X,P,P^{[1]},\ldots,P^{[d]})\equiv0.
  \label{eq:frobenius-plane-identity}
\end{equation}
Consequently, for every \(D\geq0\), the polynomial \(Q\) has at least
\begin{equation}
  q^{\lfloor D/M\rfloor+\lfloor(D-1)/M\rfloor+2}
  \label{eq:frobenius-plane-root-count}
\end{equation}
differential roots of degree at most \(D\).
\end{theorem}

\begin{proof}
By \Cref{lem:higher-derivative-kernel}, every derivative \(P^{[j]}\) with
\(2\leq j\leq d\) is zero.  Each monomial of \(Q\) contains one of the
corresponding variables, which proves
\eqref{eq:frobenius-plane-identity}.  Under the degree restriction, one may
choose \(H_0\) and \(H_1\) independently with
\[
  \deg H_0\leq\left\lfloor\frac DM\right\rfloor,
  \qquad
  \deg H_1\leq\left\lfloor\frac{D-1}{M}\right\rfloor.
\]
The decomposition is unique, so these choices give exactly the number of
distinct polynomials displayed in \eqref{eq:frobenius-plane-root-count},
where \(\lfloor-1/M\rfloor=-1\) when \(D=0\).
\end{proof}

The next instance shows that this differential root set can be exponentially
larger than the agreement list.

\begin{corollary}[A singleton list with many differential roots]
\label{cor:fixed-characteristic-enumeration-obstruction}
Fix a prime \(p\), fix \(d\geq2\), and let \(M\) be as in
\Cref{thm:frobenius-plane-obstruction}.  Let \(q=n=p^s\), with \(n\)
sufficiently large.  Evaluate at every point of \(\F_q\), take
\(k=\lfloor n/2\rfloor\), \(D=k-1\), and take the all-zero received word.
For every \(\delta\in(0,1/2]\), the agreement threshold
\(A=k+\lceil\delta n\rceil\) is at most \(n\), and its list consists only
of the zero polynomial. On the
other hand, every nonzero higher-derivative-supported \(Q\) has at least
\begin{equation}
  q^{n/M}
  \label{eq:fixed-characteristic-spurious-root-count}
\end{equation}
differential roots of degree at most \(D\).
\end{corollary}

\begin{proof}
A nonzero polynomial of degree less than \(k\) has at most \(k-1\) zeros,
whereas \(A>k-1\).  Thus only the zero polynomial reaches the agreement
threshold.  If \(p=2\), then \(2M\mid n\) for all sufficiently large
\(s\), and the exponent in \eqref{eq:frobenius-plane-root-count} is
exactly \(n/M\).  If \(p\) is odd, write \(n=ML\), where \(L\) is odd.
Substituting \(D=(n-3)/2\) shows that the same exponent is \(L\) when
\(M=3\) and \(L+1\) when \(M\geq5\).  In every case it is at least
\(n/M\), proving \eqref{eq:fixed-characteristic-spurious-root-count}.
\end{proof}

\paragraph{Scope of the obstruction.}
\Cref{cor:fixed-characteristic-enumeration-obstruction} rules out a
specific algorithmic architecture over a full-field family in every fixed
characteristic: construct a higher-derivative-supported interpolant,
enumerate all of its bounded-degree differential roots, and only then filter
by agreement.  It does not give a large Reed--Solomon list.  The exhibited
instance has a singleton list.  It also does not rule out a support that
contains monomials outside the ideal \((Y_2,\ldots,Y_d)\), a solver that
uses the agreement constraints while extracting roots, or another capacity
decoder over fixed-characteristic fields.  The interpolation and multiplicity statements
in \Cref{sec:hidden-derivative-interpolation} remain valid in every
characteristic.
 \section{Why a positive capacity gap is necessary}
\label{sec:sharpness}

The decoder in \Cref{M-thm:main} of the main paper requires a strictly positive agreement gap
above the message dimension.  This requirement is not merely an artifact of
the interpolation argument.  We give a bad received word for every evaluation
set.  The same construction also quantifies what happens when the gap tends to
zero with the block length.

Large lists and negative results for Reed--Solomon list decoding have a
substantial literature; see, for example,
\cite{RR03,GR06,CW07,BKR10}. We prove the exact formulas below directly and
do not rely on those earlier constructions. We make no priority claim for the
elementary exact-capacity center. The coefficient-fiber formulation is
included because it applies to every evaluation set and exposes the
inverse-gap dependence needed here. The coefficient-pigeonhole construction
uses no property of the evaluation set beyond distinctness; later examples
use the full field or its multiplicative group.

\subsection{Field-independent upper bounds}

We first record two upper bounds that apply to every evaluation set and
field. The subset count is elementary; the second is the standard
Johnson-type incidence estimate. See~\cite[Chapter~7]{GRS25} for
background. We include the proofs to fix the finite conventions.

\begin{theorem}[Combinatorial agreement bounds]
\label{thm:agreement-upper-bounds}
Let \(1\leq k\leq A\leq n\). Then
\begin{equation}
  |\List_{k-1}(y,A)|
  \leq
  \left\lfloor\frac{\binom nk}{\binom Ak}\right\rfloor.
  \label{eq:subset-agreement-bound}
\end{equation}
If \(A^2>n(k-1)\), then also
\begin{equation}
  |\List_{k-1}(y,A)|
  \leq
  \left\lfloor
    \frac{n(A-k+1)}{A^2-n(k-1)}
  \right\rfloor.
  \label{eq:johnson-agreement-bound}
\end{equation}
\end{theorem}

\begin{proof}
A \(k\)-subset of coordinates can lie in the agreement sets of at most one
listed polynomial: two degree-\(<k\) polynomials agreeing with \(y\) there
would agree with each other at \(k\) points. Every list member contributes
at least \(\binom Ak\) such subsets, while only \(\binom nk\) exist. This
proves \eqref{eq:subset-agreement-bound}.

For the second bound, put \(L=|\List_{k-1}(y,A)|\), and choose exactly \(A\)
agreement positions \(S_j\) for each list member. Let
\(m_i=|\{j:i\in S_j\}|\). Distinct chosen sets meet in at most \(k-1\)
positions, so
\begin{equation}
  \sum_{i=1}^n m_i(m_i-1)
  \leq L(L-1)(k-1).
  \label{eq:pairwise-agreement-intersections}
\end{equation}
Since \(\sum_i m_i=LA\), Cauchy--Schwarz gives
\(\sum_i m_i^2\geq L^2A^2/n\). Substitution in
\eqref{eq:pairwise-agreement-intersections} yields
\[
  L\left(\frac{A^2}{n}-(k-1)\right)\leq A-k+1.
\]
The coefficient of \(L\) is positive by hypothesis, and rearranging proves
\eqref{eq:johnson-agreement-bound}.
\end{proof}

These bounds constrain the agreement list rather than the differential-root
set. They may be intersected with any algebraic root bound and remain valid
when the present differential solver is ineffective.

\subsection{A coefficient-pigeonhole construction}

Fix pairwise distinct evaluation points
\(\alpha_1,\ldots,\alpha_n\in\F_q\).  For a subset \(S\subseteq[n]\), write
\begin{equation}
  V_S(X):=\prod_{i\in S}(X-\alpha_i).
  \label{eq:subset-vanishing-polynomial}
\end{equation}
The high coefficients of these monic vanishing polynomials determine a
received word with a large agreement list.

\begin{proposition}[Coefficient-pigeonhole bad ball]
\label{prop:coefficient-pigeonhole}
Let \(1\leq k\leq A\leq n\).  There is a received word
\(y\in\F_q^n\) for which
\begin{equation}
  \bigl|\List_{k-1}(y,A)\bigr|
  \geq
  \left\lceil\frac{\binom{n}{A}}{q^{A-k}}\right\rceil.
  \label{eq:coefficient-pigeonhole-bound}
\end{equation}
This holds for every choice of the distinct evaluation points.
\end{proposition}

\begin{proof}
For every \(A\)-element subset \(S\subseteq[n]\), expand
\begin{equation}
  V_S(X)=X^A+\sum_{j=0}^{A-1}v_{S,j}X^j
  \label{eq:vanishing-polynomial-coefficients}
\end{equation}
and assign to \(S\) the coefficient signature
\begin{equation}
  \sigma(S):=(v_{S,k},v_{S,k+1},\ldots,v_{S,A-1})
  \in\F_q^{A-k}.
  \label{eq:high-coefficient-signature}
\end{equation}
There are at most \(q^{A-k}\) signatures.  Hence some signature \(\sigma\)
is shared by a family \(\mathcal C\) of at least the quantity on the
right-hand side of \Cref{eq:coefficient-pigeonhole-bound} distinct subsets.

Let
\begin{equation}
  H_\sigma(X):=X^A+\sum_{j=k}^{A-1}\sigma_jX^j,
  \qquad
  y_i:=H_\sigma(\alpha_i).
  \label{eq:bad-ball-center-polynomial}
\end{equation}
For every \(S\in\mathcal C\), define
\begin{equation}
  P_S(X):=H_\sigma(X)-V_S(X).
  \label{eq:bad-ball-message-polynomial}
\end{equation}
The coefficients of degrees \(k,k+1,\ldots,A\) cancel, so
\(\deg P_S<k\).  Moreover,
\[
  P_S(\alpha_i)=y_i
  \quad\Longleftrightarrow\quad
  V_S(\alpha_i)=0
  \quad\Longleftrightarrow\quad
  i\in S.
\]
Thus \(P_S\) agrees with \(y\) in exactly the \(A\) positions in \(S\), and
\(P_S\in\List_{k-1}(y,A)\).

It remains to count distinct polynomials rather than subsets.  If
\(P_S=P_T\), then \Cref{eq:bad-ball-message-polynomial} gives \(V_S=V_T\).
The two polynomials are monic and their roots among the distinct
\(\alpha_i\)'s are exactly the elements indexed by their respective subsets,
so \(S=T\).  The map \(S\mapsto P_S\) is injective on \(\mathcal C\), which
proves the proposition.
\end{proof}

The same argument gives an exact description whenever the center is the
evaluation of a monic degree-$A$ polynomial whose coefficients of degrees
at least $k$ have been fixed.

\begin{proposition}[Coefficient-fiber description]
\label{prop:coefficient-fiber}
Let $k\le A\le n$, and let
\[
  F(X)=X^A+\sum_{j=k}^{A-1}c_jX^j.
\]
For the received word $y_i=F(\alpha_i)$, the degree-$<k$ list at
agreement $A$ is in bijection with the $A$-element subsets
$S\subseteq[n]$ for which the monic root polynomial
$V_S=\prod_{i\in S}(X-\alpha_i)$ has coefficient $c_j$ in every degree
$k\le j<A$. Every listed polynomial has exactly $A$ agreements.
\end{proposition}

\begin{proof}
For such a subset, $P_S=F-V_S$ has degree less than $k$, and its agreement
set is exactly $S$. Conversely, if a degree-$<k$ polynomial $P$ has at
least $A$ agreements, then $F-P$ is monic of degree $A$ and has at
least $A$ distinct roots among the evaluation points. It has exactly those
roots, so it equals their monic root polynomial $V_S$. The coefficients in
degrees at least $k$ are therefore the prescribed $c_j$, and the two
maps are inverse.
\end{proof}

\begin{corollary}[Two exact coefficient fibers]
\label{cor:exact-coefficient-fibers}
Let $p$ be prime.
\begin{enumerate}[label=(\roman*)]
  \item On the full evaluation set $\F_p$, let $2\le A\le p-1$,
        $k=A-1$, and $y_x=x^A$. The list at agreement $A$ has exactly
        \[
          \frac1p\binom pA
        \]
        members.
  \item On $\F_p^\times$, let $A=2r\le p-1$, $k=2r-1$, and
        $y_x=x^{2r}$. The list at agreement $2r$ has exactly
        \[
          \frac{\binom{p-1}{2r}+p-1}{p}
        \]
        members.
\end{enumerate}
\end{corollary}

\begin{proof}
By \Cref{prop:coefficient-fiber}, both lists count subsets whose elements
sum to zero. For part (i), translation by $c\in\F_p$ changes the sum of an
$A$-subset by $Ac$. Since $A\ne0$ in $\F_p$, the $p$ sum fibers
are equinumerous.

For part (ii), let $N_a$ be the number of zero-sum $a$-subsets of
$\F_p^\times$. Splitting a zero-sum $a$-subset of $\F_p$ according to
whether it contains zero gives
\[
  N_a+N_{a-1}=\frac1p\binom pa,
  \qquad N_0=1.
\]
Pascal's identity solves this recurrence as
\[
  N_a=\frac{\binom{p-1}{a}+(p-1)(-1)^a}{p}.
\]
Substitute the even integer $a=2r$.
\end{proof}

\subsection{The exact capacity boundary}

At agreement \(A=k\), no pigeonhole loss occurs.  In fact, the construction
has exactly one nearby message for every \(k\)-element subset of coordinates.

\begin{proposition}[Exact maximum list at capacity]
\label{prop:exact-capacity-list}
Let \(1\leq k\leq n\), and set
\begin{equation}
  y^{\mathrm{cap}}
  :=(\alpha_1^k,\ldots,\alpha_n^k)\in\F_q^n.
  \label{eq:capacity-center}
\end{equation}
Then
\begin{equation}
  \max_{y\in\F_q^n}
  \bigl|\List_{k-1}(y,k)\bigr|
  =
  \bigl|\List_{k-1}(y^{\mathrm{cap}},k)\bigr|
  =\binom{n}{k}.
  \label{eq:exact-capacity-list-size}
\end{equation}
\end{proposition}

\begin{proof}
For every \(k\)-element subset \(S\subseteq[n]\), the polynomial
\begin{equation}
  P_S(X):=X^k-V_S(X)
  \label{eq:capacity-message}
\end{equation}
has degree less than \(k\) and agrees with \(y^{\mathrm{cap}}\) on every
position in \(S\).  Distinct subsets give distinct polynomials by the
injectivity argument in the proof of \Cref{prop:coefficient-pigeonhole}.

Conversely, let \(P\) have degree less than \(k\) and agree with
\(y^{\mathrm{cap}}\) in at least \(k\) positions.  The monic polynomial
\(X^k-P(X)\) has degree \(k\), and the agreement positions are roots of this
polynomial.  It has at most \(k\) distinct roots, so it has exactly \(k\) of
the evaluation points as roots.  If \(S\) indexes these roots, monicity gives
\(X^k-P=V_S\), and hence \(P=P_S\).  The displayed family therefore exhausts
the list.

The universal upper bound is \eqref{eq:subset-agreement-bound} with \(A=k\),
which proves the maximum identity.
\end{proof}

For a fixed rate \(R\in(0,1)\) and \(k/n\to R\),
\Cref{eq:exact-capacity-list-size} is
\(\exp((h(R)+o(1))n)\), where \(h\) is the binary entropy function in natural
units.  Thus no polynomial list-size bound is possible at radius \(1-R\),
even when the evaluation set is fixed in advance.

\subsection{A shrinking-gap obstruction}

The coefficient-pigeonhole loss is exact enough to locate an obstruction
inside a vanishing neighborhood of capacity.  Define
\begin{equation}
  h(x):=-x\log x-(1-x)\log(1-x),
  \qquad 0\log 0:=0.
  \label{eq:binary-entropy}
\end{equation}

\begin{proposition}[Entropy lower bound for a shrinking gap]
\label{prop:shrinking-gap-entropy}
For integers \(1\leq k\leq k+s\leq n\),
\begin{equation}
  \log\max_{y\in\F_q^n}
  \bigl|\List_{k-1}(y,k+s)\bigr|
  \geq
  n h\left(\frac{k+s}{n}\right)
  -s\log q-\log(n+1).
  \label{eq:shrinking-gap-entropy-bound}
\end{equation}
Consequently, consider any sequence of evaluation sets over fields
\(\F_{q_n}\), with message dimensions \(k_n\) and nonnegative gaps \(s_n\),
such that \(k_n+s_n\le n\) and
\begin{equation}
  \frac{k_n}{n}\longrightarrow R\in(0,1),
  \qquad
  \gamma_n:=\frac{s_n}{n}\longrightarrow0,
  \qquad
  \limsup_{n\to\infty}\gamma_n\log q_n\leq c.
  \label{eq:shrinking-gap-asymptotic-hypotheses}
\end{equation}
Then the worst-case list size \(L_n\) at agreement \(k_n+s_n\) satisfies
\begin{equation}
  \liminf_{n\to\infty}\frac{1}{n}\log L_n
  \geq \max\{0,h(R)-c\}.
  \label{eq:shrinking-gap-asymptotic-bound}
\end{equation}
In particular, if \(\gamma_n\log q_n=o(1)\), then
\(L_n\geq\exp((h(R)-o(1))n)\).
\end{proposition}

\begin{proof}
Apply \Cref{prop:coefficient-pigeonhole} with \(A=k+s\).  For
\(a=(k+s)/n\), the standard type bound gives
\begin{equation}
  \binom{n}{k+s}\geq\frac{\exp(nh(a))}{n+1}.
  \label{eq:binomial-type-lower-bound}
\end{equation}
For completeness, when \(0<a<1\), the \((k+s)\)-th summand is a mode of the
binomial expansion
\(1=(a+(1-a))^n\).  Since there are at most \(n+1\) summands, this mode is at
least \(1/(n+1)\).  Dividing it by
\(a^{k+s}(1-a)^{n-k-s}=\exp(-nh(a))\) proves
\Cref{eq:binomial-type-lower-bound}; the endpoint cases follow directly.
Combining this estimate with
\Cref{eq:coefficient-pigeonhole-bound} and taking logarithms proves
\Cref{eq:shrinking-gap-entropy-bound}.

Under \Cref{eq:shrinking-gap-asymptotic-hypotheses}, continuity of \(h\)
gives \(h(k_n/n+\gamma_n)\to h(R)\).  Divide
\Cref{eq:shrinking-gap-entropy-bound} by \(n\), take the lower limit, and use
\(L_n\geq1\) to obtain \Cref{eq:shrinking-gap-asymptotic-bound}.
\end{proof}

The last proposition gives a necessary scale, not only an endpoint example.
If the worst-case lists are subexponential and \(k_n/n\to R\in(0,1)\), then
any vanishing agreement gap must satisfy
\begin{equation}
  \frac{s_n}{n}\log q_n\geq h(R)-o(1),
  \qquad\text{or equivalently}\qquad
  s_n\geq\bigl(h(R)-o(1)\bigr)\frac{n}{\log q_n}.
  \label{eq:necessary-shrinking-gap}
\end{equation}
For fields of size \(q_n=\Theta(n)\), a gap of
\(o(1/\log n)\) above rate still leaves
\(\exp((h(R)-o(1))n)\) codewords in one Hamming ball.  This includes gaps
of order \(1/\log^2 n\) and \(n^{-c}\) for every fixed \(c>0\).  The obstruction does
not rule out polynomial lists once the gap reaches the scale
\(h(R)/\log q_n\); it identifies the strongest conclusion supplied by this
construction.

\subsection{Dependence of the list size on the gap}

The same construction gives a much larger inverse-gap obstruction when the
block length is chosen on the scale \(2^{1/\delta}\).

\begin{theorem}[A doubly exponential inverse-gap bad ball]
\label{thm:doubly-exponential-gap-list}
For every \(0<\delta\le1/16\), there is a prime \(q=n\), a dimension \(k\)
with
\[
  \left|\frac kn-\frac12\right|\le\frac{\delta}{16},
\]
and a received word \(y\in\F_q^n\) such that, for
\(A=k+\lceil\delta n\rceil\),
\begin{equation}
  |\List_{k-1}(y,A)|
  \ge
  \exp\left(\frac{\delta}{16}\,2^{1/\delta}\right).
  \label{eq:doubly-exponential-gap-list}
\end{equation}
The center \(y\) is existential, obtained from the coefficient-pigeonhole
construction.
\end{theorem}

\begin{proof}
Put \(T=2^{1/\delta}\) and \(N=\lceil T/16\rceil\). Bertrand's postulate
gives a prime \(q\) with \(N\le q<2N\), so
\begin{equation}
  \frac{T}{16}\le n=q\le\frac{T}{4}.
  \label{eq:doubly-exponential-length}
\end{equation}
For \(\delta\le1/16\), these inequalities also give \(q\ge8/\delta\).
Take \(k=\lfloor q/2\rfloor\), put \(s=\lceil\delta q\rceil\), and set
\(A=k+s\). Since \(q\) is odd,
\[
  x:=\frac Aq-\frac12=\frac{s-1/2}{q}
\]
satisfies \(0<x\le17\delta/16<1/4\). The binary entropy function from
\eqref{eq:binary-entropy} has
\(|h''(1/2+x)|\le16/3\) for \(|x|\le1/4\). Taylor's theorem therefore gives
\[
  h(A/q)\ge\log2-\frac83x^2\ge\log2-4\delta^2.
\]
Also \(s\le\delta q+1\), and
\(\log q\le(\log2)/\delta-\log4\). Applying
\Cref{prop:shrinking-gap-entropy} yields
\begin{align}
  \log\max_y|\List_{k-1}(y,A)|
  &\ge qh(A/q)-s\log q-\log(q+1) \notag\\
  &\ge\delta q(\log4-4\delta)-2\log q-1.
  \label{eq:doubly-exponential-entropy-step}
\end{align}
The inequality \(\delta^2T\ge256>256\log2+128\delta\), valid throughout
\(0<\delta\le1/16\), implies
\(2\log q+1\le\delta q/8\). Hence
\[
  \log\max_y|\List_{k-1}(y,A)|
  \ge\delta q(\log4-1/4-1/8)
  >\delta q
  \ge\frac{\delta T}{16}.
\]
Exponentiating proves \eqref{eq:doubly-exponential-gap-list}. Finally,
\[
  \left|\frac kq-\frac12\right|=\frac1{2q}\le\frac{\delta}{16}.
\]
\end{proof}

\begin{remark}[Quantifiers in the lower bound]
\label{rem:doubly-exponential-quantifiers}
The lengths in \Cref{thm:doubly-exponential-gap-list} satisfy
\(n=\Theta(2^{1/\delta})\). The theorem does not give this list size at
every sufficiently large length for a fixed gap, and these lengths may lie
below the threshold \(N_\delta\) in \Cref{M-thm:main} of the main paper. An unrestricted
\(\delta\)-dependent prefactor in an eventual list bound can absorb the
example. Thus the theorem rules out all-length bounds polynomial in
\(n\) and \(1/\delta\), but it does not by itself lower-bound the derivative
order in \Cref{M-thm:main} of the main paper.
\end{remark}

The next construction, developed in this work, makes the inverse-gap
obstruction explicit. It uses the multiplicative cosets of a prime field
rather than pigeonholing coefficient signatures.

\begin{theorem}[Explicit multiplicative-coset lists]
\label{thm:multiplicative-coset-lists}
For every $0<\delta\le1/8$, let
\[
  t=\left\lceil\frac1\delta\right\rceil-1.
\]
There are infinitely many prime-field Reed--Solomon codes with
$q=n+1$, evaluation set $\F_q^\times$, and rate in
$[1/5,1/2]$, together with a received word whose list at agreement
threshold $k+\lceil\delta n\rceil$ has at least
\begin{equation}
  \frac{2^t}{t+1}
  \label{eq:coset-list-lower-bound}
\end{equation}
members.
\end{theorem}

\begin{proof}
Put $r=\lfloor t/2\rfloor$. The definition of \(t\) gives \(\delta t<1\).
By Dirichlet's theorem on primes in arithmetic
progressions~\cite{Dirichlet1837}, infinitely many primes satisfy
$q\equiv1\pmod t$. Choose one
large enough that
\[
  h=\frac{q-1}{t}
  \ge\max\{2,(1-\delta t)^{-1}\},
\]
and put $n=q-1=th$. The cyclic group
$\F_q^\times$ has a subgroup of order $h$; its $t$ cosets are the root
sets of
\[
  X^h-\beta_1,\ldots,X^h-\beta_t
\]
for the $t$ distinct $h$-th powers $\beta_i$.

For an $r$-element subset $S\subseteq[t]$, define
\[
  G_S(X)=\prod_{i\in S}(X^h-\beta_i),
  \qquad
  P_S(X)=X^{rh}-G_S(X).
\]
The leading terms cancel, so $\deg P_S\le(r-1)h$. Set
$k=(r-1)h+1$ and take the received word $y_\alpha=\alpha^{rh}$ on
$\F_q^\times$. Then $P_S(\alpha)=y_\alpha$ precisely on the $r$
selected cosets, giving exactly $rh$ agreements. Moreover,
\[
  \delta n=\delta th\le h-1,
\]
so $\lceil\delta n\rceil\le h-1$, and hence
$k+\lceil\delta n\rceil\le rh$. Distinct subsets give distinct monic
polynomials $G_S$, and therefore distinct messages $P_S$. Finally,
\[
  \binom tr\ge\frac{2^t}{t+1}
\]
because the largest of the $t+1$ binomial coefficients is at least their
average. For $t\ge4$ and sufficiently large $h$, the rate
$k/n=(r-1)/t+1/(th)$ lies in the asserted interval.
\end{proof}

Since \(t=\delta^{-1}+O(1)\), this gives the explicit bound
\[
  \frac{2^t}{t+1}
  =
  \exp\left(\frac{\log2}{\delta}-O(\log(1/\delta))\right)
\]
at infinitely many lengths, complementing the larger existential bound in
\Cref{thm:doubly-exponential-gap-list}.

\begin{corollary}[Exponential dependence on the inverse gap]
\label{cor:exponential-inverse-gap-list}
There are absolute constants \(c,\delta_0>0\) such that, for every
\(0<\delta<\delta_0\), some prime-field Reed--Solomon code of rate
\(R\) satisfying \(|R-1/2|\le\delta\) has a received word with at least
\[
  \exp(c/\delta)
\]
degree-\(<k\) polynomials agreeing in at least
\(k+\lceil\delta n\rceil\) positions.  In particular, no uniform
worst-case list bound \(\poly(1/\delta)\) is possible for arbitrary
prime-field evaluation sets.
\end{corollary}

\begin{proof}
Let \(N=\lceil1/\delta\rceil\).  Bertrand's postulate gives a prime
\(p\) with \(N\le p<2N\).  Set \(q=n=p\), use all elements of
\(\F_p\) as evaluation points, and take \(k=\lfloor p/2\rfloor\).
Put \(s=\lceil\delta p\rceil\) and \(A=k+s\).  Since
\(\delta p<2\delta(1/\delta+1)\), one has \(s\le3\) for all sufficiently
small \(\delta\), and \(A\le n\).  By
\Cref{prop:coefficient-pigeonhole},
\[
  \max_y|\List_{k-1}(y,A)|
  \ge
  \frac{\binom p{k+s}}{p^s}.
\]
The ratio \((k+s)/p\) lies in \([2/5,3/5]\) for large \(p\).  The type
bound \eqref{eq:binomial-type-lower-bound} therefore gives
\[
  \log\max_y|\List_{k-1}(y,A)|
  \ge
  p\min_{x\in[2/5,3/5]}h(x)-3\log p-\log(p+1).
\]
For sufficiently large \(p\), this is at least \(c_0p\) for an absolute
\(c_0>0\).  Since \(p\ge1/\delta\), the claim follows after decreasing
the constant.
\end{proof}
\section{Refreshed application certificate values}
\label{sec:refreshed-application-certificates}

The application checkers distributed with this companion now select the
tight first-order Taylor-degree model explicitly. The certificate engine's
default remains the legacy model, so historical commands and stored baselines
retain their original meaning.

\paragraph{ProveKit over cubic Goldilocks.}
For the second code in the retuned lookup configuration, with
$(n,k,A)=(2048,128,453)$ and support $(38,23,133)$, the squarefree list
bound is
\[
 \Lambda\leq\frac{2659384185}{326}.
\]
The machine-readable certificate records the safe generic ceiling
$8157621$. Since $\Lambda$ is an integer, the separation calculation can
instead use the sharper floor $8157620$, which gives
$137.0819759979$ bits with one out-of-domain evaluation. The
successive-stage alternative has exact bound $15314557817/163$ and ceiling
$93954343$, giving $130.0304965539$ bits in the same calculation. These are
analytical bounds for the fixed retuned configuration, not measured proof
sizes.

\paragraph{ZisK compressed final proof.}
At $51$ queries and $22$ grinding bits, the two batching stages contribute
the exact error
\[
 \frac{4073087021134687044}
 {(2^{64}-2^{32}+1)^3},
\]
or $130.1791716533$ bits. The corresponding query term has
$128.0003360286$ bits under the checker's independent-uniform-output model.
The native saved proofs remain $286013$ and $272783$ bytes. By contrast, the
analytical field-vector accounting is $254032$ and $242272$ bytes; the latter
figures are not serialized proof-file measurements.

\paragraph{LambdaVM CPU certificate.}
For the list that selects the $38$ main-column polynomials after the two
early evaluations, the tight certificate gives
\[
 \Lambda\leq\frac{367473974837}{3230},
 \qquad \lceil\Lambda\rceil=113769033.
\]
With the existing cubic-Goldilocks field, $20$ grinding bits, $208$ CPU
queries, and the unchanged non-CPU profiles, the complete changed-component
bound has $128.0627324732$ bits. The analytical field-and-hash accounting
saves $55992$ bytes. The measured \texttt{rkyv} reduction below is $59832$
bytes because it includes the serializer's actual query records, metadata,
alignment, and the two new evaluation vectors.

\section{Measured LambdaVM proof savings and overhead}
\label{sec:lambda-measurements}

The first-order bounds also support a concrete change to LambdaVM's CPU
proof: reduce its query count from $219$ to $208$ and add two early
evaluations of its $38$ main columns. We implement this experimental profile
on a fork. On the same $32768$-instruction assembly benchmark and empty
input, the complete CPU subproof shrinks from $1\,233\,024$ to
$1\,173\,192$ serialized bytes: a saving of $59\,832$ bytes
($59.832$ decimal KB, $4.85\%$)~\cite{LambdaVMMeasured2026}.
This comparison includes every CPU commitment, claim, and opening, including
the added evaluations, and excludes shared VM-level metadata.
All $25$ non-CPU subproofs keep $219$ queries. The complete VM proof
decreases by the same $59\,832$ bytes, from $36\,868\,888$ to $36\,809\,056$.
Both proofs are generated, deserialized, and verified through the release
CLI with their matching, caller-selected profiles. This is one measured
workload, not an average over programs.

\paragraph{Separating the two changes.}
To identify the cost of the added evaluations, we cross their presence
with the CPU query count. The resulting four cells have the following
complete VM-proof sizes under \texttt{rkyv} serialization.
\begin{center}
\begin{tabular}{@{}lrrr@{}}
\toprule
Cell & CPU queries & Early evaluations & Bytes\\
\midrule
Control & $219$ & $0$ & $36\,868\,888$\\
Evaluations only & $219$ & $2$ & $36\,870\,744$\\
Queries only & $208$ & $0$ & $36\,807\,200$\\
Actual profile & $208$ & $2$ & $36\,809\,056$\\
\bottomrule
\end{tabular}
\end{center}
The two early vectors add $1856$ bytes; removing eleven query records saves
$61\,688$ bytes. These effects are additive, giving the $59\,832$-byte net
saving. They include the serializer's metadata and alignment.
The middle two cells are measurement instruments, not proposed protocol
profiles; in particular, the $208$-query cell without early evaluations
is not certified by the paper's argument.

\paragraph{End-to-end timing and added work.}
The normal-grinding experiment uses an Apple M4 Max with $16$ threads,
release builds, and grinding factor $20$. We rotate the four cells through
a balanced execution order and compare the actual profile with the control
within each round. The reported $95\%$ intervals use a round-level
percentile bootstrap of the paired mean~\cite{LambdaVMOverhead2026}.
\begin{center}
\begin{tabular}{@{}lrrr@{}}
\toprule
Measurement & Rounds per cell & Change (ms) & $95\%$ interval (ms)\\
\midrule
Proving & $20$ & $-15.2$ & $[-49.4,+18.9]$\\
Whole verification & $50$ & $+0.02$ & $[-0.62,+0.52]$\\
CPU verifier components & $50$ & $-0.081$ & $[-0.117,-0.048]$\\
\bottomrule
\end{tabular}
\end{center}
Neither end-to-end comparison resolves a latency change on this benchmark.
The component measurements explain why verification can nonetheless do
less work: at $208$ queries, evaluations add $0.167$~ms, whereas removing
eleven queries with evaluations enabled saves $0.291$~ms. These are
conditional effects from the four-cell experiment, not the difference
between the actual profile and control; that net difference is the
$-0.081$~ms component entry above.

The prover's extra arithmetic is visible in diagnostic runs with grinding
disabled: the actual profile retires $264.4$ million more instructions,
about $0.1845\%$ of the control. Three-round single-thread isolation
attributes roughly $305$~ms to the evaluation and DEEP work. This serial
cost is not a $16$-thread latency estimate, since the table proofs run
concurrently. The largest new transient buffer is $3$~MiB; process RSS,
roughly $1.9$~GiB, does not resolve it. Separate encode, decode, and archive
access measurements differ by less than $0.05$~ms. These codec runs were
not interleaved, so their differences are descriptive.

\paragraph{Scope and reproduction.}
The experimental profile enforces one $32768$-row CPU chunk and retains
cubic Goldilocks, blowup two, and the baseline grinding and batching.
The mathematical certificate controls the changed CPU component; neither
the certificate nor these measurements establish a new full-VM security
theorem. The implementation and complete-proof reproduction record are
pinned in~\cite{LambdaVMMeasured2026}. The supplied timing data and
measurement report are distributed with this companion in
\texttt{experiments/lambdavm-anchors}. Its checker verifies recorded byte
arithmetic and recomputes paired mean timing contrasts; it does not rerun
the prover or regenerate the reported bootstrap intervals.

 \section*{Acknowledgements}
We thank Kai Zhe Zheng for close discussion and feedback, including a suggestion to name the central interpolation polynomial in BCPZZ the hidden-derivative explainer.
 \section*{Disclosure}
The second and third authors are Research Partners at a16z crypto, which is
an investor in various blockchain-based platforms, as well as in the crypto
ecosystem more broadly. General a16z disclosures are available at
\url{https://www.a16z.com/disclosures/}.
 \bibliographystyle{alpha}
\newcommand{\etalchar}[1]{$^{#1}$}
\begin{thebibliography}{BCH{\etalchar{+}}26}

\bibitem[Dao26a]{LambdaVMMeasured2026}
Quang Dao.
\newblock LambdaVM CPU RS anchors: complete-proof size experiment.
\newblock Implementation 7310d1e8, measurement record c40323c8, September 2026.
\newblock \url{https://github.com/quangvdao/lambda_vm/blob/c40323c85432d939567a96b6b00ec1e7f7067db0/docs/experiments/cpu_rs_anchors_v1.md}.

\bibitem[Dao26b]{LambdaVMOverhead2026}
Quang Dao.
\newblock LambdaVM CPU Reed--Solomon anchor overhead: measurements and reproducibility.
\newblock September 2026.
\newblock \url{https://github.com/quangvdao/rs-beyond-johnson/tree/main/experiments/lambdavm-anchors}.

\bibitem[BCH{\etalchar{+}}26]{BCHKS25}
Eli {Ben-Sasson}, Dan Carmon, Ulrich Hab{\"o}ck, Swastik Kopparty, and
  Shubhangi Saraf.
\newblock On proximity gaps of {Reed--Solomon} codes.
\newblock In {\em Proceedings of the 58th Annual ACM Symposium on Theory of
  Computing (STOC 2026)}, pages 1157--1167. ACM, 2026.
\newblock \url{https://doi.org/10.1145/3798129.3800827}. Full version:
  \url{https://eprint.iacr.org/2025/2055}.

\bibitem[BCP{\etalchar{+}}26]{BCPZZ26}
Joshua Brakensiek, Yeyuan Chen, Aaron Putterman, Zihan Zhang, and Kai~Zhe
  Zheng.
\newblock Algorithmic list decoding of {Reed--Solomon} codes up to capacity.
\newblock Technical Report TR26-164, revision 1, Electronic Colloquium on
  Computational Complexity, 2026.
\newblock September 5, 2026.
  \url{https://eccc.weizmann.ac.il/report/2026/164/}.

\bibitem[BKR10]{BKR10}
Eli {Ben-Sasson}, Swastik Kopparty, and Jaikumar Radhakrishnan.
\newblock Subspace polynomials and limits to list decoding of {Reed--Solomon}
  codes.
\newblock {\em IEEE Transactions on Information Theory}, 56(1):113--120, 2010.

\bibitem[CW07]{CW07}
Qi~Cheng and Daqing Wan.
\newblock On the list and bounded distance decodability of {Reed--Solomon}
  codes.
\newblock {\em SIAM Journal on Computing}, 37(1):195--209, 2007.

\bibitem[Dir37]{Dirichlet1837}
Peter Gustav~Lejeune Dirichlet.
\newblock {Beweis des Satzes, dass jede unbegrenzte arithmetische Progression,
  deren erstes Glied und Differenz ganze Zahlen ohne gemeinschaftlichen Factor
  sind, unendlich viele Primzahlen enth{\"a}lt}.
\newblock {\em Abhandlungen der K{\"o}niglich Preussischen Akademie der
  Wissenschaften}, pages 45--81, 1837.
\newblock Volume for 1837, printed in 1839.

\bibitem[DKSS13]{DKSS13}
Zeev Dvir, Swastik Kopparty, Shubhangi Saraf, and Madhu Sudan.
\newblock Extensions to the method of multiplicities, with applications to
  {Kakeya} sets and mergers.
\newblock {\em SIAM Journal on Computing}, 42(6):2305--2328, 2013.

\bibitem[GR06]{GR06}
Venkatesan Guruswami and Atri Rudra.
\newblock Limits to list decoding {Reed--Solomon} codes.
\newblock {\em IEEE Transactions on Information Theory}, 52(8):3642--3649,
  2006.

\bibitem[GRS25]{GRS25}
Venkatesan Guruswami, Atri Rudra, and Madhu Sudan.
\newblock Essential coding theory, 2025.
\newblock Draft dated August 26, 2025.

\bibitem[GS99]{GS99}
Venkatesan Guruswami and Madhu Sudan.
\newblock Improved decoding of {Reed--Solomon} and algebraic-geometry codes.
\newblock {\em IEEE Transactions on Information Theory}, 45(6):1757--1767,
  1999.

\bibitem[Kop15]{Kopparty15}
Swastik Kopparty.
\newblock List-decoding multiplicity codes.
\newblock {\em Theory of Computing}, 11(5):149--182, 2015.

\bibitem[RR03]{RR03}
Gitit Ruckenstein and Ron~M. Roth.
\newblock Bounds on the list-decoding radius of {Reed--Solomon} codes.
\newblock {\em SIAM Journal on Discrete Mathematics}, 17(2):171--195, 2003.

\bibitem[Sch80]{Schwartz80}
Jacob~T. Schwartz.
\newblock Fast probabilistic algorithms for verification of polynomial
  identities.
\newblock {\em Journal of the ACM}, 27(4):701--717, 1980.

\end{thebibliography}
 
\appendix
\section{Exact counting identities}
\label{app:exact-counts}

We collect two elementary evaluations used in the finite and coarse interpolation estimates.

\begin{lemma}[Closed form for the staircase count]
\label{lem:staircase-closed-form}
For integers $D\ge1$ and $t$, the function in \eqref{eq:staircase-count} satisfies
\[
  \mathsf N_D(t)=0
  \quad\text{if }t\le0,
\]
and, if $t>0$, then
\begin{equation}
  \mathsf N_D(t)
  =(s+1)t-\frac{Ds(s+1)}2,
  \qquad
  s=\left\lfloor\frac{t-1}{D}\right\rfloor.
  \label{eq:staircase-closed-form}
\end{equation}
\end{lemma}

\begin{proof}
When $t\le0$, no nonnegative pair $(a,b)$ satisfies $a+Db<t$.  Suppose $t>0$.  For a fixed $b$, the
admissible values are $a=0,\ldots,t-Db-1$, so there are $t-Db$ of them.  Such a value exists exactly for
$0\le b\le s$.  Summing the arithmetic progression gives \eqref{eq:staircase-closed-form}.
\end{proof}

We next evaluate the aggregate residual-rank count used in the enlarged-map analysis. This is the specialization of
\eqref{eq:exact-local-rank} with $M=m=d^3$ after replacing every $\Lambda_d(W+s)$ by its common upper
bound.  Define
\[
  h_s=\left\lceil\frac{m-s}{d}\right\rceil,
  \qquad
  \beta_s=(s+1)(m+1)-(s-h_s+1)_+(m-h_s+1)_+.
\]

\begin{lemma}[Exact aggregate residual-rank count]
\label{lem:exact-beta-sum}
Let $d\ge2$ and $m=d^3$.  Then
\begin{equation}
  \sum_{s=0}^{m-1}\beta_s
  =
  \frac{
    d^2(d^2-d+1)(2d^2+d+1)(4d^2-3d+5)
  }{12}.
  \label{eq:exact-beta-sum}
\end{equation}
Consequently,
\begin{equation}
  \frac1{m^3}\sum_{s=0}^{m-1}\beta_s
  =
  \frac{(d^2-d+1)(2d^2+d+1)(4d^2-3d+5)}{12d^7}
  <\frac{2}{3d}.
  \label{eq:exact-beta-normalized}
\end{equation}
\end{lemma}

\begin{proof}
Parameterize $s$ uniquely as
\[
  s=m-dh+t,
  \qquad
  1\le h\le d^2,
  \qquad
  0\le t<d.
\]
Then $h_s=h$.  Moreover,
\[
  s-h+1=d^3-(d+1)h+t+1.
\]
The identity
\[
  d^3+1=(d+1)(d^2-d+1)
\]
shows that this quantity is nonnegative exactly when $h\le d^2-d+1$; at the last such value it equals
$t$.  Since $m-h+1>0$, the positive part can therefore be removed after restricting the second sum to that
range.  Hence
\begin{align}
  \sum_{s=0}^{m-1}\beta_s
  &=(m+1)\sum_{h=1}^{d^2}\sum_{t=0}^{d-1}(m-dh+t+1)\notag\\
  &\quad-
  \sum_{h=1}^{d^2-d+1}\sum_{t=0}^{d-1}
  (m-(d+1)h+t+1)(m-h+1).
  \label{eq:beta-two-sums}
\end{align}
Using
\[
  \sum_{t=0}^{d-1}t=\frac{d(d-1)}2,
  \qquad
  \sum_{h=1}^{H}h=\frac{H(H+1)}2,
  \qquad
  \sum_{h=1}^{H}h^2=\frac{H(H+1)(2H+1)}6
\]
in \eqref{eq:beta-two-sums}, with $H=d^2-d+1$, gives \eqref{eq:exact-beta-sum} after collecting factors.

It remains to prove the strict inequality.  After multiplying by $12d^7$, it is equivalent to
\[
  (d^2-d+1)(2d^2+d+1)(4d^2-3d+5)<8d^6.
\]
The difference between the right and left sides is
\[
  d^4(10d-21)+d^2(11d-14)+3d-5.
\]
It equals \(17\) at \(d=2\), and every displayed group is positive for
\(d\ge3\).
\end{proof}
 \section{Source-level repairs to the low-rate proof}
\label{app:source-audit}

This appendix records the points at which our reconstruction differs from the
formal statements and notation of Brakensiek, Chen, Putterman, Zhang, and
Zheng~\cite{BCPZZ26}. The hidden-derivative identity and the construction of
the local kernel subspaces remain unchanged. We additionally prove that these
subspaces exhaust the kernel of the enlarged map. The repairs to the published
argument concern one quantifier in the global dimension comparison, one
rounding step, and one rank-label slip.

\subsection{The dimension quantifier in Proposition~3.13}

In the source notation, Lemma~3.2 gives the dimension estimate
\begin{equation}
  \dim\Qspace>
  \frac{\theta^3}{384}|\mathcal B|(k-1)m^3,
  \label{eq:source-dimension-lower-bound}
\end{equation}
while Lemma~3.12 bounds the rank imposed at one received point by
\begin{equation}
  |\mathcal B|m^3d^{-2\theta/(5+\theta)}.
  \label{eq:source-local-rank-bound}
\end{equation}
The source proceeds by verifying the following sufficient inequality:
\begin{equation}
  \frac{\theta^3}{384}\frac{k-1}{n}
  d^{2\theta/(5+\theta)}>1.
  \label{eq:source-required-dimension-inequality}
\end{equation}
This is the interface used between the source's dimension and rank bounds.

Proposition~3.13 assumes only
\(k/n\leq(1-\theta)\varepsilon\), but its proof then uses
\(k=\lfloor(1-\theta)\varepsilon n\rfloor\) to lower-bound
\((k-1)/n\)~\cite[Section 3.3, Proposition 3.13]{BCPZZ26}.  An upper bound on
\(k/n\) does not imply the lower bound required by
\Cref{eq:source-required-dimension-inequality}.  The proof therefore
establishes the proposition at the design dimension used inside the proof,
not simultaneously at every smaller message dimension.

There are two direct repairs.  The local statement can be restricted to
\(k=\lfloor(1-\theta)\varepsilon n\rfloor\), together with the finite-size
condition needed for the asserted lower bound on \((k-1)/n\).  To recover the
claimed ``rate at most'' quantifier, run interpolation and root finding at an
\emph{ambient} dimension \(K\) satisfying
\begin{equation}
  k\leq K<A,
  \qquad
  \frac{K-1}{n}\geq\rho_*>0,
  \label{eq:source-ambient-repair-conditions}
\end{equation}
and replace \(k-1\) by \(K-1\) throughout the interpolation support and the
root-finding call.  Every degree-less-than-\(k\) message is an ambient
candidate, and candidates of degree at least \(k\) are removed after root
finding.  The source proof can take
\(K=\lfloor(1-\theta)\varepsilon n\rfloor\).  For the one-sided support in this companion, the adaptive choice in
\Cref{eq:adaptive-ambient} supplies both inequalities in
\Cref{eq:source-ambient-repair-conditions} uniformly over the actual
message dimension. The main paper uses a different ambient baseline in
its weighted-support certificate (\Cref{M-thm:weighted-support}).

\subsection{The ceiling in Lemma~3.2}

The source support in equations~(17) and~(18) contains the condition
\begin{equation}
  |c|\leq
  \left\lceil\left(1+\frac{3\theta}{4}\right)m\right\rceil.
  \label{eq:source-ceiling-cap}
\end{equation}
The proof of Lemma~3.2 replaces this ceiling by a floor and concludes
\(|c|\leq(1+3\theta/4)m\), which is not valid as written
~\cite[Section 3.1, equations (17)--(19)]{BCPZZ26}.  The extra unit can be
absorbed without changing the support or the final parameters.

Indeed, the rate assumption and \(A=\lceil\varepsilon n\rceil\) give, for
\(k>1\),
\begin{equation}
  \frac{A}{k-1}>\frac{1}{1-\theta}.
  \label{eq:source-strong-agreement-ratio}
\end{equation}
Combining \Cref{eq:source-ceiling-cap,eq:source-strong-agreement-ratio} yields
the exact slack estimate
\begin{align}
  \frac{mA}{k-1}-|c|
  &>
  \frac{m}{1-\theta}
  -\left(1+\frac{3\theta}{4}\right)m-1 \notag\\
  &=
  \frac{\theta m}{4}
  +\frac{\theta^2m}{1-\theta}-1.
  \label{eq:source-ceiling-repair}
\end{align}
Consequently, the source's equation~(19) follows whenever
\begin{equation}
  \theta^2m\geq1-\theta.
  \label{eq:source-rounding-sufficient-condition}
\end{equation}
This condition is already forced by the source parameter regime.  To see
this, its equation~(20) implies \(\varepsilon<\theta^{15}\), because its base
is smaller than \(\theta^3\) and its exponent is larger than five.  Hence
\(d=\lceil\varepsilon^{-3/\theta}\rceil>
\theta^{-45/\theta}\) and
\(m=d^3>\theta^{-135/\theta}\), which gives
\(\theta^2m>1>1-\theta\).  Replacing the two incorrect rounding lines by
\Cref{eq:source-ceiling-repair} repairs Lemma~3.2 with its original ceiling.

\subsection{Rank notation}

The source factors the local constraint map as
\(\Phi_{\alpha,y}=\Gamma\circ\Psi_{\alpha,y}\), and therefore
\(\rank(\Phi_{\alpha,y})\leq\rank(\Gamma)\).  In the proof of Lemma~3.12,
rank-nullity and the kernel spaces from Lemma~3.11 are then applied to
\(\Gamma\), but the first displayed rank in that calculation is written as
\(\rank(\Phi)\)~\cite[Section 3.2.2, equations (21)--(24)]{BCPZZ26}.  The
correct chain is
\begin{equation}
  \rank(\Phi_{\alpha,y})
  \leq\rank(\Gamma)
  \leq
  \dim\mathcal V-
  \dim\left(\bigoplus_{r=0}^{m-1}\mathcal K_r\right).
  \label{eq:source-phi-gamma-repair}
\end{equation}
The remainder of the proof bounds the last expression, so this is a label
correction rather than a change to the kernel argument.  The corresponding
factorization and rank calculation appear in
\Cref{eq:local-map-factorization,lem:exact-local-rank}. The source proves the
kernel inclusion needed for this bound; the reverse inclusion and resulting
exact formula for \(\rank(\Gamma)\) are new to this work.

\subsection{Hypotheses at the interfaces}

For later use, we make explicit four hypotheses that are implicit or only
stated elsewhere in the source.

\begin{enumerate}[label=\textup{(\roman*)}]
  \item Proposition~3.13 must fix pairwise distinct evaluation points.  Its
        final multiplicity argument counts \(A\) distinct roots, and this is
        exactly where distinctness is used.
  \item The source interpolation support contains \(mA/(k-1)\), so its
        standalone statement requires \(k\geq2\).  Lemma~3.12 uses
        \(d\geq3\) in its last numerical estimate.  The later root-finding
        invocation additionally requires \(k>d\).
  \item The local conditions in source equation~(15) are imposed only for
        indices with \(db<m\).  Statements that quantify over every
        \(b\geq0\) must either retain this restriction or adopt the convention
        that no divisibility condition is imposed once \(db\geq m\).
  \item The interpolation argument through Proposition~3.13 is
        characteristic-free.  Primality of \(q\), the inequalities
        \(q\geq k>d\), the individual degree conditions
        \(\deg_{Y_i}Q<q\), and the weighted-degree condition belong to the
        source root-finding theorem
        ~\cite[Theorem 2.1]{BCPZZ26}, not to the local interpolation lemma.
\end{enumerate}

With these interfaces stated, the source equations~(13)--(16) correspond to
\Cref{eq:backward-taylor-truncated,eq:local-substitution,eq:local-constraints};
the global multiplicity step of Proposition~3.13 corresponds to
\Cref{prop:global-interpolation-principle}.  Our changes begin at the choice of
ambient degree and interpolation support, not at the hidden-derivative local
identity.
 \section{A fixed-gap boundary asymptotic}
\label{app:boundary-asymptotics}

The main paper uses a weighted support to obtain its explicit derivative
order. For comparison, this appendix records a sharper asymptotic analysis
of the one-sided support from \Cref{sec:exact-certificate}. It identifies a
natural Gumbel boundary law and an exact one-dimensional variational
constant. The result is modular and is not used in the proof of
\Cref{M-thm:main} of the main paper. Its limit keeps the rate and agreement
fixed while $d\to\infty$; it does not justify a joint limit as the gap tends
to zero.

Recall the staircase function \(\mathsf G\) from
\Cref{eq:normalized-staircase-function}.
For an integer $s$ and $D\ge1$, the exact identity
\begin{equation}
  \mathsf N_D(s)=D\,\mathsf G(s/D)
  \label{eq:boundary-staircase-exact}
\end{equation}
follows by summing, for each $j\ge0$, the $s-Dj$ possible
$X$-exponents. Fix constants $0<\rho <\eta\le1$ and put
\begin{equation}
  h=\frac{\eta}{\rho },
  \qquad
  \theta=1-\frac1h.
  \label{eq:boundary-h-theta}
\end{equation}
For fixed $d,m,M,W$, if $D/n\to \rho $ and $A/n\to\eta$, then
\begin{equation}
  \frac{\dim\Qspace}{n}
  \longrightarrow
  \rho \mathscr D_d(h),
  \qquad
  \mathscr D_d(h):=
  \sum_{\omega(c)\le W}\sum_{b=0}^{M}
  \mathsf G(mh-b-|c|).
  \label{eq:boundary-fixed-limit}
\end{equation}
This is an immediate application of
\Cref{eq:boundary-staircase-exact} to the exact dimension formula and
continuity of \(\mathsf G\). Thus the limiting one-sided certificate is
\begin{equation}
  \rho \mathscr D_d(h)>\mathsf R_d(m,M,W).
  \label{eq:boundary-limit-certificate}
\end{equation}

For $\tau>0$ and an integer $j\ge0$, define
\begin{equation}
  I_j(\tau):=
  \int_{\tau}^{\infty}\log^j(u/\tau)e^{-u}\,du,
  \qquad
  C(h):=\min_{\tau>0}\frac{6\tau^{-1/h}}{I_3(\tau)}.
  \label{eq:boundary-I-C}
\end{equation}

\begin{theorem}[One-sided boundary ratio]
\label{thm:boundary-ratio}
Fix $0<\rho <\eta\le1$, $h=\eta/\rho $, and $\tau>0$. Let
$H=H_{d-1}$, let $\gamma$ denote Euler's constant, and set
\begin{equation}
  \xi=-h(\gamma+\log\tau),
  \qquad
  \alpha=h-\frac{\xi}{H}.
  \label{eq:boundary-alpha}
\end{equation}
Choose
\begin{equation}
  m=d^3,\qquad
  M=\left\lfloor\frac{m}{\sqrt H}\right\rfloor,\qquad
  W=\left\lfloor\frac{\alpha dm}{H}\right\rfloor.
  \label{eq:boundary-parameters}
\end{equation}
Then, as $d\to\infty$,
\begin{equation}
  \frac{\rho \mathscr D_d(h)}{\mathsf R_d(m,M,W)}
  =
  (1+o(1))
  \frac{\eta I_3(\tau)\tau^{1/h}}6
  \frac{d^\theta}{H_{d-1}}.
  \label{eq:boundary-ratio}
\end{equation}
The error is uniform when $(\rho ,\eta,\tau)$ ranges over a compact subset of
$\{0<\rho <\eta\le1,\ \tau>0\}$.
\end{theorem}

\begin{proof}
We give the probabilistic and endpoint estimates that determine every
factor in \Cref{eq:boundary-ratio}. Let
\[
  \mathcal S_W=
  \left\{x\in\R_{\ge0}^{d-1}:
  \sum_{i=1}^{d-1}ix_i\le W\right\},
  \qquad
  V_W=\frac{W^{d-1}}{((d-1)!)^2}.
\]
If $X$ is uniform on $\mathcal S_W$ and
$E_0,\ldots,E_{d-1}$ are independent unit exponential random variables,
then the standard exponential representation of the uniform simplex gives
\begin{equation}
  X_i=\frac{WE_i}{i\sum_{j=0}^{d-1}E_j}.
  \label{eq:boundary-exponential-representation}
\end{equation}
Consequently, for
\[
  T_d=
  \frac d{\sum_{j=0}^{d-1}E_j}
  \sum_{i=1}^{d-1}\frac{E_i}{i},
\]
we have $|X|/m=(\alpha/H+o(1))T_d$.
More precisely, the definition of $W$ gives
\[
  \frac{W}{dm}=\frac\alpha H+O\left(\frac1{dm}\right),
\]
so the coefficient error is negligible on the boundary scale used below.

The sum $\sum_{i=1}^{d-1}E_i/i$ has the same distribution as the maximum
of $d-1$ independent unit exponentials. Hence
\begin{equation}
  T_d-H\Longrightarrow Z:=G_{\rm Gum}-\gamma,
  \label{eq:boundary-gumbel-limit}
\end{equation}
where $\Pr(G_{\rm Gum}\le z)=\exp(-e^{-z})$. The convergence holds with
uniformly integrable cubes. Indeed, the maximum $M_N$ satisfies
\[
  \Pr(M_N-\log N>t)\le e^{-t},
  \qquad
  \Pr(M_N-\log N<-t)\le e^{-e^t},
\]
and the factor $d/\sum E_j$ differs from one in $L^6$ by $O(d^{-1/2})$.
It follows from \Cref{eq:boundary-alpha} that
\begin{equation}
  H\left(h-\frac{|X|}{m}\right)
  \Longrightarrow \xi-hZ
  \label{eq:boundary-slack-limit}
\end{equation}
with uniformly integrable positive cubes.

Now put $a_m=m/H$ and
\[
  F_{m,M}(t):=\sum_{b=0}^{M}\mathsf G(t-b).
\]
Since $M/a_m=\sqrt H+o(1)$, comparison with a Riemann integral gives,
uniformly on compact $x$-sets,
\begin{equation}
  a_m^{-3}F_{m,M}(a_mx)
  \longrightarrow\frac{x_+^3}{6}.
  \label{eq:boundary-cubic-riemann}
\end{equation}
Indeed, the elementary bounds
\[
  \frac{x_+^2}{2}\leq \mathsf G(x)\leq\frac{x_+^2}{2}+x_+
\]
reduce the normalized sum to a Riemann sum for
$\frac12\int_0^\infty(x-t)_+^2\,dt$. They also give the uniform envelope
$a_m^{-3}F_{m,M}(a_mx)\leq C(1+x_+^3)$.

We next justify the lattice transfer, including the floors. Put
$n_d=d-1$, $B_d=\binom d2$, and
\[
  U_W:=\bigcup_{\omega(c)\leq W}\bigl(c+[0,1)^{n_d}\bigr).
\]
Then $\mathcal S_W\subseteq U_W\subseteq\mathcal S_{W+B_d}$. Since
$F_{m,M}$ is nondecreasing and
$|c|\leq|x|<|c|+n_d$ on the cube based at $c$, we have the exact
sandwich
\begin{align}
  \int_{\mathcal S_W}F_{m,M}(mh-|x|)\,dx
  &\leq \mathscr D_d(h) \notag\\
  &\leq
  \int_{\mathcal S_{W+B_d}}
    F_{m,M}(mh-|x|+n_d)\,dx.
  \label{eq:boundary-dimension-sandwich}
\end{align}
Here
\[
  \frac{B_d}{W}=O\left(\frac{H}{d^2}\right),
  \qquad
  \frac{n_dH}{m}=O\left(\frac{H}{d^2}\right),
  \qquad
  \frac{\operatorname{Vol}(\mathcal S_{W+B_d})}{V_W}
  =1+O\left(\frac Hd\right)=1+o(1).
\]
Increasing the simplex radius by $B_d$ perturbs the normalized boundary
slack by $O_{L^3}(dH^2/m)=o(1)$, using the exponential representation and
the cubic moment bound above. The additive cube shift perturbs it by
$O(dH/m)=o(1)$. The envelope following
\Cref{eq:boundary-cubic-riemann}, together with the uniform integrability
in \Cref{eq:boundary-slack-limit}, therefore shows that the two normalized
integrals in \Cref{eq:boundary-dimension-sandwich} have the same limit.

If $U$ is unit exponential, then $-\log U$ is standard Gumbel, and
\[
  \xi-hZ=h\log(U/\tau).
\]
Thus \Cref{eq:boundary-slack-limit,eq:boundary-cubic-riemann} yield
\begin{equation}
  \frac{\mathscr D_d(h)}{V_W}
  =
  (1+o(1))
  \frac{m^3h^3}{6H^3}I_3(\tau).
  \label{eq:boundary-dimension-asymptotic}
\end{equation}

It remains to evaluate the exact enlarged-map rank. Write its index as
$r=m-s$ and put $u_s=\lceil s/d\rceil$. The unit-cube bounds give,
uniformly for $1\leq s\leq m$,
\begin{equation}
  \left(\frac{W+m-s}{W}\right)^{d-1}
  \leq\frac{\Lambda_d(W+m-s)}{V_W}
  \leq
  \left(\frac{W+m-s+B_d}{W}\right)^{d-1}.
  \label{eq:boundary-rank-lattice-sandwich}
\end{equation}
Since $d(m/W)^2=O(H^2/d)$ and $dB_d/W=O(H/d)$, expanding the
logarithms in \Cref{eq:boundary-rank-lattice-sandwich} yields the uniform
estimate
\begin{equation}
  \frac{\Lambda_d(W+m-s)}{V_W}
  =\exp\left(\frac H\alpha-\frac{Hs}{\alpha m}+o(1)\right).
  \label{eq:boundary-rank-lattice-asymptotic}
\end{equation}

The exact residual factor from \Cref{eq:exact-local-rank} is
\[
  \beta_s:=(m-s+1)(M+1)
  -\positive{m-s-u_s+1}\positive{M-u_s+1}.
\]
For fixed $0<\varepsilon<L$, uniformly on
$\varepsilon m/H\leq s\leq Lm/H$, both positive parts are active and
direct expansion gives
\begin{equation}
  \beta_s=(1+o(1))\frac{ms}{d}.
  \label{eq:boundary-beta-window}
\end{equation}
For every $1\leq s\leq m$, the exact definition also gives
\begin{equation}
  \beta_s\leq\left(\frac sd+1\right)(m+M+2).
  \label{eq:boundary-beta-global}
\end{equation}
To see this when both positive parts are active, use
$\beta_s=u_s(m-s+M+2)-u_s^2$; if one is inactive, the same upper bound
follows from $u_s>m-s+1$ or $u_s>M+1$.

Set $\lambda=H/(\alpha m)$. Insert
\Cref{eq:boundary-rank-lattice-asymptotic,eq:boundary-beta-window} in the
rank sum first on
$\varepsilon m/H\leq s\leq Lm/H$. The global bound
\Cref{eq:boundary-beta-global} controls the complementary ranges. After
normalization by $m^3/(dH^2)$, the range $s<\varepsilon m/H$ contributes
$O(\varepsilon^2)$, while the range $s>Lm/H$ contributes
$O((L+1)e^{-L/(2\alpha)})$. The terms arising from the $+1$ and $M$ in
\Cref{eq:boundary-beta-global} are $o(1)$ because
$dH/m=o(1)$ and $M/m=o(1)$. Letting first $d\to\infty$, then
$\varepsilon\downarrow0$ and $L\to\infty$, reduces the rank to the
geometric sum
\[
  \sum_{s=1}^{m}s e^{-\lambda s}
  =(1+o(1))\lambda^{-2}
  =(1+o(1))\frac{\alpha^2m^2}{H^2}.
\]
Here the first equality follows from
$\sum_{s\geq1}sz^s=z/(1-z)^2$ with $z=e^{-\lambda}$; truncation after
$m$ has relative error
$O((1+\lambda m)e^{-\lambda m})
=O(Hd^{-1/\alpha+o(1)})=o(1)$.
Consequently,
\begin{equation}
  \frac{\mathsf R_d(m,M,W)}{V_W}
  =
  (1+o(1))
  \frac{m^3\alpha^2e^{H/\alpha}}{dH^2}.
  \label{eq:boundary-rank-asymptotic}
\end{equation}

Divide \Cref{eq:boundary-dimension-asymptotic} by
\Cref{eq:boundary-rank-asymptotic} and multiply by $\rho $. Since
$H=\log d+\gamma+o(1)$, $\alpha=h-\xi/H$, and $\rho h=\eta$,
\[
  \frac{d}{e^{H/\alpha}}
  =(1+o(1))d^{1-1/h}\tau^{1/h}.
\]
Substitution proves \Cref{eq:boundary-ratio}. Every comparison used above
is uniform on the stated compact sets.
\end{proof}

\begin{lemma}[The variational constant]
\label{lem:boundary-constant}
For every $h\ge1$, the minimum defining $C(h)$ exists and has a unique
minimizer \(\tau_h\in(0,1/h)\), characterized by
\begin{equation}
  I_3(\tau_h)=3hI_2(\tau_h).
  \label{eq:boundary-stationarity}
\end{equation}
Moreover, $C(h)$ is strictly decreasing in $h$.
\end{lemma}

\begin{proof}
The objective tends to infinity as $\tau\downarrow0$, since
$I_3(\tau)=O((1+|\log\tau|)^3)$, and also as $\tau\to\infty$, since
$I_3(\tau)\le6e^{-\tau}/\tau^3$. It therefore has an interior minimizer.
For $j\geq1$, differentiation under the integral gives
$I_j'(\tau)=-(j/\tau)I_{j-1}(\tau)$, so every stationary point satisfies
\Cref{eq:boundary-stationarity}.

To prove uniqueness, put
\[
  r(\tau):=\frac{I_3(\tau)}{3I_2(\tau)}.
\]
After the change of variables \(u=\tau v\), this is one third of the
expectation of \(\log V\) under the probability density on \(v>1\)
proportional to \(\log^2(v)e^{-\tau v}\). Differentiating this expectation
with respect to \(\tau\) gives
\[
  r'(\tau)=-\frac13\operatorname{Cov}_\tau(\log V,V)<0,
\]
because both functions are strictly increasing and the distribution is
nondegenerate. As \(\tau\to\infty\), integration by parts gives
\(I_3/I_2<3/\tau\to0\). As \(\tau\downarrow0\), restricting the integral
for \(I_3\) to \(u\in[1,2]\) gives
\(I_3=\Omega(\log^3(1/\tau))\), while direct expansion of
\(\log(u/\tau)\) gives \(I_2=O(\log^2(1/\tau))\). Hence
\(r(\tau)\to\infty\) at zero and \(r(\tau)\to0\) at infinity. Thus
\(r(\tau)=h\) has exactly one solution,
and the logarithmic derivative of the objective changes from negative to
positive there. This solution is the unique minimizer. Integration by
parts gives \(I_3(\tau)<3I_2(\tau)/\tau\), hence
\(\tau_h<1/h\).

Finally, if \(h_2>h_1\), evaluate the \(h_2\)-objective at
\(\tau_{h_1}<1\). The factor \(\tau^{-1/h}\) strictly decreases with \(h\)
at that point, so \(C(h_2)<C(h_1)\).
\end{proof}

\begin{corollary}[Fixed-gap one-sided certificate]
\label{cor:boundary-uniform-certificate}
Fix $\delta\in(0,1)$ and $\zeta>0$, and put
\[
  \rho_0=\frac{\delta(1-\delta)}2,
  \qquad
  h_0=(1-\delta)^{-1}.
\]
Choose a minimizer $\tau_0$ for $C(h_0)$. There is
$d_0=d_0(\delta,\zeta)$ with the following property. Let
$d\geq d_0$ be an integer satisfying
\begin{equation}
  \frac{d^\delta}{H_{d-1}}
  >(1+\zeta)C((1-\delta)^{-1}),
  \label{eq:boundary-uniform-condition}
\end{equation}
and set $H=H_{d-1}$, $m=d^3$, and $M=\lfloor m/\sqrt H\rfloor$. There is
$N=N(\delta,\zeta,d)$ such that the following holds for every $n\geq N$
and $1\leq k\leq n$. Set $A:=k+\lceil\delta n\rceil$. If $A\leq n$,
define
\begin{gather*}
  K:=\max\{k,\lceil\rho_0n\rceil\},
  \qquad D:=K-1,
  \qquad h_{n,k}:=\frac AD,\\
  \xi_{n,k}:=-h_{n,k}(\gamma+\log\tau_0),
  \qquad
  \alpha_{n,k}:=h_{n,k}-\frac{\xi_{n,k}}H,\\
  W_{n,k}:=
  \left\lfloor\frac{\alpha_{n,k}dm}{H}\right\rfloor.
\end{gather*}
Then $d<D$, and the exact finite inequality
\Cref{eq:finite-interpolation-certificate} holds with parameters
$(D,A,d,m,M,W_{n,k})$. Thus the derivative order is rate-independent,
while the support threshold $W_{n,k}$ is explicitly rate-adaptive.
\end{corollary}

\begin{proof}
For orientation, first associate to an actual rate $x\in[0,1-\delta]$ the
limiting quantities
\[
  \rho _x=\max\{x,\rho_0\},\qquad
  \eta_x=x+\delta,\qquad
  h_x=\eta_x/\rho _x,\qquad
  \theta_x=1-1/h_x.
\]
Then $h_x\geq h_0$. If
\[
  C_{\tau_0}(h):=\frac{6\tau_0^{-1/h}}{I_3(\tau_0)},
\]
the proof of \Cref{lem:boundary-constant} gives
$C_{\tau_0}(h_x)\leq C(h_0)$.

Suppose first that $x\geq\rho_0$. Then
$\theta_x=\delta/\eta_x$. Once $\delta\log d\geq1$,
\[
  \frac{d}{d\eta}
  \log\bigl(\eta d^{\delta/\eta-\delta}\bigr)
  =\frac{\eta-\delta\log d}{\eta^2}\leq0
  \qquad(\delta\leq\eta\leq1).
\]
The expression inside the logarithm is one at $\eta=1$, so
\begin{equation}
  \eta_xd^{\theta_x-\delta}\geq1.
  \label{eq:boundary-unpadded-uniform}
\end{equation}
If $x<\rho_0$, then $\eta_x\geq\delta$ and
$\theta_x\geq(1+\delta)/2$. For all sufficiently large $d$,
\begin{equation}
  \eta_xd^{\theta_x-\delta}
  \geq\delta d^{(1-\delta)/2}\geq1.
  \label{eq:boundary-padded-uniform}
\end{equation}
Combining \Cref{eq:boundary-uniform-condition} with
\Cref{eq:boundary-unpadded-uniform,eq:boundary-padded-uniform} gives,
uniformly in $x$,
\[
  \frac{\eta_xd^{\theta_x}}H
  >(1+\zeta)C_{\tau_0}(h_x).
\]
By \Cref{thm:boundary-ratio}, the left-to-right ratio in the limiting
certificate is therefore larger than one with a uniform positive margin,
provided $d_0$ is sufficiently large. The parameters
$(\rho _x,\eta_x,h_x)$ range over a compact set with $\rho _x\geq\rho_0$, so the
uniformity assertion in that theorem applies.

We now apply this comparison directly to each finite instance, avoiding any
continuity assertion across the floor in $W_{n,k}$. Write
\[
  \widehat \rho =\frac Dn,
  \qquad
  \widehat\eta=\frac An,
  \qquad
  \widehat h=\frac AD,
  \qquad
  \widehat\theta=1-\frac1{\widehat h},
  \qquad x=\frac kn.
\]
For every sufficiently large $n$,
$\widehat \rho <\rho _x$ and $\widehat\eta\geq\eta_x$. Hence
$\widehat h\geq h_x$, $\widehat\theta\geq\theta_x$, and
\[
  \widehat\eta d^{\widehat\theta-\delta}
  \geq\eta_xd^{\theta_x-\delta}\geq1.
\]
Also $C_{\tau_0}(\widehat h)\leq C_{\tau_0}(h_x)$. It follows that
\begin{equation}
  \frac{\widehat\eta d^{\widehat\theta}}H
  >(1+\zeta)C_{\tau_0}(\widehat h).
  \label{eq:boundary-finite-leading-margin}
\end{equation}
The triples $(\widehat \rho ,\widehat\eta,\tau_0)$ lie in one compact subset of
$\{0<\rho <\eta\leq1,\ \tau>0\}$ for all large $n$. The uniform version of
\Cref{thm:boundary-ratio}, applied with the actual floored
$W_{n,k}$, therefore turns \Cref{eq:boundary-finite-leading-margin} into
\[
  \widehat \rho \mathscr D_d(\widehat h)
  >\mathsf R_d(m,M,W_{n,k})
\]
with a uniform positive margin. Choose $d_0$ large enough also to ensure
$\alpha_{n,k}>0$ throughout this compact set.

Finally, the exact staircase identity gives
\begin{align*}
  \frac{\dim\Qspace(D,A,d,m,M,W_{n,k})}{n}
  &=\widehat \rho
  \sum_{\omega(c)\leq W_{n,k}}\sum_{b=0}^{M}
  \\
  &\qquad{}\cdot
  G\left(
    m\widehat h-(1-D^{-1})b
    -\sum_{j=2}^{d}(1-jD^{-1})c_j
  \right).
\end{align*}
The argument of \(G\) in each exact summand is
\[
  m\widehat h-b-|c|
  +\frac{b+\sum_{j=2}^djc_j}{D},
\]
which is at least \(m\widehat h-b-|c|\). Since \(\mathsf G\) is nondecreasing, the
exact dimension is therefore at least
\(\widehat \rho \mathscr D_d(\widehat h)\), with no limiting error. The strict
comparison above proves the finite certificate. Increasing \(N\) if
necessary also gives \(d<D\).
\end{proof}

The corollary identifies the sharp leading constant of this particular
one-sided enlarged-map certificate at its worst rate. It does not show that
the actual local rank, another interpolation support, or the weighted support
cannot do better. In fact, the weighted support is what gives the stronger
explicit $\exp(1.5/\delta)$ derivative order in the main paper.
 \end{document}
